%% ============================================================
%%  UNIVERSIDAD NACIONAL DE INGENIERÍA — FACULTAD DE CIENCIAS
%%  Escuela Profesional de Física
%%  Cuaderno: Mecánica Cuántica I
%%  Autor: Vera Vivanco, Jesús Rodolfo (20252703K)
%% ============================================================

\documentclass[a4paper, 11pt, oneside]{article}

%% ------------------------------------------------------------
%%  CODIFICACIÓN, IDIOMA Y FUENTES
%% ------------------------------------------------------------
\usepackage[utf8]{inputenc}
\usepackage[T1]{fontenc}
\usepackage[spanish, es-tabla]{babel}
\usepackage{csquotes}
\usepackage{lmodern}

%% ------------------------------------------------------------
%%  GEOMETRÍA Y MÁRGENES
%% ------------------------------------------------------------
\usepackage[top=2.5cm, bottom=2.5cm, left=2.8cm, right=2.5cm]{geometry}

%% ------------------------------------------------------------
%%  MATEMÁTICAS AVANZADAS
%% ------------------------------------------------------------
\usepackage{amsmath}
\usepackage{amssymb}
\usepackage{mathtools}
\usepackage{bm}
\usepackage{cancel}

%% ------------------------------------------------------------
%%  TABLAS Y FIGURAS
%% ------------------------------------------------------------
\usepackage{graphicx}
\graphicspath{{Figures/}}
\usepackage{float}
\usepackage{caption}
\usepackage{subcaption}
\usepackage{booktabs}
\usepackage{array}
\usepackage{tikz}
\usepackage{pgfplots}
\pgfplotsset{compat=1.17}
\usetikzlibrary{calc,babel,arrows.meta,positioning,angles,quotes,decorations.pathreplacing}
\usepgfplotslibrary{groupplots}

%% ------------------------------------------------------------
%%  COLORES E HIPERVÍNCULOS
%% ------------------------------------------------------------
\usepackage[dvipsnames]{xcolor}
\definecolor{uni-blue}{RGB}{0, 56, 101}
\definecolor{box-blue}{RGB}{240, 245, 252}
\definecolor{box-border}{RGB}{0, 56, 101}
\definecolor{alert-bg}{RGB}{255, 248, 235}
\definecolor{alert-border}{RGB}{185, 95, 10}
\definecolor{example-bg}{RGB}{245, 248, 245}
\definecolor{example-border}{RGB}{35, 120, 60}

\usepackage[
    hypertexnames = false,
    colorlinks    = true,
    linkcolor     = uni-blue,
    citecolor     = uni-blue,
    urlcolor      = uni-blue
]{hyperref}

\definecolor{uniblue}{RGB}{0,64,128}
\definecolor{unired}{RGB}{128,0,0}
\definecolor{unigreen}{RGB}{0,120,60}
\definecolor{unigray}{RGB}{110,110,110}
\definecolor{azulUNI}{RGB}{0,56,101}
\definecolor{naranjaUNI}{RGB}{185,95,10}
\definecolor{verdeUNI}{RGB}{35,120,60}
\definecolor{rojoUNI}{RGB}{170,30,30}
\definecolor{grisUNI}{gray}{0.45}
\tikzset{
  caja/.style={draw=uniblue,fill=uniblue!8,rounded corners=3pt,
    align=center,inner sep=5pt,font=\small,minimum height=1.0cm,text width=3.3cm},
  flecha/.style={-{Stealth[length=2.5mm]},thick,uniblue},
}
%% ------------------------------------------------------------
%%  CAJAS DIDÁCTICAS (tcolorbox)
%% ------------------------------------------------------------
\usepackage[most]{tcolorbox}

\newtcolorbox{teoremabox}[1][]{
    enhanced,
    colback=box-blue,
    colframe=box-border,
    fonttitle=\bfseries\color{white},
    colbacktitle=uni-blue,
    title={#1},
    arc=3mm,
    boxrule=1pt,
    left=8pt, right=8pt, top=7pt, bottom=7pt,
    before skip=10pt, after skip=10pt,
    breakable
}

\newtcolorbox{alertabox}[1][]{
    enhanced,
    colback=alert-bg,
    colframe=alert-border,
    fonttitle=\bfseries\color{alert-border},
    colbacktitle=alert-bg,
    title={$\blacktriangleright$ #1},
    arc=2mm,
    boxrule=1.2pt,
    left=8pt, right=8pt, top=7pt, bottom=7pt,
    before skip=10pt, after skip=10pt,
    breakable
}

\newtcolorbox{ejemplobox}[1][]{
    enhanced,
    colback=example-bg,
    colframe=example-border,
    fonttitle=\bfseries\color{white},
    colbacktitle=example-border,
    title={#1},
    arc=2mm,
    boxrule=1pt,
    left=8pt, right=8pt, top=7pt, bottom=7pt,
    before skip=10pt, after skip=10pt,
    breakable
}

\newtcolorbox{ejerciciobox}[1][]{
    enhanced,
    colback=white,
    colframe=uni-blue,
    fonttitle=\bfseries\color{white},
    colbacktitle=uni-blue,
    title={#1},
    arc=2mm,
    boxrule=1pt,
    left=8pt, right=8pt, top=7pt, bottom=7pt,
    before skip=10pt, after skip=10pt,
    breakable
}

%% ------------------------------------------------------------
%%  COMANDOS CUÁNTICOS PERSONALIZADOS
%% ------------------------------------------------------------
\newcommand{\ket}[1]{\left| #1 \right\rangle}
\newcommand{\bra}[1]{\left\langle #1 \right|}
\newcommand{\braket}[2]{\left\langle #1 \middle| #2 \right\rangle}
\newcommand{\matrixel}[3]{\left\langle #1 \middle| #2 \middle| #3 \right\rangle}
\newcommand{\comm}[2]{\left[ #1, #2 \right]}
\newcommand{\acomm}[2]{\left\{ #1, #2 \right\}}
\newcommand{\ev}[1]{\left\langle #1 \right\rangle}
\newcommand{\abs}[1]{\left| #1 \right|}
\newcommand{\norm}[1]{\left\| #1 \right\|}
\newcommand{\iu}{\mathrm{i}}
\newcommand{\eu}{\mathrm{e}}
\newcommand{\dif}{\mathrm{d}}
\newcommand{\adj}{^\dagger}

%% ------------------------------------------------------------
%%  ENCABEZADOS Y PIES DE PÁGINA
%% ------------------------------------------------------------
\usepackage{fancyhdr}
\pagestyle{fancy}
\fancyhf{}
\fancyhead[L]{\small\color{uni-blue}\textbf{\cursoNombre}\quad\small\color{gray!70!black}\textbar\quad\small\textnormal{\cursoCodigo}}
\fancyhead[R]{\small\textit{Fundamentos y Momento Angular}}
\fancyfoot[C]{\small\thepage}
\renewcommand{\headrulewidth}{0.4pt}

%% ------------------------------------------------------------
%%  ESTILO DE SECCIONES
%% ------------------------------------------------------------
\usepackage{titlesec}
\newcommand{\sectionbreak}{\clearpage}
\titleformat{\section}
    [block]{\Large\bfseries\color{uni-blue}}{\thesection.}{0.6em}{}
    [\vspace{2pt}{\color{uni-blue}\titlerule}\vspace{4pt}]
\titleformat{\subsection}
    [block]{\large\bfseries}{\thesubsection}{0.6em}{}
\titleformat{\subsubsection}
    [block]{\normalsize\bfseries\color{uni-blue}}{\thesubsubsection}{0.6em}{}

\titlespacing*{\section}       {0pt}{1.8em}{0.8em}
\titlespacing*{\subsection}    {0pt}{1.2em}{0.4em}
\titlespacing*{\subsubsection} {0pt}{1.0em}{0.3em}
\setcounter{secnumdepth}{3}
\setcounter{tocdepth}{3}

%% ------------------------------------------------------------
%%  ÍNDICE GENERAL
%% ------------------------------------------------------------
\usepackage{tocloft}
\addto\captionsspanish{
    \renewcommand{\contentsname}{Índice General de Contenidos}
}
\renewcommand{\cftsecfont}{\normalfont\bfseries\color{uni-blue}}
\renewcommand{\cftsecpagefont}{\normalfont\bfseries\color{uni-blue}}
\renewcommand{\cftsecdotsep}{\cftdotsep}
\renewcommand{\cftsubsecdotsep}{\cftdotsep}
\renewcommand{\cftsubsubsecdotsep}{\cftdotsep}
\setlength{\cftbeforesecskip}{6pt}
\setlength{\cftbeforesubsecskip}{2pt}
\setlength{\cftbeforesubsubsecskip}{1pt}

%% ------------------------------------------------------------
%%  ESPACIADO Y PÁRRAFOS
%% ------------------------------------------------------------
\usepackage{setspace}
\setstretch{1.15}
\setlength{\parskip}{0.4em}
\setlength{\parindent}{0pt}

%% ------------------------------------------------------------
%%  BIBLIOGRAFÍA (BibLaTeX + Biber)
%% ------------------------------------------------------------
\usepackage[backend=biber, style=numeric, sorting=none]{biblatex}
\addbibresource{refs.bib}

%% ------------------------------------------------------------
%%  ANEXOS
%% ------------------------------------------------------------
\usepackage[toc, page]{appendix}
\renewcommand{\appendixname}{Anexo}
\renewcommand{\appendixtocname}{Anexos}
\renewcommand{\appendixpagename}{Anexos}

%% ------------------------------------------------------------
%%  MICROTYPE
%% ------------------------------------------------------------
\usepackage[nopatch={toc,footnote}]{microtype}

%% ============================================================
%%  DATOS DEL DOCUMENTO
%% ============================================================
\newcommand{\cursoNombre}{Mecánica Cuántica I}
\newcommand{\cursoCodigo}{CF3A1}
\newcommand{\docTitulo}{Fundamentos Operacionales, Álgebra de Conmutadores, Incertidumbre Generalizada y Momento Angular Tensorial}
\newcommand{\docSubtitulo}{De los Espacios de Hilbert al Álgebra de Lie $\mathfrak{so}(3)$ y Conjuntos Completos de Observables Compatibles (CCOC)}
\newcommand{\docFecha}{Lima, Octubre de 2026}
\newcommand{\docAutor}{Vera Vivanco, Jesús Rodolfo}
\newcommand{\docCodigo}{20252703K}

%% ============================================================
\begin{document}

%% ------------------------------------------------------------
%%  PORTADA INSTITUCIONAL
%% ------------------------------------------------------------
\begin{titlepage}
    \centering
    \vspace*{0.5cm}
    {\Large\bfseries UNIVERSIDAD NACIONAL DE INGENIERÍA}\\[0.3cm]
    {\large\bfseries FACULTAD DE CIENCIAS}\\[0.1cm]
    {\large Escuela Profesional de Física}\\[1.0cm]
    \includegraphics[width=0.25\linewidth]{uni_logo}\\[1.0cm]
    {\normalsize\textbf{\cursoNombre} (\cursoCodigo)}\\[0.5cm]
    {\LARGE\bfseries\color{uni-blue} \docTitulo}\\[0.4cm]
    {\large\itshape \docSubtitulo}\\[1.2cm]
    \begin{minipage}{0.8\textwidth}
        \centering
        \textbf{Autor:}\\[0.2cm]
        {\large \docAutor}\\[0.1cm]
        \texttt{Código: \docCodigo}\\[0.6cm]
        \textit{Basado en referencias canónicas (\cite{zettili, sakurai, cohen, dirac, ballentine}).}
    \end{minipage}
    \vfill
    {\normalsize \docFecha}
\end{titlepage}

\pagenumbering{roman}

%% ------------------------------------------------------------
%%  ÍNDICE GENERAL
%% ------------------------------------------------------------
\tableofcontents
\newpage
\pagenumbering{arabic}

%% ============================================================
%%  SECCIONES MODULARES
%% ============================================================
%% ============================================================
%%  SECCIÓN 1 — Introducción y Hoja de Ruta Conceptual
%% ============================================================

\section{Introducción y Hoja de Ruta Conceptual}
\label{sec:introduccion}

\subsection{Motivación: Del Formalismo Ondulatorio al Enfoque Operacional}

En los cursos introductorios de física moderna y mecánica cuántica elemental, la descripción de los sistemas microscópicos suele presentarse a través de la \textit{mecánica ondulatoria} formulada por Erwin Schrödinger en 1926. En dicho marco, el estado de una partícula queda determinado por una función de onda continua $\psi(\vec{r}, t)$ en el espacio de coordenadas, y las magnitudes físicas observables se asocian a operadores diferenciales (por ejemplo, el momento lineal $\hat{p}_x = -\iu\hbar\,\partial/\partial x$).

Si bien este enfoque resulta intuitivo y sumamente eficaz para potenciales continuos unidimensionales (como el pozo infinito o el oscilador armónico en coordenadas espaciales), presenta limitaciones conceptuales severas cuando se busca generalizar la teoría a:
\begin{enumerate}
    \item Grados de libertad intrínsecos sin análogo espacial clásico, tales como el \textbf{espín} del electrón ($s = 1/2$), el cual no puede describirse mediante una simple función de onda escalar espacial $\psi(\vec{r})$, sino que requiere asociar a cada punto del espacio un vector de dos componentes (un espinor en $\mathcal{L}^2(\mathbb{R}^3) \otimes \mathbb{C}^2$, especificando las amplitudes para la proyección del espín $\pm\hbar/2$ a lo largo del eje polar $z$, denotadas por $\ket{\uparrow}$ o $\ket{\downarrow}$).
    \item Transformaciones geométricas continuas (rotaciones espaciales finitas, traslaciones, inversiones temporales y paridad).
    \item La comprensión algebraica profunda de las leyes de conservación y de las relaciones de incompatibilidad cuántica.
\end{enumerate}

Para superar estas limitaciones, Paul Dirac (1930) \cite{dirac} y John von Neumann (1932) reformularon la mecánica cuántica como una teoría geométrica y algebraica abstracta desarrollada en \textbf{espacios de Hilbert complejos} $\mathcal{H}$. En esta perspectiva:
\begin{itemize}
    \item El estado físico de un sistema no es una función de coordenadas arbitraria, sino un vector de estado o \textbf{ket} $\ket{\psi} \in \mathcal{H}$.
    \item Los observables físicos son \textbf{operadores lineales autoadjuntos} (hermíticos) $\hat{A} = \hat{A}\adj$ que actúan sobre dichos vectores.
    \item Las magnitudes medibles y sus incertidumbres estadísticas emergen directamente del producto interno $\braket{\phi}{\psi}$ y de las propiedades de conmutación $[\hat{A}, \hat{B}] = \hat{A}\hat{B} - \hat{B}\hat{A}$.
\end{itemize}

\subsection{El Puente Epistemológico: Álgebra Lineal para Físicos y la Mecánica Cuántica}

Para que esta transición sea verdaderamente comprensible y no parezca un conjunto de definiciones arbitrarias, conviene conectar estos conceptos con la estructura formal estudiada en \textit{Álgebra Lineal para Físicos} \cite{sakurai}:
\begin{enumerate}
    \item \textbf{Estados Físicos frente a Aparatos de Medida:}
    En un espacio vectorial $V$, los vectores representan estados o configuraciones del sistema. Sin embargo, ningún instrumento de laboratorio mide un vector ket completo: un detector siempre registra una magnitud observable macroscópica (un número real).
    Por el principio de superposición, la proyección sobre un estado particular se formula mediante un \textbf{funcional lineal} $f: V \to \mathbb{C}$ en el \textbf{espacio dual} $V^*$.
    Dirac sintetizó esta dualidad experimental de forma insuperable:
    \begin{equation}
        \boxed{\text{\textbf{Estados del Sistema:}} \quad \ket{\psi} \in \mathcal{H} \qquad \longleftrightarrow \qquad \text{\textbf{Canales de Detección:}} \quad \bra{\phi} \in \mathcal{H}^*}.
    \end{equation}
    El emparejamiento bracket $\braket{\phi}{\psi} \in \mathbb{C}$ representa la \textbf{amplitud de probabilidad} cuántica. La lectura física medible (o probabilidad de transición registrada por el detector) viene dada por la regla de Born como el número real positivo $P(\phi|\psi) = |\braket{\phi}{\psi}|^2 \in [0, 1]$.

    \item \textbf{El Teorema de Riesz como Puente Métrico:}
    La introducción de un producto interno hermítico $\braket{\cdot}{\cdot}$ establece un isomorfismo antilineal canónico entre el espacio de estados y su dual. A cada ket $\ket{\psi}$ le corresponde unívocamente un bra $\bra{\psi} = \ket{\psi}\adj$.

    \item \textbf{Operadores como Preguntas Experimentales y el Teorema Espectral:}
    ¿Por qué los observables son operadores hermíticos $\hat{A} = \hat{A}\adj$?
    Porque por el \textbf{Teorema Espectral de operadores normales}, un operador es autoadjunto si y solo si todos sus autovalores son estrictamente \textbf{reales} ($\lambda_n \in \mathbb{R}$, las únicas lecturas posibles en la carátula del instrumento) y sus autoestados forman una base ortonormal completa de proyectores unidimensionales $\hat{P}_n = \ket{n}\bra{n}$ que descomponen el espacio:
    \begin{equation}
        \hat{A} = \sum_n \lambda_n \ket{n}\bra{n}, \qquad \sum_n \ket{n}\bra{n} = \hat{I}.
    \end{equation}

    \item \textbf{Exponenciales de Operadores y Grupos de Lie:}
    Al igual que en sistemas de ecuaciones diferenciales acopladas $\dot{\vec{x}} = A\vec{x} \implies \vec{x}(t) = \eu^{At}\vec{x}(0)$, la ecuación de Schrödinger $\iu\hbar\frac{\dif}{\dif t}\ket{\psi} = \hat{H}\ket{\psi}$ se integra formalmente mediante el operador unitario $\ket{\psi(t)} = \eu^{-\iu\hat{H}t/\hbar}\ket{\psi(0)}$.
    De la misma manera que una matriz antisimétrica genera rotaciones en $\mathbb{R}^3$ ($P(t) = \eu^{\Omega t} \in \mathrm{SO}(3)$), el operador hermítico de momento $\hat{p}$ genera traslaciones espaciales $\eu^{-\iu a\hat{p}/\hbar}$ y el momento angular $\hat{\vec{L}}$ genera rotaciones cuánticas $\eu^{-\iu\vec{\theta}\cdot\hat{\vec{L}}/\hbar}$.
\end{enumerate}

El presente documento constituye una guía rigurosa, exhaustiva y pedagógica de la primera unidad de un curso de Mecánica Cuántica I. Nuestro propósito es desmitificar cada paso algebraico, evitando los habituales \enquote{saltos de fe} matemáticos y dotando al estudiante de intuición física y solidez formal \cite{zettili, sakurai}.

\subsection{Objetivos de Aprendizaje}

Al culminar el estudio de este cuaderno, el estudiante será capaz de:
\begin{itemize}
    \item \textbf{Dominar la notación de Dirac:} Manejar con destreza vectores ket, funcionales bra, relaciones de clausura (completitud) y representaciones matriciales de operadores.
    \item \textbf{Deducir el Principio de Incertidumbre Generalizado:} Construir rigurosamente la desigualdad de Cauchy--Schwarz en $\mathcal{H}$, descomponer el producto de operadores en sus partes simétrica (anticonmutador) y antisimétrica (conmutador), y derivar la desigualdad de \textbf{Robertson--Schrödinger}, comprendiendo el significado exacto de la covarianza cuántica y su reducción al principio canónico de Heisenberg.
    \item \textbf{Manipular Funciones Analíticas de Operadores:} Aplicar el cálculo diferencial con parámetros auxiliares para demostrar el \textbf{Lema de Hadamard} (trenzado de operadores) y el teorema de \textbf{Baker--Campbell--Hausdorff (BCH)}, aplicando ambos a generadores de traslación temporal y espacial.
    \item \textbf{Operar con Tensores en Momento Angular:} Utilizar la notación de índices de Einstein, el delta de Kronecker $\delta_{ij}$ y el pseudotensor de Levi-Civita $\varepsilon_{ijk}$ para deducir el álgebra de Lie $\mathfrak{so}(3)$, resolver la aparente paradoja $\hat{\vec{L}} \times \hat{\vec{L}} = \iu\hbar\hat{\vec{L}} \neq \vec{0}$, y demostrar la invariancia del operador Casimir $\hat{L}^2$.
    \item \textbf{Estructurar Conjuntos Completos (CCOC):} Justificar algebraicamente por qué las componentes del momento angular no pueden medirse simultáneamente y cómo la base $\{ \ket{\ell, m} \}$ resuelve la degeneración espectral mediante el conjunto $\{\hat{H}, \hat{L}^2, \hat{L}_z\}$.
\end{itemize}

\subsection{Hoja de Ruta Conceptual}

El desarrollo temático del documento sigue una ilación rigurosa donde cada concepto se apoya directamente en el anterior, tal como se esquematiza en la Tabla~\ref{tab:hoja_ruta}.

\begin{table}[H]
    \centering
    \caption{Secuencia lógica y conceptual de las secciones del documento.}
    \label{tab:hoja_ruta}
    \vspace{0.2cm}
    \resizebox{\linewidth}{!}{%
    \begin{tabular}{@{}clp{9.5cm}@{}}
        \toprule
        \textbf{Sección} & \textbf{Eje Temático} & \textbf{Hito Conceptual y Matemático Clave} \\
        \midrule
        \textbf{2} & Espacios de Hilbert y Dirac & Axiomática de kets/bras, diccionario matricial y producto externo, postulado de colapso, Teorema 2.2 (Espectro y completitud), Teorema 2.3 (Diagonalización simultánea y ejemplo $3\times 3$), taller de conmutadores fundamentales ($[\hat{x},\hat{p}_x]=\iu\hbar\hat{I}$, operadores $\hat{a},\hat{a}\adj$) y dispersión cuántica. \\
        \addlinespace[3pt]
        \textbf{3} & Incertidumbre Generalizada & Cauchy--Schwarz en $\mathcal{H}$, descomposición $[\Delta\hat{A}, \Delta\hat{B}]$ y $\{\Delta\hat{A}, \Delta\hat{B}\}$, covarianza cuántica, desigualdad de Robertson--Schrödinger y saturación en paquetes gaussianos. \\
        \addlinespace[3pt]
        \textbf{4} & Álgebra de Lie y Lemas BCH & Exponenciales operacionales, adjunto de funciones de operadores, Lema de Hadamard, BCH simplificado (c-números canónicos) frente a la serie general de Lie para momento angular, traslaciones espaciales, imagen de Heisenberg y Teorema de Noether cuántico. \\
        \addlinespace[3pt]
        \textbf{5} & Momento Angular y Producto Tensorial & Notación de Levi-Civita $\varepsilon_{ijk}$, conexión con $\mathfrak{so}(3)$ en Álgebra Lineal, producto cruz cuántico $\hat{\vec{L}}\times\hat{\vec{L}}=\iu\hbar\hat{\vec{L}}$, invariancia del Casimir $\hat{L}^2$, operadores escalera $\hat{L}_\pm$, matrices explícitas para $\ell=1$, laplaciano esférico, CCOC central y producto tensorial $\mathcal{H}_1 \otimes \mathcal{H}_2$ con acoplamiento de momentos angulares (bloque $M=0$). \\
        \addlinespace[3pt]
        \textbf{6} & Banco de Problemas Resueltos y Propuestos & 5 problemas resueltos paso a paso (rotaciones finitas, operadores vectoriales, incertidumbre en $\ket{\ell, m}$, efecto Zeeman y precesión de Larmor vía Heisenberg, oscilador 2D isótropo y modos circulares $\ket{R}, \ket{L}$) y 5 ejercicios propuestos de nivel avanzado. \\
        \addlinespace[3pt]
        \textbf{7} & Bibliografía de Referencia & Obras canónicas de consulta recomendadas (Zettili, Sakurai, Cohen-Tannoudji, Dirac, Ballentine). \\
        \addlinespace[3pt]
        \textbf{Anexos} & Herramientas y Solucionario & Demostración matricial de Levi-Civita, Formulario Maestro de conmutadores y Solucionario Analítico detallado de todos los ejercicios propuestos (1 al 5). \\
        \bottomrule
    \end{tabular}%
    }
\end{table}

\begin{figure}[H]
\centering
\begin{tikzpicture}[node distance=0.9cm and 1.0cm, caja/.append style={text width=4.6cm, font=\footnotesize}]
  \node[caja] (h) {\textbf{2. Hilbert y Dirac}\\ kets, bras, $\hat A^\dagger=\hat A$,\\ $\langle\hat A\rangle$, $\Delta\hat A$};
  \node[caja,right=of h] (u) {\textbf{3. Incertidumbre}\\ Cauchy--Schwarz,\\ Robertson--Schr\"odinger};
  \node[caja,below=of h] (b) {\textbf{4.4--4.5 BCH}\\ $\eu^{\hat A}\hat B\eu^{-\hat A}$,\\ traslaciones};
  \node[caja,right=of b] (n) {\textbf{4.7 Simetr\'ias y Noether}\\ $[\hat G,\hat H]=0$\\ $\Rightarrow$ conservaci\'on};
  \node[caja,below=of b] (l) {\textbf{5. Momento angular}\\ $\varepsilon_{ijk}$, $\mathfrak{so}(3)$,\\ $\hat{\vec L}\times\hat{\vec L}=\iu\hbar\hat{\vec L}$};
  \node[caja,right=of l,fill=unired!10,draw=unired] (c) {\textbf{CCOC}\\ $\{\hat H,\hat L^2,\hat L_z\}$\\ $\ket{n,\ell,m_\ell}$ (sin esp\'in)};
  \node[caja,below=of l,text width=9.6cm] (p) {\textbf{6. Problemas y ejercicios}\; (rotaciones finitas, operadores vectoriales, $\Delta L_x\Delta L_y$)};

  \draw[flecha] (h) -- (u);
  \draw[flecha] (h) -- (b);
  \draw[flecha] (u.east) -- ++(0.6,0) |- node[pos=0.25,right,font=\scriptsize,align=left] {incompati-\\bilidad} (c.east);
  \draw[flecha] (b) -- (n);
  \draw[flecha] (b) -- (l);
  \draw[flecha] (n) -- (l);
  \draw[flecha] (l) -- (c);
  \draw[flecha] (l) -- (p);
  \draw[flecha] (c.south) |- ($(p.east)+(0.2,0)$);
\end{tikzpicture}
\caption{Hoja de ruta: dependencia l\'ogica entre las secciones.}
\label{fig:mapa_dependencias}
\end{figure}


\begin{alertabox}[Recomendación Didáctica de Estudio]
    La física cuántica operacional no se aprende contemplando pasivamente las fórmulas, sino reproduciendo cada línea de cálculo con lápiz y papel. A lo largo del texto encontrarás \textbf{Cajas de Alerta Conceptual} señalizadas con $\blacktriangleright$. En ellas se discuten explícitamente los \enquote{puntos ciegos} matemáticos (por qué un término se anula, por qué se inserta un operador identidad específico o qué hipótesis físicas justifican una aproximación). No saltes estas cajas: en ellas reside la clave para dominar los exámenes teóricos.
\end{alertabox}

%% ============================================================
%%  SECCIÓN 2 — Espacios de Hilbert, Notación de Dirac y Operadores
%% ============================================================

\section{Espacios de Hilbert, Notación de Dirac y Operadores Físicos}
\label{sec:espacio_hilbert}

\subsection{El Espacio de Hilbert como Escenario Cuántico}

En la formulación rigurosa de von Neumann, el estado físico instantáneo de un sistema cuántico aislado queda descrito por un rayo vectorial (un vector normalizado salvo una fase global arbitraria $\eu^{\iu\alpha}$) que habita en un \textbf{espacio de Hilbert complejo}, denotado por $\mathcal{H}$ \cite{dirac, sakurai}.

\begin{teoremabox}[Definición Formal: Espacio de Hilbert $\mathcal{H}$]
    Un espacio de Hilbert $\mathcal{H}$ es un espacio vectorial lineal sobre el cuerpo de los números complejos $\mathbb{C}$, dotado de un producto interno $\braket{\cdot}{\cdot}: \mathcal{H} \times \mathcal{H} \to \mathbb{C}$, que induce una norma $\norm{\psi} = \sqrt{\braket{\psi}{\psi}}$, y que es \textbf{completo} respecto a la métrica inducida por dicha norma (es decir, toda sucesión de Cauchy converge a un elemento dentro del propio espacio).
\end{teoremabox}

En mecánica cuántica no relativista de una partícula en tres dimensiones, $\mathcal{H}$ suele identificarse con el espacio funcional $\mathcal{L}^2(\mathbb{R}^3)$ de funciones de cuadrado integrable con respecto a la medida de Lebesgue:
\begin{equation}
    \int_{\mathbb{R}^3} |\psi(\vec{r})|^2\,\dif^3\vec{r} < \infty.
    \label{eq:l2_norm}
\end{equation}
Para grados de libertad discretos (como el espín $s = 1/2$), $\mathcal{H}$ es un espacio de dimensión finita isomorfo a $\mathbb{C}^2$.

\subsection{Dualidad Bra-Ket y Axiomática del Producto Interno}

Paul Dirac introdujo una notación sumamente elegante y potente que desacopla la física de cualquier sistema particular de coordenadas:
\begin{itemize}
    \item Los elementos de $\mathcal{H}$ se denominan \textbf{kets} y se denotan por $\ket{\psi}$.
    \item Los elementos del espacio dual continuo $\mathcal{H}^*$ (el espacio de funcionales lineales continuos que transforman kets en números complejos $\mathbb{C}$) se denominan \textbf{bras} y se denotan por $\bra{\phi}$.
\end{itemize}

Por el \textit{Teorema de Representación de Riesz}, existe un isomorfismo antilineal canónico entre $\mathcal{H}$ y su dual $\mathcal{H}^*$: a todo ket $\ket{\psi}$ le corresponde biunívocamente un bra $\bra{\psi}$, denominado su adjunto o conjugado hermítico:
\begin{equation}
    \ket{\psi}\adj = \bra{\psi}, \qquad (\alpha\ket{\phi} + \beta\ket{\psi})\adj = \alpha^*\bra{\phi} + \beta^*\bra{\psi}.
    \label{eq:riesz_adjoint}
\end{equation}

El producto interno entre un bra $\bra{\phi}$ y un ket $\ket{\psi}$ se denota por el \textit{bracket} $\braket{\phi}{\psi} \in \mathbb{C}$ y satisface los siguientes axiomas fundamentales:
\begin{enumerate}
    \item \textbf{Linealidad en el segundo argumento (ket):}
    \begin{equation}
        \braket{\phi}{\alpha\psi_1 + \beta\psi_2} = \alpha\braket{\phi}{\psi_1} + \beta\braket{\phi}{\psi_2}, \quad \forall\,\alpha, \beta \in \mathbb{C}.
    \end{equation}
    \item \textbf{Antilinealidad en el primer argumento (bra):}
    \begin{equation}
        \braket{\alpha\phi_1 + \beta\phi_2}{\psi} = \alpha^*\braket{\phi_1}{\psi} + \beta^*\braket{\phi_2}{\psi}.
    \end{equation}
    \item \textbf{Simetría conjugada hermítica:}
    \begin{equation}
        \braket{\phi}{\psi} = \braket{\psi}{\phi}^*.
        \label{eq:hermitian_symmetry}
    \end{equation}
    \item \textbf{Positividad definida:}
    \begin{equation}
        \braket{\psi}{\psi} \ge 0, \quad \text{y} \quad \braket{\psi}{\psi} = 0 \iff \ket{\psi} = 0.
        \label{eq:positive_definite}
    \end{equation}
\end{enumerate}

La condición física de normalización para un estado cuántico que representa una certeza de probabilidad total es:
\begin{equation}
    \norm{\psi}^2 = \braket{\psi}{\psi} = 1.
\end{equation}

\subsection{Bases Ortonormales y la Relación de Clausura}

Si $\{\ket{n}\}$ constituye una base discreta y ortonormal de $\mathcal{H}$, sus elementos satisfacen la condición de ortonormalidad:
\begin{equation}
    \braket{m}{n} = \delta_{mn},
    \label{eq:orthonormality}
\end{equation}
donde $\delta_{mn}$ es el delta de Kronecker. Cualquier estado arbitrario $\ket{\psi}$ puede expandirse de manera única como una combinación lineal:
\begin{equation}
    \ket{\psi} = \sum_n c_n \ket{n}.
\end{equation}
Multiplicando por la izquierda por el bra $\bra{m}$ y aplicando la Ec.~\eqref{eq:orthonormality}:
\begin{equation}
    \braket{m}{\psi} = \sum_n c_n \braket{m}{n} = \sum_n c_n \delta_{mn} = c_m.
\end{equation}
Sustituyendo estos coeficientes $c_n = \braket{n}{\psi}$ nuevamente en la expansión de $\ket{\psi}$:
\begin{equation}
    \ket{\psi} = \sum_n \ket{n}\braket{n}{\psi} = \left( \sum_n \ket{n}\bra{n} \right) \ket{\psi}.
\end{equation}
Como esta igualdad se cumple para cualquier vector $\ket{\psi} \in \mathcal{H}$, el operador entre paréntesis debe ser idéntico al operador identidad $\hat{I}$. Esta es la célebre \textbf{Relación de Clausura (o completitud)}:
\begin{equation}
    \boxed{\sum_n \ket{n}\bra{n} = \hat{I}}.
    \label{eq:closure_relation}
\end{equation}
La relación de clausura es una de las herramientas de cálculo más poderosas de la física cuántica: permite \enquote{interpolar} una base en cualquier producto interno o elemento de matriz, transformando ecuaciones abstractas de operadores en sumas algebraicas directas \cite{zettili}.

\subsection{Operadores Lineales y el Operador Adjunto Hermítico}

Un operador lineal $\hat{A}$ transforma vectores de $\mathcal{H}$ en otros vectores del mismo espacio:
\begin{equation}
    \hat{A}(\alpha\ket{\psi_1} + \beta\ket{\psi_2}) = \alpha\hat{A}\ket{\psi_1} + \beta\hat{A}\ket{\psi_2}.
\end{equation}
El \textbf{adjunto hermítico} de $\hat{A}$, denotado por $\hat{A}\adj$, se define formalmente a través del producto interno como el único operador que satisface:
\begin{equation}
    \braket{\phi}{\hat{A}\psi} = \braket{\hat{A}\adj\phi}{\psi}, \quad \forall\,\ket{\phi}, \ket{\psi} \in \mathcal{H}.
    \label{eq:def_adjoint}
\end{equation}
Equivalentemente, en la notación de brackets, si $\ket{\chi} = \hat{A}\ket{\psi}$, entonces el bra asociado es $\bra{\chi} = \bra{\psi}\hat{A}\adj$. De esta definición se deducen directamente las siguientes propiedades esenciales:
\begin{enumerate}
    \item $(\hat{A} + \hat{B})\adj = \hat{A}\adj + \hat{B}\adj$.
    \item $(\alpha\hat{A})\adj = \alpha^* \hat{A}\adj$, con $\alpha \in \mathbb{C}$.
    \item Inversión de orden en productos de operadores:
    \begin{equation}
        (\hat{A}\hat{B})\adj = \hat{B}\adj \hat{A}\adj.
        \label{eq:adjoint_product}
    \end{equation}
    \textit{Demostración:} Aplicando la definición dos veces sucesivas,
    $\braket{\phi}{(\hat{A}\hat{B})\psi} = \braket{\hat{A}\adj\phi}{\hat{B}\psi} = \braket{\hat{B}\adj(\hat{A}\adj\phi)}{\psi} = \braket{(\hat{B}\adj\hat{A}\adj)\phi}{\psi}$.
    \item Involución: $(\hat{A}\adj)\adj = \hat{A}$.
\end{enumerate}

\begin{teoremabox}[Clasificación de Operadores según su Adjunción]
    \begin{itemize}
        \item \textbf{Operador Hermítico (Autoadjunto):} Cumple $\hat{A}\adj = \hat{A}$.
        En mecánica cuántica, \textbf{todo observable físico medible se representa por un operador lineal hermítico}. El Teorema Espectral garantiza que todos sus autovalores son estrictamente reales ($\lambda_n \in \mathbb{R}$) y que sus autovectores correspondientes a autovalores distintos son mutuamente ortogonales.
        \item \textbf{Operador Antihermítico (o Antiautoadjunto):} Cumple $\hat{C}\adj = -\hat{C}$.
        Sus autovalores son imaginarios puros o cero, y su valor esperado en cualquier estado físico es estrictamente imaginario puro: $\ev{\hat{C}} \in \iu\mathbb{R}$.
        \item \textbf{Operador Unitario:} Cumple $\hat{U}\adj\hat{U} = \hat{U}\hat{U}\adj = \hat{I}$, lo que equivale a $\hat{U}\adj = \hat{U}^{-1}$. Los operadores unitarios preservan productos internos y normas, describiendo simetrías y la evolución temporal continua del sistema.
    \end{itemize}
\end{teoremabox}

\subsubsection{Diccionario Matricial del Bra-Ket y el Producto Externo}
\label{subsec:diccionario_matricial_braket}

Para conectar sin fisuras la notación abstracta de Dirac con el cálculo matricial concreto de Álgebra Lineal, fijemos una base ortonormal discreta $\{\ket{1}, \ket{2}, \dots, \ket{N}\}$ de un espacio de Hilbert $\mathcal{H}$ de dimensión $N$.

\begin{enumerate}
    \item \textbf{Kets como Vectores Columna y Bras como Vectores Fila Conjugados:}
    Todo ket $\ket{\psi} = \sum_{n=1}^N c_n\ket{n}$ queda representado unívocamente por el vector columna de sus componentes $c_n = \braket{n}{\psi}$, mientras que su bra adjunto $\bra{\psi} = \sum_{n=1}^N c_n^*\bra{n}$ corresponde al vector fila transpuesto conjugado:
    \begin{equation}
        \ket{\psi} \;\doteq\; \begin{pmatrix} c_1 \\ c_2 \\ \vdots \\ c_N \end{pmatrix}, \qquad
        \bra{\psi} = \ket{\psi}\adj \;\doteq\; \begin{pmatrix} c_1^* & c_2^* & \cdots & c_N^* \end{pmatrix}.
    \end{equation}

    \item \textbf{Producto Interno $\braket{\phi}{\psi}$ frente a Producto Externo $\ket{\psi}\bra{\phi}$:}
    En la notación de Dirac, el orden en que se yuxtaponen un bra y un ket cambia radicalmente la naturaleza matemática del objeto:
    \begin{itemize}
        \item \textbf{Producto Interno (Bra-Ket $\braket{\phi}{\psi}$):} Es una fila $1\times N$ multiplicada por una columna $N\times 1$, lo que arroja un \textbf{número complejo escalar} ($\mathbb{C}$):
        \begin{equation}
            \braket{\phi}{\psi} = \begin{pmatrix} d_1^* & d_2^* & \cdots & d_N^* \end{pmatrix} \begin{pmatrix} c_1 \\ c_2 \\ \vdots \\ c_N \end{pmatrix} = \sum_{n=1}^N d_n^* c_n \in \mathbb{C}.
        \end{equation}
        \item \textbf{Producto Externo (Ket-Bra $\ket{\psi}\bra{\phi}$):} Es una columna $N\times 1$ multiplicada por una fila $1\times N$, lo que produce un \textbf{operador lineal} representado por una \textbf{matriz cuadrada de $N\times N$}:
        \begin{equation}
            \ket{\psi}\bra{\phi} \;\doteq\; \begin{pmatrix} c_1 \\ c_2 \\ \vdots \\ c_N \end{pmatrix} \begin{pmatrix} d_1^* & d_2^* & \cdots & d_N^* \end{pmatrix} = \begin{pmatrix}
                c_1 d_1^* & c_1 d_2^* & \cdots & c_1 d_N^* \\
                c_2 d_1^* & c_2 d_2^* & \cdots & c_2 d_N^* \\
                \vdots & \vdots & \ddots & \vdots \\
                c_N d_1^* & c_N d_2^* & \cdots & c_N d_N^*
            \end{pmatrix}.
        \end{equation}
        En efecto, al actuar sobre cualquier ket $\ket{\chi}$, el operador $(\ket{\psi}\bra{\phi})$ produce $\left(\ket{\psi}\bra{\phi}\right)\ket{\chi} = \ket{\psi}\braket{\phi}{\chi} = (\braket{\phi}{\chi})\ket{\psi}$, que es un vector en la dirección de $\ket{\psi}$.
    \end{itemize}

    \item \textbf{¿Cómo se construye la matriz de un operador $\hat{A}$ en una base dada?}
    Insertando dos veces la relación de clausura $\hat{I} = \sum_m \ket{m}\bra{m}$ a ambos lados del operador $\hat{A}$:
    \begin{equation}
        \hat{A} = \hat{I}\hat{A}\hat{I} = \left(\sum_{m=1}^N \ket{m}\bra{m}\right) \hat{A} \left(\sum_{n=1}^N \ket{n}\bra{n}\right) = \sum_{m=1}^N \sum_{n=1}^N \underbrace{\matrixel{m}{\hat{A}}{n}}_{A_{mn}} \ket{m}\bra{n}.
    \end{equation}
    El número complejo $A_{mn} \equiv \matrixel{m}{\hat{A}}{n}$ es exactamente el \textbf{elemento de matriz en la fila $m$ y columna $n$}:
    \begin{equation}
        \boxed{A_{mn} = \matrixel{m}{\hat{A}}{n} = \braket{m}{\hat{A}n} \quad (\text{fila } m, \text{ columna } n)}.
    \end{equation}
    \textit{Regla práctica para construir la matriz:} Aplique el operador $\hat{A}$ sobre el $n$-ésimo vector base $\ket{n}$; los coeficientes del vector resultante constituyen la \textbf{$n$-ésima columna} de la matriz $A$.
\end{enumerate}

\begin{table}[H]
    \centering
    \caption{Diccionario de traducción entre el Álgebra Lineal Matricial y la Notación de Dirac.}
    \label{tab:diccionario_dirac}
    \vspace{0.15cm}
    \resizebox{\linewidth}{!}{%
    \begin{tabular}{@{}lll@{}}
        \toprule
        \textbf{Concepto Matemático} & \textbf{Notación de Dirac (Abstracta)} & \textbf{Representación Matricial en Base $\{\ket{n}\}$} \\
        \midrule
        Vector de estado & Ket $\ket{\psi}$ & Vector columna $N\times 1$: $(c_1, c_2, \dots, c_N)^T$ \\
        Funcional dual (adjunto) & Bra $\bra{\psi} = \ket{\psi}\adj$ & Vector fila conjugado $1\times N$: $(c_1^*, c_2^*, \dots, c_N^*)$ \\
        Producto interno (escalar) & Bracket $\braket{\phi}{\psi}$ & Producto fila por columna: $\sum_n d_n^* c_n \in \mathbb{C}$ \\
        Producto externo (operador) & Ket-Bra $\ket{\psi}\bra{\phi}$ & Matriz cuadrada $N\times N$ de entradas $c_m d_n^*$ \\
        Elemento de matriz de $\hat{A}$ & Sandwich $\matrixel{m}{\hat{A}}{n}$ & Entrada $A_{mn}$ (fila $m$, columna $n$) \\
        Operador adjunto $\hat{A}\adj$ & $\matrixel{m}{\hat{A}\adj}{n} = \matrixel{n}{\hat{A}}{m}^*$ & Matriz transpuesta conjugada: $(A\adj)_{mn} = A_{nm}^*$ \\
        Traza de un operador & $\mathrm{Tr}(\hat{A}) = \sum_n \matrixel{n}{\hat{A}}{n}$ & Suma de la diagonal principal: $\sum_n A_{nn}$ \\
        \bottomrule
    \end{tabular}%
    }
\end{table}

\subsection{Valor Esperado, Varianza y el Operador Desviación}

Cuando un sistema se encuentra en un estado físico normalizado $\ket{\psi}$ ($\braket{\psi}{\psi}=1$), la medición de un observable $\hat{A}$ no arroja en general un valor determinista, sino una distribución estadística de resultados.

\subsubsection{El Teorema Espectral y el Postulado de Medida de von Neumann}
¿De dónde surge la fórmula $\ev{\hat{A}} = \matrixel{\psi}{\hat{A}}{\psi}$? No es una definición ad-hoc, sino una consecuencia directa del \textbf{Teorema Espectral para operadores autoadjuntos} en espacios de Hilbert.

Dado un observable hermítico $\hat{A}$, existe una base ortonormal completa de autovectores $\{\ket{n}\}$ con autovalores reales $\{a_n\}$ tales que $\hat{A}\ket{n} = a_n \ket{n}$. Utilizando la relación de clausura $\sum_n \ket{n}\bra{n} = \hat{I}$, el operador $\hat{A}$ admite la descomposición espectral única en proyectores ortogonales unidimensionales $\hat{P}_n \equiv \ket{n}\bra{n}$:
\begin{equation}
    \hat{A} = \hat{A}\hat{I} = \hat{A}\left( \sum_n \ket{n}\bra{n} \right) = \sum_n a_n \ket{n}\bra{n} = \sum_n a_n \hat{P}_n.
    \label{eq:spectral_decomposition}
\end{equation}
donde los proyectores satisfacen $\hat{P}_n^2 = \hat{P}_n$, $\hat{P}_n\adj = \hat{P}_n$ y $\hat{P}_n\hat{P}_m = \delta_{nm}\hat{P}_n$.

Cuando el sistema está en un estado normalizado arbitrario $\ket{\psi}$, podemos expandirlo en dicha base ortonormal:
\begin{equation}
    \ket{\psi} = \sum_n c_n \ket{n}, \quad c_n = \braket{n}{\psi}.
\end{equation}
Por el \textbf{Postulado de Medida de Born y von Neumann}:
\begin{enumerate}
    \item Los únicos resultados posibles al medir el observable $\hat{A}$ son sus autovalores $\{a_n\}$.
    \item La probabilidad de obtener el resultado discreto $a_n$ está dada por el valor esperado del proyector $\hat{P}_n$:
    \begin{equation}
        P(a_n) = \matrixel{\psi}{\hat{P}_n}{\psi} = \braket{\psi}{n}\braket{n}{\psi} = |c_n|^2.
    \end{equation}
    Como $\norm{\psi}^2 = 1$, la probabilidad total está conservada: $\sum_n P(a_n) = \sum_n |c_n|^2 = 1$.
    \item \textbf{Colapso del Estado (Postulado de Proyección de von Neumann):} Si al efectuar la medición sobre $\ket{\psi}$ el aparato registra el autovalor específico $a_n$, el estado físico del sistema sufre un cambio abrupto e irreversible, proyectándose instantáneamente sobre el subespacio propio correspondiente:
    \begin{equation}
        \ket{\psi} \;\xrightarrow{\quad\text{medición de } a_n\quad}\; \frac{\hat{P}_n\ket{\psi}}{\sqrt{\matrixel{\psi}{\hat{P}_n}{\psi}}} = \frac{c_n}{|c_n|}\ket{n} \;\sim\; \ket{n}.
    \end{equation}
    Esto garantiza la reproducibilidad física: si se vuelve a medir el mismo observable $\hat{A}$ inmediatamente después sobre el sistema recién colapsado, se obtendrá con certeza absoluta ($100\%$) el mismo valor $a_n$.
\end{enumerate}

\subsubsection{Valor Esperado}
El valor esperado $\ev{\hat{A}}$ es el promedio estadístico ponderado de todos los resultados posibles en un ensamble estadístico de mediciones idénticas:
\begin{align}
    \ev{\hat{A}} &\equiv \sum_n a_n P(a_n) = \sum_n a_n \braket{\psi}{n}\braket{n}{\psi} \nonumber\\
    &= \bra{\psi} \left( \sum_n a_n \ket{n}\bra{n} \right) \ket{\psi} = \boxed{\matrixel{\psi}{\hat{A}}{\psi}}.
    \label{eq:expectation_value}
\end{align}
¡La fórmula $\matrixel{\psi}{\hat{A}}{\psi}$ es exactamente el promedio probabilístico ponderado por el espectro del operador!

Si $\hat{A}$ es hermítico, su valor esperado es obligatoriamente un número real:
\begin{equation}
    \ev{\hat{A}}^* = \left(\matrixel{\psi}{\hat{A}}{\psi}\right)^* = \matrixel{\psi}{\hat{A}\adj}{\psi} = \matrixel{\psi}{\hat{A}}{\psi} = \ev{\hat{A}} \implies \ev{\hat{A}} \in \mathbb{R}.
\end{equation}

\subsubsection{El Operador Desviación \texorpdfstring{$\Delta\hat{A}$}{Delta A}}
Para cuantificar la dispersión de las medidas alrededor de la media, definimos el \textbf{operador desviación} (o fluctuación cuántica):
\begin{equation}
    \boxed{\Delta\hat{A} \equiv \hat{A} - \ev{\hat{A}}\hat{I}}.
    \label{eq:deviation_operator}
\end{equation}
Observemos con cuidado la naturaleza matemática de este operador: $\hat{A}$ es un operador, mientras que $\ev{\hat{A}}$ es un escalar (un c-número real). Para poder restarlos con sentido algebraico riguroso, multiplicamos formalmente el escalar por el operador identidad $\hat{I}$ (frecuentemente omitido en la notación abreviada, pero implícito en todo cálculo).

\begin{alertabox}[Propiedad Fundamental del Operador Desviación]
    Como $\hat{A}$ es hermítico y $\ev{\hat{A}} \in \mathbb{R}$, el operador desviación es automáticamente \textbf{hermítico}:
    \begin{equation}
        (\Delta\hat{A})\adj = (\hat{A} - \ev{\hat{A}}\hat{I})\adj = \hat{A}\adj - \ev{\hat{A}}^*\hat{I}\adj = \hat{A} - \ev{\hat{A}}\hat{I} = \Delta\hat{A}.
    \end{equation}
    Además, su valor esperado promedio es idénticamente nulo:
    \begin{equation}
        \ev{\Delta\hat{A}} = \ev{\hat{A} - \ev{\hat{A}}\hat{I}} = \ev{\hat{A}} - \ev{\hat{A}}\ev{\hat{I}} = \ev{\hat{A}} - \ev{\hat{A}} = 0.
    \end{equation}
\end{alertabox}

\subsubsection{Varianza Cuántica}
La incertidumbre cuántica (o desviación estándar) $\Delta A$ se define como la raíz cuadrada de la varianza:
\begin{equation}
    (\Delta A)^2 \equiv \ev{(\Delta\hat{A})^2} = \matrixel{\psi}{(\Delta\hat{A})^2}{\psi}.
    \label{eq:variance_def}
\end{equation}
Expandiendo el cuadrado del operador:
\begin{align}
    (\Delta\hat{A})^2 &= (\hat{A} - \ev{\hat{A}}\hat{I})(\hat{A} - \ev{\hat{A}}\hat{I}) \nonumber\\
    &= \hat{A}^2 - \hat{A}\ev{\hat{A}}\hat{I} - \ev{\hat{A}}\hat{I}\hat{A} + \ev{\hat{A}}^2 \hat{I}^2 \nonumber\\
    &= \hat{A}^2 - 2\ev{\hat{A}}\hat{A} + \ev{\hat{A}}^2 \hat{I}.
\end{align}
Tomando el valor esperado en el estado $\ket{\psi}$:
\begin{align}
    (\Delta A)^2 &= \ev{\hat{A}^2 - 2\ev{\hat{A}}\hat{A} + \ev{\hat{A}}^2 \hat{I}} \nonumber\\
    &= \ev{\hat{A}^2} - 2\ev{\hat{A}}\ev{\hat{A}} + \ev{\hat{A}}^2 = \boxed{\ev{\hat{A}^2} - \ev{\hat{A}}^2}.
    \label{eq:variance_expanded}
\end{align}
Esta fórmula es formalmente análoga a la varianza estadística clásica $\sigma_A^2 = \overline{A^2} - (\overline{A})^2$, pero con una diferencia conceptual titánica: la varianza cuántica no proviene de imperfecciones o ignorancia instrumental, sino de la dispersión ontológica intrínseca del estado cuántico $\ket{\psi}$.

\subsection{Ejemplo Resuelto 1: Cálculo de Dispersión Cuántica en un Sistema de Dos Niveles}

\begin{ejemplobox}[Ejemplo Resuelto 1: Fluctuación Cuántica en una Matriz de Pauli]
    \textbf{Enunciado:} Considere un sistema de dos niveles (espín $1/2$) con base ortonormal $\{\ket{+z}, \ket{-z}\}$. El observable físico $\hat{S}_x$ está dado por la matriz de Pauli:
    \begin{equation}
        \hat{S}_x = \frac{\hbar}{2}\begin{pmatrix} 0 & 1 \\ 1 & 0 \end{pmatrix}.
    \end{equation}
    Suponga que el sistema se encuentra preparado en el estado cuántico $\ket{\psi} = \frac{1}{\sqrt{5}}\begin{pmatrix} 2 \\ 1 \end{pmatrix}$.
    \begin{enumerate}
        \item Verifique que $\ket{\psi}$ está normalizado y calcule $\ev{\hat{S}_x}$.
        \item Calcule $\ev{\hat{S}_x^2}$ y determine la incertidumbre $\Delta S_x$.
        \item Construya explícitamente el operador desviación $\Delta\hat{S}_x$ en forma matricial y verifique directamente que $\matrixel{\psi}{(\Delta\hat{S}_x)^2}{\psi} = (\Delta S_x)^2$.
    \end{enumerate}

    \textbf{Solución Paso a Paso:}
    
    \textit{Paso 1: Normalización y Valor Esperado:}
    Calculamos el producto interno:
    \begin{equation}
        \braket{\psi}{\psi} = \frac{1}{\sqrt{5}}\begin{pmatrix} 2 & 1 \end{pmatrix} \frac{1}{\sqrt{5}}\begin{pmatrix} 2 \\ 1 \end{pmatrix} = \frac{1}{5}(2^2 + 1^2) = \frac{5}{5} = 1.
    \end{equation}
    El estado está correctamente normalizado. Evaluamos ahora la acción de $\hat{S}_x$ sobre $\ket{\psi}$:
    \begin{equation}
        \hat{S}_x \ket{\psi} = \frac{\hbar}{2}\begin{pmatrix} 0 & 1 \\ 1 & 0 \end{pmatrix} \frac{1}{\sqrt{5}}\begin{pmatrix} 2 \\ 1 \end{pmatrix} = \frac{\hbar}{2\sqrt{5}}\begin{pmatrix} 1 \\ 2 \end{pmatrix}.
    \end{equation}
    Tomando el producto interno con $\bra{\psi}$:
    \begin{equation}
        \ev{\hat{S}_x} = \braket{\psi}{\hat{S}_x\psi} = \frac{1}{\sqrt{5}}\begin{pmatrix} 2 & 1 \end{pmatrix} \frac{\hbar}{2\sqrt{5}}\begin{pmatrix} 1 \\ 2 \end{pmatrix} = \frac{\hbar}{10}(2\cdot 1 + 1\cdot 2) = \frac{4\hbar}{10} = \frac{2}{5}\hbar = 0{,}40\,\hbar.
    \end{equation}

    \textit{Paso 2: Varianza Cuántica:}
    Notemos que para las matrices de Pauli $\sigma_x^2 = \hat{I}_2$. Por ende:
    \begin{equation}
        \hat{S}_x^2 = \left(\frac{\hbar}{2}\right)^2 \begin{pmatrix} 0 & 1 \\ 1 & 0 \end{pmatrix}\begin{pmatrix} 0 & 1 \\ 1 & 0 \end{pmatrix} = \frac{\hbar^2}{4}\begin{pmatrix} 1 & 0 \\ 0 & 1 \end{pmatrix} = \frac{\hbar^2}{4}\hat{I}.
    \end{equation}
    En consecuencia, $\ev{\hat{S}_x^2} = \frac{\hbar^2}{4}\ev{\hat{I}} = \frac{\hbar^2}{4}$. Aplicando la fórmula de la varianza:
    \begin{equation}
        (\Delta S_x)^2 = \ev{\hat{S}_x^2} - \ev{\hat{S}_x}^2 = \frac{\hbar^2}{4} - \left(\frac{2}{5}\hbar\right)^2 = \frac{\hbar^2}{4} - \frac{4}{25}\hbar^2 = \frac{25 - 16}{100}\hbar^2 = \frac{9}{100}\hbar^2.
    \end{equation}
    Extrayendo la raíz cuadrada:
    \begin{equation}
        \Delta S_x = \sqrt{\frac{9}{100}\hbar^2} = \frac{3}{10}\hbar = 0{,}30\,\hbar.
    \end{equation}

    \textit{Paso 3: Matriz del Operador Desviación:}
    Por definición, $\Delta\hat{S}_x = \hat{S}_x - \ev{\hat{S}_x}\hat{I}$:
    \begin{equation}
        \Delta\hat{S}_x = \frac{\hbar}{2}\begin{pmatrix} 0 & 1 \\ 1 & 0 \end{pmatrix} - \frac{2}{5}\hbar\begin{pmatrix} 1 & 0 \\ 0 & 1 \end{pmatrix} = \hbar \begin{pmatrix} -2/5 & 1/2 \\ 1/2 & -2/5 \end{pmatrix} = \frac{\hbar}{10}\begin{pmatrix} -4 & 5 \\ 5 & -4 \end{pmatrix}.
    \end{equation}
    Actuando con este operador sobre $\ket{\psi}$:
    \begin{equation}
        \Delta\hat{S}_x \ket{\psi} = \frac{\hbar}{10\sqrt{5}}\begin{pmatrix} -4 & 5 \\ 5 & -4 \end{pmatrix}\begin{pmatrix} 2 \\ 1 \end{pmatrix} = \frac{\hbar}{10\sqrt{5}}\begin{pmatrix} -8 + 5 \\ 10 - 4 \end{pmatrix} = \frac{\hbar}{10\sqrt{5}}\begin{pmatrix} -3 \\ 6 \end{pmatrix}.
    \end{equation}
    Calculando la norma cuadrada de este vector resultante:
    \begin{equation}
        \matrixel{\psi}{(\Delta\hat{S}_x)^2}{\psi} = \norm{\Delta\hat{S}_x\psi}^2 = \frac{\hbar^2}{100\cdot 5}\left( (-3)^2 + 6^2 \right) = \frac{\hbar^2}{500}(9 + 36) = \frac{45}{500}\hbar^2 = \frac{9}{100}\hbar^2,
    \end{equation}
    lo cual coincide exactamente con el valor calculado en el Paso 2.
\end{ejemplobox}

\begin{figure}[H]
\centering
\begin{subfigure}[t]{0.50\textwidth}
\centering
\begin{tikzpicture}[scale=2.0,
   proj/.style={x={(1cm,0cm)},y={(0.45cm,0.3cm)},z={(0cm,1cm)}}]
  \draw[gray!70] (0,0) circle (1);
  \begin{scope}[proj]
    \draw[gray!60,dashed,domain=0:360,samples=100,smooth,variable=\t]
      plot ({cos(\t)},{sin(\t)},0);
    \draw[gray!60,dashed,domain=0:360,samples=100,smooth,variable=\t]
      plot ({cos(\t)},0,{sin(\t)});
    \draw[->,gray] (0,0,0) -- (1.35,0,0) node[right] {$x$};
    \draw[->,gray] (0,0,0) -- (0,1.35,0) node[right] {$y$};
    \draw[->,gray] (0,0,0) -- (0,0,1.45) node[above] {$z$};
    \node[font=\scriptsize,below] at (0,0,-1.0) {$\ket{{-}z}$};
    \node[font=\scriptsize,above right] at (0,0,1.0) {$\ket{{+}z}$};
    \node[font=\scriptsize,unigreen,below right] at (1.0,0,0) {$\ket{{+}x}$};
    \draw[dotted,unired] (0.8,0,0.6) -- (0.8,0,0);
    \draw[dotted,unired] (0.8,0,0) -- (0,0,0);
    \draw[-{Stealth},very thick,unired] (0,0,0) -- (0.8,0,0.6);
    \fill[unired] (0.8,0,0.6) circle (0.025);
    \node[unired,font=\small,right] at (0.8,0,0.6) {$\vec r$};
  \end{scope}
\end{tikzpicture}
\caption{$\vec r=(0{,}8;\,0;\,0{,}6)$, con $\langle S_x\rangle=\tfrac{\hbar}{2}r_x=0{,}4\,\hbar$.}
\end{subfigure}\hfill
\begin{subfigure}[t]{0.46\textwidth}
\centering
\begin{tikzpicture}
\begin{axis}[ybar,width=\linewidth,height=5.6cm,ymin=0,ymax=1.1,
  bar width=22pt,symbolic x coords={m,p},xtick=data,
  xticklabels={$-\hbar/2$,$+\hbar/2$},
  ylabel={Probabilidad},xlabel={Resultado de $S_x$},
  nodes near coords,nodes near coords align={vertical},
  enlarge x limits=0.5]
  \addplot[fill=uniblue!60,draw=uniblue] coordinates {(m,0.1) (p,0.9)};
\end{axis}
\end{tikzpicture}
\caption{$P(\pm)=\frac{1\pm0{,}8}{2}$; $\Delta S_x=0{,}3\,\hbar$.}
\end{subfigure}
\caption{Ejemplo Resuelto 1: estado $\ket\psi=(2,1)^T/\sqrt5$ sobre la esfera de Bloch y distribuci\'on de resultados al medir $\hat S_x$.}
\label{fig:bloch}
\end{figure}

\subsection{Teorema Espectral y Completitud (Teorema 2.2)}
\label{subsec:teorema_2_2}

En el estudio del álgebra lineal cuántica, la descomposición espectral representa el puente definitivo entre los operadores lineales y la física de las mediciones experimentales. El postulado de que los observables físicos son operadores autoadjuntos descansa en dos propiedades matemáticas inquebrantables.

\begin{teoremabox}[Teorema 2.2: Espectro y Completitud de Operadores Hermíticos]
    Sea $\hat{A}$ un operador lineal hermítico ($\hat{A}\adj = \hat{A}$) definido en un espacio de Hilbert $\mathcal{H}$ de dimensión finita $N$ (o un operador autoadjunto compacto en un espacio infinito-dimensional). Entonces:
    \begin{enumerate}
        \item \textbf{Realidad de los Autovalores:} Todos los autovalores de $\hat{A}$ son estrictamente reales:
        \begin{equation}
            \hat{A}\ket{\phi_n} = a_n\ket{\phi_n} \implies a_n \in \mathbb{R}.
        \end{equation}
        \item \textbf{Ortogonalidad de Autoestados:} Los autovectores asociados a autovalores distintos son mutuamente ortogonales:
        \begin{equation}
            a_n \neq a_m \implies \braket{\phi_n}{\phi_m} = 0.
        \end{equation}
        \item \textbf{Completitud y Relación de Clausura:} Los autovectores normalizados $\{\ket{\phi_n}\}_{n=1}^N$ conforman una base ortonormal completa de $\mathcal{H}$, de modo que el operador identidad admite la resolución espectral:
        \begin{equation}
            \sum_{n=1}^N \ket{\phi_n}\bra{\phi_n} = \hat{I}.
        \end{equation}
        \item \textbf{Descomposición Espectral y Diagonalización:} El operador $\hat{A}$ puede expandirse de forma única en términos de sus proyectores ortogonales unidimensionales $\hat{P}_n \equiv \ket{\phi_n}\bra{\phi_n}$:
        \begin{equation}
            \hat{A} = \sum_{n=1}^N a_n \ket{\phi_n}\bra{\phi_n} = \sum_{n=1}^N a_n \hat{P}_n.
            \label{eq:descomposicion_espectral_teorica}
        \end{equation}
        En dicha base propia, la matriz representante de $\hat{A}$ es estrictamente diagonal:
        \begin{equation}
            A_{nm} = \matrixel{\phi_n}{\hat{A}}{\phi_m} = a_n \delta_{nm}.
        \end{equation}
    \end{enumerate}
\end{teoremabox}

\textbf{Demostración Didáctica y Rigurosa:}\\
\textit{1. Demostración de que los autovalores son reales:}
Sea $\ket{\phi_n}$ un autovector normalizado de $\hat{A}$ con autovalor $a_n$:
\begin{equation}
    \hat{A}\ket{\phi_n} = a_n\ket{\phi_n}.
\end{equation}
Multiplicamos por la izquierda por el bra $\bra{\phi_n}$:
\begin{equation}
    \matrixel{\phi_n}{\hat{A}}{\phi_n} = a_n\braket{\phi_n}{\phi_n} = a_n.
\end{equation}
Tomando el complejo conjugado y usando la simetría hermítica del producto interno y la condición $\hat{A}\adj = \hat{A}$:
\begin{equation}
    a_n^* = \left(\matrixel{\phi_n}{\hat{A}}{\phi_n}\right)^* = \matrixel{\phi_n}{\hat{A}\adj}{\phi_n} = \matrixel{\phi_n}{\hat{A}}{\phi_n} = a_n.
\end{equation}
Como $a_n^* = a_n$, el autovalor $a_n$ debe ser forzosamente un número real ($a_n \in \mathbb{R}$).

\textit{2. Demostración de la ortogonalidad para autovalores distintos:}
Consideremos dos autovectores $\ket{\phi_n}$ y $\ket{\phi_m}$ tales que:
\begin{equation}
    \hat{A}\ket{\phi_n} = a_n\ket{\phi_n}, \qquad \hat{A}\ket{\phi_m} = a_m\ket{\phi_m}, \quad \text{con } a_n \neq a_m.
\end{equation}
Evaluemos el elemento de matriz $\matrixel{\phi_m}{\hat{A}}{\phi_n}$ de dos maneras distintas:
\begin{align}
    \text{Por la derecha (acción sobre ket):} \quad &\matrixel{\phi_m}{\hat{A}}{\phi_n} = \bra{\phi_m}(a_n\ket{\phi_n}) = a_n\braket{\phi_m}{\phi_n}. \\
    \text{Por la izquierda (acción sobre bra):} \quad &\matrixel{\phi_m}{\hat{A}}{\phi_n} = \braket{\hat{A}\adj\phi_m}{\phi_n} = \braket{\hat{A}\phi_m}{\phi_n} = (a_m\bra{\phi_m})\ket{\phi_n} = a_m\braket{\phi_m}{\phi_n},
\end{align}
donde hemos utilizado que $a_m^* = a_m$. Restando ambas igualdades:
\begin{equation}
    (a_n - a_m)\braket{\phi_m}{\phi_n} = 0.
\end{equation}
Dado que por hipótesis los autovalores son distintos ($a_n \neq a_m \implies a_n - a_m \neq 0$), el producto interno debe anularse obligatoriamente:
\begin{equation}
    \boxed{\braket{\phi_m}{\phi_n} = 0 \quad (\forall\,n \neq m)}.
\end{equation}

\textit{3. Completitud y diagonalización (inducción sobre la dimensión):}\\
Por el teorema fundamental del álgebra, el polinomio característico de $\hat{A}$ posee al menos una raíz en $\mathbb{C}$, garantizando la existencia de un primer autovector normalizado $\ket{\phi_1}$ con autovalor real $a_1$. Sea $W = \mathrm{span}\{\ket{\phi_1}\}$ y $W^\perp$ su complemento ortogonal en $\mathcal{H}$. Para todo $\ket{\psi} \in W^\perp$, se cumple $\braket{\phi_1}{\hat{A}\psi} = \matrixel{\phi_1}{\hat{A}}{\psi} = \braket{\hat{A}\phi_1}{\psi} = a_1\braket{\phi_1}{\psi} = 0$, lo que demuestra que el subespacio $W^\perp$ de dimensión $N-1$ es invariante bajo $\hat{A}$. Por hipótesis inductiva sobre la dimensión, la restricción de $\hat{A}$ a $W^\perp$ es hermítica y admite una base ortonormal de autovectores. Si un autovalor $a_n$ posee multiplicidad geométrica $g_n > 1$ (subespacio propio degenerado $E_{a_n}$), el procedimiento de ortonormalización de Gram--Schmidt permite elegir siempre una base ortonormal interna para $E_{a_n}$. La unión de dichas bases conforma una base ortonormal completa de todo el espacio $\mathcal{H}$, demostrando la relación de clausura $\sum_{n=1}^N \ket{\phi_n}\bra{\phi_n} = \hat{I}$ y la diagonalizabilidad de $\hat{A}$.

\subsection{Conmutación de Observables y Diagonalización Simultánea (Teorema 2.3)}
\label{subsec:teorema_2_3}

Uno de los resultados más trascendentales de la física cuántica establece la relación íntima entre la \textbf{compatibilidad física de dos mediciones} y la \textbf{conmutación algebraica de sus operadores}.

\begin{teoremabox}[Teorema 2.3: Operadores Compatibles y Base Común]
    Dos operadores hermíticos $\hat{A}$ y $\hat{B}$ admiten una base ortonormal completa de autovectores simultáneos comunes si y solo si \textbf{conmutan entre sí}:
    \begin{equation}
        [\hat{A}, \hat{B}] = \hat{A}\hat{B} - \hat{B}\hat{A} = 0.
    \end{equation}
\end{teoremabox}

\textbf{Demostración de la Condición Necesaria ($\Leftarrow$):}\\
Si suponemos que existe una base ortonormal completa común $\{\ket{n}\}$ tal que $\hat{A}\ket{n} = a_n\ket{n}$ y $\hat{B}\ket{n} = b_n\ket{n}$ para todo $n$, entonces:
\begin{equation}
    [\hat{A}, \hat{B}]\ket{n} = (\hat{A}\hat{B} - \hat{B}\hat{A})\ket{n} = (a_n b_n - b_n a_n)\ket{n} = 0.
    \label{a1-eq:comm_necessary}
\end{equation}
Como los autovalores escalares conmutan ($a_n b_n = b_n a_n$) y $\{\ket{n}\}$ es una base completa que genera todo $\mathcal{H}$, el operador $[\hat{A}, \hat{B}]$ anula cualquier combinación lineal $\ket{\psi} = \sum_n c_n\ket{n}$, concluyéndose idénticamente que $[\hat{A}, \hat{B}] = 0$.

\textbf{Demostración de la Condición Suficiente ($\Rightarrow$):}\\
Para probar que $[\hat{A}, \hat{B}] = 0$ implica la existencia de una base común y prepararnos para resolver problemas con operadores compatibles, debemos distinguir con absoluto rigor didáctico los dos escenarios posibles: el caso sin degeneración y el caso con degeneración.

\subsubsection{Caso 1: Espectro No Degenerado ($g = 1$)}

Supongamos primero que todos los autovalores del operador $\hat{A}$ son simples (no degenerados), es decir, a cada autovalor $a_n$ le corresponde un subespacio propio unidimensional generado por el autovector $\ket{\phi_n}$:
\begin{equation}
    \hat{A}\ket{\phi_n} = a_n\ket{\phi_n}, \quad \text{con } a_n \neq a_m \text{ para } n \neq m.
\end{equation}
Apliquemos el operador $\hat{B}$ sobre el vector $\ket{\phi_n}$ y examinemos qué le hace el operador $\hat{A}$ al vector resultante $\hat{B}\ket{\phi_n}$. Como $[\hat{A}, \hat{B}] = 0$:
\begin{equation}
    \hat{A}(\hat{B}\ket{\phi_n}) = (\hat{A}\hat{B})\ket{\phi_n} = (\hat{B}\hat{A})\ket{\phi_n} = \hat{B}(\hat{A}\ket{\phi_n}) = \hat{B}(a_n\ket{\phi_n}) = a_n(\hat{B}\ket{\phi_n}).
    \label{eq:conmutacion_no_deg}
\end{equation}
¡Observemos con atención la Ec.~\eqref{eq:conmutacion_no_deg}! Nos dice que el vector $\ket{\chi} \equiv \hat{B}\ket{\phi_n}$ es un autovector de $\hat{A}$ con el mismo autovalor $a_n$.\\
Pero como el subespacio propio $E_{a_n}$ asociado al autovalor $a_n$ es unidimensional (no degenerado), \textbf{todo vector en dicho subespacio debe ser forzosamente un múltiplo escalar del vector base $\ket{\phi_n}$}:
\begin{equation}
    \hat{B}\ket{\phi_n} = b_n\ket{\phi_n}, \quad \text{para algún escalar } b_n \in \mathbb{C}.
\end{equation}
Como $\hat{B}$ es hermítico, el escalar $b_n$ es necesariamente real ($b_n \in \mathbb{R}$). Por lo tanto, el vector $\ket{\phi_n}$ es simultáneamente autovector de $\hat{A}$ (con autovalor $a_n$) y de $\hat{B}$ (con autovalor $b_n$):
\begin{equation}
    \boxed{\hat{A}\ket{\phi_n} = a_n\ket{\phi_n} \quad \text{y} \quad \hat{B}\ket{\phi_n} = b_n\ket{\phi_n}}.
\end{equation}
En el caso no degenerado, la base ortonormal común es \textbf{única} (salvo fases globales de cada vector).

\subsubsection{Caso 2: Espectro Degenerado ($g > 1$) y Subespacios Invariantes}

¿Qué ocurre si el operador $\hat{A}$ presenta degeneración espectral? Supongamos que a un determinado autovalor $a_n$ le corresponde un subespacio propio $E_{a_n}$ de dimensión $g_n > 1$ (multiplicidad algebraica/geométrica):
\begin{equation}
    E_{a_n} \equiv \left\{ \ket{\psi} \in \mathcal{H} : \hat{A}\ket{\psi} = a_n\ket{\psi} \right\}, \quad \dim(E_{a_n}) = g_n > 1.
\end{equation}
Sea $\{\ket{\phi_{n, i}}\}_{i=1}^{g_n}$ una base ortonormal cualquiera elegida para dicho subespacio.
\begin{enumerate}
    \item \textbf{Invariancia del subespacio propio bajo $\hat{B}$:}
    Si $\ket{\psi} \in E_{a_n}$, entonces por la conmutación $[\hat{A}, \hat{B}] = 0$:
    \begin{equation}
        \hat{A}(\hat{B}\ket{\psi}) = \hat{B}\hat{A}\ket{\psi} = \hat{B}(a_n\ket{\psi}) = a_n(\hat{B}\ket{\psi}) \implies \hat{B}\ket{\psi} \in E_{a_n}.
    \end{equation}
    Esto demuestra un hecho geométrico fundamental: \textbf{el operador $\hat{B}$ no saca los vectores de $E_{a_n}$; el subespacio $E_{a_n}$ es estrictamente invariante bajo $\hat{B}$}.
    \item \textbf{¿Es $\ket{\phi_{n, i}}$ automáticamente un autovector de $\hat{B}$? ¡NO!}\\
    En general, la acción de $\hat{B}$ sobre un vector base $\ket{\phi_{n, i}}$ es una combinación lineal arbitraria de los vectores que forman la base de $E_{a_n}$:
    \begin{equation}
        \hat{B}\ket{\phi_{n, i}} = \sum_{j=1}^{g_n} B_{ji}^{(n)} \ket{\phi_{n, j}} \neq \text{constante} \times \ket{\phi_{n, i}}.
    \end{equation}
    Por consiguiente, un autovector genérico de $\hat{A}$ dentro del subespacio degenerado \textbf{no es necesariamente un autovector de $\hat{B}$}.
    \item \textbf{Diagonalización interna en el bloque degenerado:}\\
    La restricción de $\hat{B}$ al subespacio $E_{a_n}$ (que denotamos por $\hat{B}|_{E_{a_n}}$) es un operador hermítico que actúa en un espacio de dimensión finita $g_n$. Por el Teorema 2.2, dicho operador restringido \textbf{puede ser diagonalizado de forma independiente dentro de $E_{a_n}$}.\\
    Existen exactamente $g_n$ combinaciones lineales ortonormales $\{\ket{\chi_{n, k}}\}_{k=1}^{g_n} \subset E_{a_n}$ tales que:
    \begin{equation}
        \hat{B}\ket{\chi_{n, k}} = b_{n, k}\ket{\chi_{n, k}}, \quad k \in \{1, 2, \dots, g_n\}.
    \end{equation}
    Y como cada $\ket{\chi_{n, k}}$ pertenece al subespacio $E_{a_n}$, automáticamente satisface:
    \begin{equation}
        \hat{A}\ket{\chi_{n, k}} = a_n \ket{\chi_{n, k}}.
    \end{equation}
    ¡Los vectores $\{\ket{\chi_{n, k}}\}$ son los autovectores comunes buscados!
    \item \textbf{No unicidad de la base:}\\
    Si algunos autovalores $b_{n, k}$ del operador $\hat{B}$ vuelven a coincidir (degeneración en $\hat{B}$), la base común simultánea \textbf{tampoco es única}. Se requerirá añadir un tercer observable compatible $\hat{C}$ ($[\hat{A}, \hat{C}] = [\hat{B}, \hat{C}] = 0$) para terminar de romper la degeneración y etiquetar cada estado de manera no ambigua. ¡Este es exactamente el origen del concepto de \textbf{Conjunto Completo de Observables Compatibles (CCOC)}!
\end{enumerate}

\subsubsection{Ejemplo Matricial Guiado: Diagonalización Simultánea con Degeneración}

Para consolidar la técnica de diagonalización simultánea, desarrollemos paso a paso un problema canónico.

\begin{ejemplobox}[Ejemplo Resuelto: Diagonalización Simultánea de Matrices $3\times 3$ con Degeneración]
    \textbf{Enunciado:} Considere dos operadores lineales representados en la base canónica de $\mathbb{R}^3$ por las matrices simétricas reales:
    \begin{equation}
        A = \begin{pmatrix} 1 & 0 & 0 \\ 0 & 1 & 0 \\ 0 & 0 & -2 \end{pmatrix}, \qquad B = \begin{pmatrix} 0 & 2 & 0 \\ 2 & 0 & 0 \\ 0 & 0 & 3 \end{pmatrix}.
    \end{equation}
    \begin{enumerate}
        \item Verifique explícitamente que $A$ y $B$ conmutan ($[A, B] = 0$).
        \item Halle los autovalores de $A$, identifique el subespacio propio degenerado $E_1$ y explique por qué la base canónica ya no es autovectorial para $B$.
        \item Restrinja $B$ al subespacio degenerado $E_1$, diagonalícelo y construya la matriz unitaria (ortogonal) de cambio de base $P$ que diagonaliza simultáneamente a ambas matrices.
    \end{enumerate}

    \textbf{Solución Paso a Paso:}

    \textit{Paso 1: Verificación de la Conmutación:}
    Efectuamos el producto matricial en ambos órdenes:
    \begin{align}
        AB &= \begin{pmatrix} 1 & 0 & 0 \\ 0 & 1 & 0 \\ 0 & 0 & -2 \end{pmatrix} \begin{pmatrix} 0 & 2 & 0 \\ 2 & 0 & 0 \\ 0 & 0 & 3 \end{pmatrix} = \begin{pmatrix} 0 & 2 & 0 \\ 2 & 0 & 0 \\ 0 & 0 & -6 \end{pmatrix}, \\
        BA &= \begin{pmatrix} 0 & 2 & 0 \\ 2 & 0 & 0 \\ 0 & 0 & 3 \end{pmatrix} \begin{pmatrix} 1 & 0 & 0 \\ 0 & 1 & 0 \\ 0 & 0 & -2 \end{pmatrix} = \begin{pmatrix} 0 & 2 & 0 \\ 2 & 0 & 0 \\ 0 & 0 & -6 \end{pmatrix}.
    \end{align}
    Puesto que $AB = BA$, se verifica idénticamente que $[A, B] = 0$.

    \textit{Paso 2: Espectro de $A$ y Subespacios Propios:}
    Como la matriz $A$ ya es diagonal, sus autovalores se leen directamente de la diagonal principal:
    \begin{equation}
        \lambda_1 = \lambda_2 = 1 \;\; (\text{deg., } g=2), \qquad \lambda_3 = -2 \;\; (\text{no deg., } g=1).
    \end{equation}
    El autovector asociado al autovalor no degenerado $\lambda_3 = -2$ es único (salvo fase):
    \begin{equation}
        \ket{v_3} = \begin{pmatrix} 0 \\ 0 \\ 1 \end{pmatrix}.
    \end{equation}
    El subespacio degenerado asociado al autovalor doble $\lambda = 1$ es el plano generado por los dos primeros vectores de la base canónica:
    \begin{equation}
        E_1 = \mathrm{span}\left\{ \ket{e_1} = \begin{pmatrix} 1 \\ 0 \\ 0 \end{pmatrix}, \, \ket{e_2} = \begin{pmatrix} 0 \\ 1 \\ 0 \end{pmatrix} \right\}.
    \end{equation}
    Notemos con absoluta claridad pedagógica: aunque $\ket{e_1}$ y $\ket{e_2}$ son autovectores de $A$ con autovalor $1$, \textbf{ninguno de ellos es autovector de $B$}, pues:
    \begin{equation}
        B\ket{e_1} = \begin{pmatrix} 0 & 2 & 0 \\ 2 & 0 & 0 \\ 0 & 0 & 3 \end{pmatrix}\begin{pmatrix} 1 \\ 0 \\ 0 \end{pmatrix} = \begin{pmatrix} 0 \\ 2 \\ 0 \end{pmatrix} = 2\ket{e_2} \neq \text{constante} \times \ket{e_1}.
    \end{equation}
    El operador $B$ \enquote{mezcla} los vectores dentro del plano $E_1$.

    \textit{Paso 3: Restricción y Diagonalización en el Bloque Degenerado:}
    La matriz de $B$ restringida al subespacio bidimensional $E_1 = \mathrm{span}\{\ket{e_1}, \ket{e_2}\}$ viene dada por el bloque superior izquierdo de $2\times 2$:
    \begin{equation}
        B|_{E_1} = \begin{pmatrix} 0 & 2 \\ 2 & 0 \end{pmatrix}.
    \end{equation}
    Hallamos sus autovalores resolviendo el polinomio característico secular:
    \begin{equation}
        \det(B|_{E_1} - \mu I_2) = \begin{vmatrix} -\mu & 2 \\ 2 & -\mu \end{vmatrix} = \mu^2 - 4 = 0 \implies \mu = \pm 2.
    \end{equation}
    Calculamos los autovectores correspondientes en el subespacio $E_1$:
    \begin{itemize}
        \item Para $\mu_1 = +2$:
        \begin{equation}
            \begin{pmatrix} -2 & 2 \\ 2 & -2 \end{pmatrix}\begin{pmatrix} c_1 \\ c_2 \end{pmatrix} = \begin{pmatrix} 0 \\ 0 \end{pmatrix} \implies c_1 = c_2 \implies \ket{u_1} = \frac{1}{\sqrt{2}}\begin{pmatrix} 1 \\ 1 \\ 0 \end{pmatrix}.
        \end{equation}
        \item Para $\mu_2 = -2$:
        \begin{equation}
            \begin{pmatrix} 2 & 2 \\ 2 & 2 \end{pmatrix}\begin{pmatrix} c_1 \\ c_2 \end{pmatrix} = \begin{pmatrix} 0 \\ 0 \end{pmatrix} \implies c_1 = -c_2 \implies \ket{u_2} = \frac{1}{\sqrt{2}}\begin{pmatrix} 1 \\ -1 \\ 0 \end{pmatrix}.
        \end{equation}
    \end{itemize}
    Notemos que $\ket{u_1}$ y $\ket{u_2}$ siguen perteneciendo a $E_1$, por lo que $A\ket{u_1} = 1\ket{u_1}$ y $A\ket{u_2} = 1\ket{u_2}$.\\
    Para el tercer autovector, evaluamos la acción de $B$ sobre $\ket{v_3}$:
    \begin{equation}
        B\ket{v_3} = \begin{pmatrix} 0 & 2 & 0 \\ 2 & 0 & 0 \\ 0 & 0 & 3 \end{pmatrix}\begin{pmatrix} 0 \\ 0 \\ 1 \end{pmatrix} = \begin{pmatrix} 0 \\ 0 \\ 3 \end{pmatrix} = 3\ket{v_3}.
    \end{equation}
    Por tanto, $\ket{u_3} \equiv \ket{v_3} = \begin{pmatrix} 0 \\ 0 \\ 1 \end{pmatrix}$ ya era autovector de $B$ con autovalor $\mu_3 = 3$.

    \textit{Paso 4: Construcción de la Matriz de Paso $P$ y Verificación Final:}
    Colocamos los autovectores comunes $\{\ket{u_1}, \ket{u_2}, \ket{u_3}\}$ como columnas de la matriz ortogonal de transformación $P$:
    \begin{equation}
        P = \begin{pmatrix}
            \frac{1}{\sqrt{2}} & \frac{1}{\sqrt{2}} & 0 \\
            \frac{1}{\sqrt{2}} & -\frac{1}{\sqrt{2}} & 0 \\
            0 & 0 & 1
        \end{pmatrix}.
    \end{equation}
    Puesto que las columnas son ortonormales, $P$ es ortogonal ($P^{-1} = P^T$). Comprobemos que diagonaliza simultáneamente a ambas matrices:
    \begin{align}
        P^{-1} A P &= P^T A P = \begin{pmatrix}
            \frac{1}{\sqrt{2}} & \frac{1}{\sqrt{2}} & 0 \\
            \frac{1}{\sqrt{2}} & -\frac{1}{\sqrt{2}} & 0 \\
            0 & 0 & 1
        \end{pmatrix} \begin{pmatrix} 1 & 0 & 0 \\ 0 & 1 & 0 \\ 0 & 0 & -2 \end{pmatrix} \begin{pmatrix}
            \frac{1}{\sqrt{2}} & \frac{1}{\sqrt{2}} & 0 \\
            \frac{1}{\sqrt{2}} & -\frac{1}{\sqrt{2}} & 0 \\
            0 & 0 & 1
        \end{pmatrix} = \begin{pmatrix} 1 & 0 & 0 \\ 0 & 1 & 0 \\ 0 & 0 & -2 \end{pmatrix}, \\
        P^{-1} B P &= P^T B P = \begin{pmatrix}
            \frac{1}{\sqrt{2}} & \frac{1}{\sqrt{2}} & 0 \\
            \frac{1}{\sqrt{2}} & -\frac{1}{\sqrt{2}} & 0 \\
            0 & 0 & 1
        \end{pmatrix} \begin{pmatrix} 0 & 2 & 0 \\ 2 & 0 & 0 \\ 0 & 0 & 3 \end{pmatrix} \begin{pmatrix}
            \frac{1}{\sqrt{2}} & \frac{1}{\sqrt{2}} & 0 \\
            \frac{1}{\sqrt{2}} & -\frac{1}{\sqrt{2}} & 0 \\
            0 & 0 & 1
        \end{pmatrix} = \begin{pmatrix} 2 & 0 & 0 \\ 0 & -2 & 0 \\ 0 & 0 & 3 \end{pmatrix}.
    \end{align}
    ¡Ambas matrices quedan rigurosamente diagonalizadas al mismo tiempo en la base $\{\ket{u_1}, \ket{u_2}, \ket{u_3}\}$!
\end{ejemplobox}

\begin{figure}[H]
\centering
\begin{tikzpicture}[
    x={(2.1cm,0cm)}, y={(1.0cm,0.8cm)}, z={(0cm,2.0cm)},
    >=Stealth, font=\small]

  % plano degenerado E_1
  \fill[azulUNI!10] (-1.1,-1.1,0)--(1.1,-1.1,0)--(1.1,1.1,0)--(-1.1,1.1,0)--cycle;
  \draw[azulUNI!45] (-1.1,-1.1,0)--(1.1,-1.1,0)--(1.1,1.1,0)--(-1.1,1.1,0)--cycle;
  \node[azulUNI,anchor=south east,font=\small] at (-0.85,1.1,0)
        {$E_1$ \; ($A=1$, degenerado)};

  % base canónica (no es propia de B)
  \draw[->,thick,grisUNI] (0,0,0)--(1,0,0) node[below right]{$|e_1\rangle$};
  \draw[->,thick,grisUNI] (0,0,0)--(0,1,0) node[above left]{$|e_2\rangle$};

  % base propia común
  \draw[->,very thick,rojoUNI]  (0,0,0)--(0.85,0.85,0)
        node[above right]{$|u_1\rangle$};
  \draw[->,very thick,azulUNI]  (0,0,0)--(0.85,-0.85,0)
        node[below]{$|u_2\rangle$};
  \draw[->,very thick,verdeUNI] (0,0,0)--(0,0,1.4)
        node[above]{$|v_3\rangle$};

  % B mezcla e_1 y e_2
  \draw[<->,dashed,grisUNI] (0.55,0,0) to[bend left=35] (0,0.55,0);

  % leyenda
  \node[align=left,anchor=north west,font=\small] at (canvas cs:x=-3.6cm,y=-1.7cm) {%
    \textcolor{rojoUNI}{$|u_1\rangle=\tfrac{1}{\sqrt2}(|e_1\rangle+|e_2\rangle)$: $A=1,\;B=+2$}\\[2pt]
    \textcolor{azulUNI}{$|u_2\rangle=\tfrac{1}{\sqrt2}(|e_1\rangle-|e_2\rangle)$: $A=1,\;B=-2$}\\[2pt]
    \textcolor{verdeUNI}{$|v_3\rangle=|e_3\rangle$: $A=-2,\;B=3$}\\[6pt]
    \textcolor{grisUNI}{$B|e_1\rangle=2|e_2\rangle$\; (\,$B$ «mezcla» dentro de $E_1$)}};
\end{tikzpicture}
\caption{Diagonalización simultánea con degeneración (Ejemplo de la Sección~2.8.3). Dentro del
plano $E_1$, $A$ es un múltiplo de la identidad, pero $B$ gira los vectores:
la base canónica $\{|e_1\rangle,|e_2\rangle\}$ no es propia de $B$. La base
común se obtiene diagonalizando $B$ restringido a $E_1$, lo que equivale a
girar $45^\circ$ dentro del plano; $|v_3\rangle$ ya era propio de ambos.}
\label{fig:diag-simultanea}
\end{figure}


\subsubsection{Ejemplos Físicos de Conjuntos Simultáneamente Diagonalizables}

Este procedimiento algebraico se manifiesta en sistemas físicos fundamentales de la física cuántica:
\begin{enumerate}
    \item \textbf{Partícula Libre en 3D:} Las tres componentes cartesianas del momento lineal conmutan entre sí:
    \begin{equation}
        [\hat{p}_x, \hat{p}_y] = [\hat{p}_y, \hat{p}_z] = [\hat{p}_z, \hat{p}_x] = 0, \quad \text{y además } [\hat{H}, \hat{p}_i] = 0 \quad \text{con } \hat{H} = \frac{\hat{p}^2}{2m}.
    \end{equation}
    Admiten como base común continua de autoestados a las \textbf{ondas planas tridimensionales}:
    \begin{equation}
        \psi_{\vec{p}}(\vec{r}) = \braket{\vec{r}}{\vec{p}} = \frac{1}{(2\pi\hbar)^{3/2}}\eu^{\frac{\iu}{\hbar}\vec{p}\cdot\vec{r}}, \quad \text{con } \hat{\vec{p}}\ket{\vec{p}} = \vec{p}\ket{\vec{p}} \quad \text{y} \quad \hat{H}\ket{\vec{p}} = \frac{p^2}{2m}\ket{\vec{p}}.
    \end{equation}
    \item \textbf{Oscilador Armónico 3D Isótropo Separable:} El Hamiltoniano total es la suma de tres osciladores 1D independientes:
    \begin{equation}
        \hat{H} = \hat{H}_x + \hat{H}_y + \hat{H}_z = \hbar\omega\left(\hat{a}_x\adj\hat{a}_x + \hat{a}_y\adj\hat{a}_y + \hat{a}_z\adj\hat{a}_z + \frac{3}{2}\right).
    \end{equation}
    Dado que $[\hat{H}_x, \hat{H}_y] = [\hat{H}_y, \hat{H}_z] = [\hat{H}_z, \hat{H}_x] = 0$, la base ortonormal común viene dada por los productos directos de autoestados de Fock unidimensionales:
    \begin{equation}
        \ket{n_x, n_y, n_z} = \ket{n_x} \otimes \ket{n_y} \otimes \ket{n_z}, \quad \text{con energía } E = \hbar\omega\left(n_x + n_y + n_z + \frac{3}{2}\right).
    \end{equation}
\end{enumerate}

\subsection{Taller de Conmutadores Fundamentales, Par Canónico y Operadores Escalera}
\label{subsec:taller_conmutadores}

Antes de abordar las demostraciones avanzadas de las secciones siguientes, resulta indispensable dominar la \enquote{carpintería algebraica} de los conmutadores, comprender por qué la posición y el momento no conmutan, y conocer el álgebra de los operadores de creación y aniquilación.

\subsubsection{Las Cuatro Reglas Operativas del Álgebra de Conmutadores}

Dados tres operadores lineales cualesquiera $\hat{A}, \hat{B}, \hat{C}$ y escalares $\alpha, \beta \in \mathbb{C}$, el conmutador $[\hat{A}, \hat{B}] \equiv \hat{A}\hat{B} - \hat{B}\hat{A}$ satisface:
\begin{enumerate}
    \item \textbf{Antisimetría:} Intercambiar el orden de los operadores invierte el signo:
    \begin{equation}
        [\hat{A}, \hat{B}] = -[\hat{B}, \hat{A}], \quad \text{en particular } [\hat{A}, \hat{A}] = 0 \text{ y } [\hat{A}, c\hat{I}] = 0.
    \end{equation}
    \item \textbf{Bilinealidad:} Es lineal respecto a la suma y multiplicación por escalares en ambos argumentos:
    \begin{equation}
        [\alpha\hat{A} + \beta\hat{B}, \, \hat{C}] = \alpha[\hat{A}, \hat{C}] + \beta[\hat{B}, \hat{C}], \qquad [\hat{A}, \, \alpha\hat{B} + \beta\hat{C}] = \alpha[\hat{A}, \hat{B}] + \beta[\hat{A}, \hat{C}].
    \end{equation}
    \item \textbf{Regla de Leibniz (Derivación de un Producto):} Permite descomponer el conmutador de un producto en conmutadores simples:
    \begin{equation}
        \boxed{[\hat{A}\hat{B}, \hat{C}] = \hat{A}[\hat{B}, \hat{C}] + [\hat{A}, \hat{C}]\hat{B}}, \qquad
        \boxed{[\hat{A}, \hat{B}\hat{C}] = [\hat{A}, \hat{B}]\hat{C} + \hat{B}[\hat{A}, \hat{C}]}.
        \label{eq:leibniz_comm_basic}
    \end{equation}
    \textit{Demostración directa:} Expandiendo el miembro derecho de la primera identidad y observando la cancelación telescópica del término intermedio:
    \begin{equation}
        \hat{A}[\hat{B}, \hat{C}] + [\hat{A}, \hat{C}]\hat{B} = \hat{A}\hat{B}\hat{C} - \bcancel{\hat{A}\hat{C}\hat{B}} + \bcancel{\hat{A}\hat{C}\hat{B}} - \hat{C}\hat{A}\hat{B} = \hat{A}\hat{B}\hat{C} - \hat{C}\hat{A}\hat{B} = [\hat{A}\hat{B}, \hat{C}].
    \end{equation}
    \item \textbf{Identidad de Jacobi:} La suma cíclica de conmutadores anidados es idénticamente nula (propiedad definitoria de toda álgebra de Lie):
    \begin{equation}
        [\hat{A}, [\hat{B}, \hat{C}]] + [\hat{B}, [\hat{C}, \hat{A}]] + [\hat{C}, [\hat{A}, \hat{B}]] = 0.
    \end{equation}
\end{enumerate}

\subsubsection{¿De dónde sale \texorpdfstring{$[\hat{x}, \hat{p}_x] = \iu\hbar\hat{I}$}{[x, px] = ihbar I}? Demostración sobre una Función de Prueba}

En la representación de coordenadas espaciales, el operador posición actúa multiplicando por $x$ ($\hat{x}\psi(x) = x\psi(x)$) y el operador momento lineal actúa derivando ($\hat{p}_x\psi(x) = -\iu\hbar\frac{\dif\psi(x)}{\dif x}$).

\begin{alertabox}[Regla de Oro para Evaluar Conmutadores Diferenciales]
    Nunca evalúe un operador diferencial \enquote{en el vacío}. Para hallar $[\hat{x}, \hat{p}_x]$, hágalo actuar sobre una función de onda diferenciable arbitraria (\textbf{función de prueba}) $\psi(x)$ y aplique la regla de la derivada de un producto:
    \begin{align}
        [\hat{x}, \hat{p}_x]\psi(x) &= (\hat{x}\hat{p}_x - \hat{p}_x\hat{x})\psi(x) = x\left(-\iu\hbar\frac{\dif\psi}{\dif x}\right) - \left(-\iu\hbar\frac{\dif}{\dif x}\right)\big(x\psi(x)\big) \nonumber\\
        &= -\iu\hbar\,x\frac{\dif\psi}{\dif x} + \iu\hbar\left( \frac{\dif x}{\dif x}\psi(x) + x\frac{\dif\psi}{\dif x} \right) \nonumber\\
        &= -\bcancel{\iu\hbar\,x\frac{\dif\psi}{\dif x}} + \iu\hbar\,\psi(x) + \bcancel{\iu\hbar\,x\frac{\dif\psi}{\dif x}} = \iu\hbar\,\psi(x).
    \end{align}
    Como esta igualdad se verifica para cualquier función $\psi(x) \in \mathcal{H}$, suprimiendo la función de prueba a ambos lados se obtiene la célebre \textbf{Relación de Conmutación Canónica de Born--Jordan--Heisenberg}:
    \begin{equation}
        \boxed{[\hat{x}, \hat{p}_x] = \iu\hbar\hat{I}}, \quad \text{y en 3D:} \quad \boxed{[\hat{x}_i, \hat{p}_j] = \iu\hbar\,\delta_{ij}\hat{I}}.
    \end{equation}
\end{alertabox}

\subsubsection{Operadores Escalera del Oscilador Armónico (\texorpdfstring{$\hat{a}, \hat{a}\adj$}{a, a+}) y Base de Fock}

En múltiples problemas cuánticos (como el Oscilador 2D del Problema Resuelto 5), el Hamiltoniano armónico unidimensional $\hat{H} = \frac{\hat{p}^2}{2m} + \frac{1}{2}m\omega^2\hat{x}^2$ se factoriza algebraicamente definiendo los operadores adimensionales de \textbf{aniquilación} ($\hat{a}$) y \textbf{creación} ($\hat{a}\adj$):
\begin{equation}
    \hat{a} \equiv \sqrt{\frac{m\omega}{2\hbar}}\left(\hat{x} + \frac{\iu}{m\omega}\hat{p}\right), \qquad
    \hat{a}\adj \equiv \sqrt{\frac{m\omega}{2\hbar}}\left(\hat{x} - \frac{\iu}{m\omega}\hat{p}\right).
\end{equation}
Utilizando $[\hat{x}, \hat{p}] = \iu\hbar\hat{I}$, su álgebra fundamental se resume en cuatro relaciones esenciales:
\begin{enumerate}
    \item \textbf{Conmutador bosónico fundamental:} $[\hat{a}, \hat{a}\adj] = \hat{I}$ (y para ejes distintos en 2D/3D, $[\hat{a}_i, \hat{a}_j\adj] = \delta_{ij}\hat{I}$, $[\hat{a}_i, \hat{a}_j] = 0$).
    \item \textbf{Operador Número y Hamiltoniano:} Definiendo el operador hermítico número de cuantos $\hat{N} \equiv \hat{a}\adj\hat{a}$, el Hamiltoniano se reescribe como $\hat{H} = \hbar\omega\left(\hat{N} + \frac{1}{2}\hat{I}\right)$, con conmutadores $[\hat{N}, \hat{a}] = -\hat{a}$ y $[\hat{N}, \hat{a}\adj] = +\hat{a}\adj$.
    \item \textbf{Acción sobre los autoestados de Fock $\{\ket{n}\}$ ($\hat{N}\ket{n} = n\ket{n}$, $n=0,1,2,\dots$):}
    \begin{equation}
        \boxed{\hat{a}\ket{n} = \sqrt{n}\,\ket{n-1} \quad (\text{con } \hat{a}\ket{0} = 0), \qquad \hat{a}\adj\ket{n} = \sqrt{n+1}\,\ket{n+1}}.
    \end{equation}
\end{enumerate}

%% ============================================================
%%  SECCIÓN 3 — Principio Generalizado de Robertson-Schrödinger
%% ============================================================

\section{El Principio Generalizado de Incertidumbre de Robertson--\texorpdfstring{Schrödinger}{Schrodinger}}
\label{sec:incertidumbre}

\subsection{La Desigualdad de Cauchy--Schwarz en el Espacio de Hilbert}

El fundamento geométrico de todas las relaciones de indeterminación cuántica reside en la desigualdad de Cauchy--Schwarz, la cual constituye el análogo en espacios vectoriales complejos de la propiedad geométrica elemental $|\vec{a}\cdot\vec{b}| \le \norm{\vec{a}}\norm{\vec{b}}$ del espacio euclídeo \cite{zettili}.

\begin{teoremabox}[Teorema: Desigualdad de Cauchy--Schwarz en $\mathcal{H}$]
    Para cualesquiera dos vectores $\ket{u}, \ket{v} \in \mathcal{H}$, se cumple estrictamente:
    \begin{equation}
        |\braket{u}{v}|^2 \le \braket{u}{u}\braket{v}{v} = \norm{u}^2 \norm{v}^2,
        \label{eq:cauchy_schwarz}
    \end{equation}
    alcanzándose la igualdad si y solo si los vectores son linealmente dependientes (es decir, $\ket{u} = c\ket{v}$ para algún escalar $c \in \mathbb{C}$, o alguno de ellos es el vector nulo).
\end{teoremabox}

\textbf{Demostración Rigurosa Sin Atajos:}\\
Si $\ket{v} = 0$, la desigualdad se cumple trivialmente ($0 \le 0$). Supongamos entonces que $\ket{v} \neq 0$, lo que implica $\braket{v}{v} > 0$ por el axioma de positividad definida. Consideremos un vector auxiliar $\ket{\phi}$ construido como una combinación lineal con un parámetro complejo arbitrario $\lambda \in \mathbb{C}$:
\begin{equation}
    \ket{\phi} = \ket{u} + \lambda\ket{v}.
\end{equation}
Por el axioma de positividad definida del producto interno, la norma al cuadrado de $\ket{\phi}$ no puede ser negativa:
\begin{align}
    0 \le \braket{\phi}{\phi} &= \braket{u + \lambda v}{u + \lambda v} \nonumber\\
    &= \braket{u}{u} + \lambda\braket{u}{v} + \lambda^*\braket{v}{u} + \lambda^*\lambda\braket{v}{v} \nonumber\\
    &= \braket{u}{u} + \lambda\braket{u}{v} + \lambda^*\braket{u}{v}^* + |\lambda|^2\braket{v}{v}.
    \label{eq:cs_expansion}
\end{align}
Dado que esta desigualdad debe ser válida para \textit{cualquier} número complejo $\lambda$, elegimos el valor que minimiza la forma cuadrática, el cual corresponde a la proyección ortogonal de $\ket{u}$ sobre $\ket{v}$:
\begin{equation}
    \lambda = -\frac{\braket{v}{u}}{\braket{v}{v}}.
    \label{eq:optimal_lambda}
\end{equation}
Tomando el conjugado complejo de $\lambda$:
\begin{equation}
    \lambda^* = -\frac{\braket{v}{u}^*}{\braket{v}{v}^*} = -\frac{\braket{u}{v}}{\braket{v}{v}}.
\end{equation}
Sustituyendo cuidadosamente $\lambda$ y $\lambda^*$ en la Ec.~\eqref{eq:cs_expansion}:
\begin{align}
    0 &\le \braket{u}{u} + \left(-\frac{\braket{v}{u}}{\braket{v}{v}}\right)\braket{u}{v} + \left(-\frac{\braket{u}{v}}{\braket{v}{v}}\right)\braket{v}{u} + \left|\frac{\braket{v}{u}}{\braket{v}{v}}\right|^2 \braket{v}{v} \nonumber\\
    &= \braket{u}{u} - \frac{\braket{v}{u}\braket{u}{v}}{\braket{v}{v}} - \frac{\braket{u}{v}\braket{v}{u}}{\braket{v}{v}} + \frac{|\braket{v}{u}|^2}{\braket{v}{v}^2}\braket{v}{v}.
\end{align}
Recordando que $\braket{v}{u}\braket{u}{v} = \braket{u}{v}^*\braket{u}{v} = |\braket{u}{v}|^2$, los términos se simplifican algebraicamente:
\begin{equation}
    0 \le \braket{u}{u} - \frac{|\braket{u}{v}|^2}{\braket{v}{v}} - \frac{|\braket{u}{v}|^2}{\braket{v}{v}} + \frac{|\braket{u}{v}|^2}{\braket{v}{v}} = \braket{u}{u} - \frac{|\braket{u}{v}|^2}{\braket{v}{v}}.
\end{equation}
Multiplicando toda la desigualdad por el número real positivo $\braket{v}{v} > 0$:
\begin{equation}
    0 \le \braket{u}{u}\braket{v}{v} - |\braket{u}{v}|^2 \implies \boxed{|\braket{u}{v}|^2 \le \braket{u}{u}\braket{v}{v}}.
\end{equation}
Queda formalmente demostrada la desigualdad de Cauchy--Schwarz sin recurrir a derivadas de funciones reales ni asumir que el producto interno sea real \cite{sakurai}. Nótese que la igualdad en la Ec.~\eqref{eq:cauchy_schwarz} se alcanza si y solo si $\braket{\phi}{\phi} = 0 \iff \ket{\phi} = 0$, lo cual equivale a la dependencia lineal $\ket{u} = -\lambda\ket{v}$ entre ambos estados (con el escalar óptimo $\lambda = -\braket{v}{u}/\braket{v}{v}$).

\begin{figure}[H]
\centering
\begin{tikzpicture}[scale=1.15]
  \coordinate (O) at (0,0);
  \coordinate (U) at (3.5,2.5);
  \coordinate (V) at (5,1);
  \coordinate (P) at (3.846,0.769);
  \draw[-{Stealth},very thick,uniblue] (O) -- (V) node[right] {$\ket v$};
  \draw[-{Stealth},very thick,unired] (O) -- (U) node[above] {$\ket u$};
  \draw[-{Stealth},very thick,unigreen] (O) -- (P)
        node[midway,below,font=\small] {$\dfrac{\braket vu}{\braket vv}\ket v$};
  \draw[dashed,thick,black] (P) -- (U)
        node[midway,left=3pt,font=\small]
        {$\ket\phi$};
  \draw (4.140,0.828) -- (4.081,1.122) -- (3.787,1.063);
  \draw pic["$\vartheta$",draw=unigray,angle radius=1.1cm,angle eccentricity=1.35]
        {angle=V--O--U};
  \node[font=\small,align=left] at (7.2,3.0)
    {$\ket\phi=\ket u-\dfrac{\braket vu}{\braket vv}\ket v$\\[4pt]
     $\braket\phi v=0$ \\[2pt] $\Rightarrow\ \braket\phi\phi\ge0$ \\[2pt]
     $\Rightarrow\ |\braket uv|^2\le\braket uu\braket vv$};
\end{tikzpicture}
\caption{Cauchy--Schwarz: $\lambda=-\braket vu/\braket vv$ hace que $\ket\phi=\ket u+\lambda\ket v$ sea el residuo ortogonal a $\ket v$; su norma no puede ser negativa.}
\label{fig:cauchy_schwarz}
\end{figure}

\subsection{Construcción de los Estados de Desviación}

Sean $\hat{A}$ y $\hat{B}$ dos observables físicos representados por operadores lineales hermíticos ($\hat{A}\adj = \hat{A}$, $\hat{B}\adj = \hat{B}$). Supongamos que el sistema físico se encuentra en un estado normalizado arbitrario $\ket{\psi}$ ($\braket{\psi}{\psi}=1$).

Definimos los estados desviados (o fluctuaciones de estado) aplicando los operadores desviación definidos en la Ec.~\eqref{eq:deviation_operator}:
\begin{equation}
    \ket{u} \equiv \Delta\hat{A}\ket{\psi} = (\hat{A} - \ev{\hat{A}}\hat{I})\ket{\psi},
    \label{eq:state_u}
\end{equation}
\begin{equation}
    \ket{v} \equiv \Delta\hat{B}\ket{\psi} = (\hat{B} - \ev{\hat{B}}\hat{I})\ket{\psi}.
    \label{eq:state_v}
\end{equation}
Evaluemos ahora las normas al cuadrado de estos vectores:
\begin{equation}
    \braket{u}{u} = \braket{\Delta\hat{A}\psi}{\Delta\hat{A}\psi} = \matrixel{\psi}{(\Delta\hat{A})\adj\Delta\hat{A}}{\psi}.
\end{equation}
Debido a que $\Delta\hat{A}$ es hermítico ($(\Delta\hat{A})\adj = \Delta\hat{A}$):
\begin{equation}
    \braket{u}{u} = \matrixel{\psi}{(\Delta\hat{A})^2}{\psi} = (\Delta A)^2.
    \label{eq:norm_u_sq}
\end{equation}
Exactamente con el mismo procedimiento para el estado $\ket{v}$:
\begin{equation}
    \braket{v}{v} = \matrixel{\psi}{(\Delta\hat{B})^2}{\psi} = (\Delta B)^2.
    \label{eq:norm_v_sq}
\end{equation}
Por otro lado, el producto interno cruzado entre ambos estados desviados resulta ser:
\begin{equation}
    \braket{u}{v} = \matrixel{\psi}{(\Delta\hat{A})\adj\Delta\hat{B}}{\psi} = \matrixel{\psi}{\Delta\hat{A}\Delta\hat{B}}{\psi} = \ev{\Delta\hat{A}\Delta\hat{B}}.
    \label{eq:inner_uv}
\end{equation}
Sustituyendo directamente los resultados de las Ecs.~\eqref{eq:norm_u_sq}, \eqref{eq:norm_v_sq} y \eqref{eq:inner_uv} en la desigualdad de Cauchy--Schwarz \eqref{eq:cauchy_schwarz}:
\begin{equation}
    \boxed{|\ev{\Delta\hat{A}\Delta\hat{B}}|^2 \le (\Delta A)^2 (\Delta B)^2}.
    \label{eq:cs_applied}
\end{equation}

\subsection{Descomposición Algebraica: Conmutador y Anticonmutador}

El producto de dos operadores hermíticos arbitrarios $\Delta\hat{A}\Delta\hat{B}$ no es en general hermítico, ya que $(\Delta\hat{A}\Delta\hat{B})\adj = \Delta\hat{B}\adj\Delta\hat{A}\adj = \Delta\hat{B}\Delta\hat{A} \neq \Delta\hat{A}\Delta\hat{B}$ a menos que conmuten.

Sin embargo, cualquier producto de operadores puede descomponerse de forma exacta y única en una \textbf{parte simétrica} y una \textbf{parte antisimétrica} mediante la identidad:
\begin{equation}
    \Delta\hat{A}\Delta\hat{B} = \frac{1}{2}\{\Delta\hat{A}, \Delta\hat{B}\} + \frac{1}{2}[\Delta\hat{A}, \Delta\hat{B}],
    \label{eq:symmetric_antisymmetric_decomp}
\end{equation}
donde $\{\hat{X}, \hat{Y}\} \equiv \hat{X}\hat{Y} + \hat{Y}\hat{X}$ es el \textbf{anticonmutador} y $[\hat{X}, \hat{Y}] \equiv \hat{X}\hat{Y} - \hat{Y}\hat{X}$ es el \textbf{conmutador}.

\subsubsection{Cálculo Detallado del Conmutador de Desviaciones}
Desarrollemos algebraicamente el conmutador $[\Delta\hat{A}, \Delta\hat{B}]$:
\begin{align}
    [\Delta\hat{A}, \Delta\hat{B}] &= [\hat{A} - \ev{\hat{A}}\hat{I}, \, \hat{B} - \ev{\hat{B}}\hat{I}] \nonumber\\
    &= [\hat{A}, \hat{B}] - [\hat{A}, \ev{\hat{B}}\hat{I}] - [\ev{\hat{A}}\hat{I}, \hat{B}] + [\ev{\hat{A}}\hat{I}, \ev{\hat{B}}\hat{I}].
\end{align}
Dado que los valores esperados $\ev{\hat{A}}$ y $\ev{\hat{B}}$ son números escalares ordinarios (c-números) y la identidad $\hat{I}$ conmuta con absolutamente cualquier operador ($[\hat{X}, c\hat{I}] = c[\hat{X}, \hat{I}] = 0$):
\begin{equation}
    \boxed{[\Delta\hat{A}, \Delta\hat{B}] = [\hat{A}, \hat{B}]}.
    \label{eq:comm_dev_eq}
\end{equation}
¡El conmutador de las fluctuaciones es exactamente idéntico al conmutador de los operadores originales! Las medias estadísticas constantes se cancelan por completo.

\subsubsection{Propiedades de Hermiticidad de las Partes}
Analicemos la naturaleza de cada operador bajo la conjugación hermítica:
\begin{itemize}
    \item \textbf{El anticonmutador $\{\Delta\hat{A}, \Delta\hat{B}\}$ es HERMÍTICO:}
    \begin{align}
        \{\Delta\hat{A}, \Delta\hat{B}\}\adj &= (\Delta\hat{A}\Delta\hat{B} + \Delta\hat{B}\Delta\hat{A})\adj = \Delta\hat{B}\adj\Delta\hat{A}\adj + \Delta\hat{A}\adj\Delta\hat{B}\adj \nonumber\\
        &= \Delta\hat{B}\Delta\hat{A} + \Delta\hat{A}\Delta\hat{B} = \{\Delta\hat{A}, \Delta\hat{B}\}.
    \end{align}
    Por consiguiente, su valor esperado es obligatoriamente un \textbf{número real puro}:
    \begin{equation}
        \ev{\{\Delta\hat{A}, \Delta\hat{B}\}} \in \mathbb{R}.
    \end{equation}
    Esta cantidad real define la \textbf{covarianza cuántica}:
    \begin{equation}
        \mathrm{Cov}(\hat{A}, \hat{B}) \equiv \frac{1}{2}\ev{\{\Delta\hat{A}, \Delta\hat{B}\}} = \frac{1}{2}\ev{\hat{A}\hat{B} + \hat{B}\hat{A}} - \ev{\hat{A}}\ev{\hat{B}}.
    \end{equation}
    
    \item \textbf{El conmutador $[\hat{A}, \hat{B}]$ es ANTIHERMÍTICO:}
    \begin{equation}
        [\hat{A}, \hat{B}]\adj = (\hat{A}\hat{B} - \hat{B}\hat{A})\adj = \hat{B}\adj\hat{A}\adj - \hat{A}\adj\hat{B}\adj = \hat{B}\hat{A} - \hat{A}\hat{B} = -[\hat{A}, \hat{B}].
    \end{equation}
    Por consiguiente, su valor esperado es estrictamente un \textbf{número imaginario puro}:
    \begin{equation}
        \ev{[\hat{A}, \hat{B}]} \in \iu\mathbb{R} \implies \ev{[\hat{A}, \hat{B}]}^* = -\ev{[\hat{A}, \hat{B}]}.
    \end{equation}
\end{itemize}

\subsection{Deducción de la Desigualdad de Robertson--\texorpdfstring{Schrödinger}{Schrodinger}}

Tomando el valor esperado de la descomposición \eqref{eq:symmetric_antisymmetric_decomp}:
\begin{equation}
    \ev{\Delta\hat{A}\Delta\hat{B}} = \underbrace{\frac{1}{2}\ev{\{\Delta\hat{A}, \Delta\hat{B}\}}}_{\text{Parte Real } X} + \underbrace{\frac{1}{2}\ev{[\hat{A}, \hat{B}]}}_{\text{Parte Imaginaria } \iu Y}.
\end{equation}
Para cualquier número complejo $z = X + \iu Y$ con $X, Y \in \mathbb{R}$, el cuadrado de su módulo se calcula de forma exacta como la suma de los cuadrados de sus partes real e imaginaria:
\begin{equation}
    |z|^2 = (\mathrm{Re}\,z)^2 + (\mathrm{Im}\,z)^2 = X^2 + Y^2.
\end{equation}
Identificando los términos:
\begin{align}
    \mathrm{Re}\left(\ev{\Delta\hat{A}\Delta\hat{B}}\right) &= \frac{1}{2}\ev{\{\Delta\hat{A}, \Delta\hat{B}\}}, \\
    \mathrm{Im}\left(\ev{\Delta\hat{A}\Delta\hat{B}}\right) &= \frac{1}{2\iu}\ev{[\hat{A}, \hat{B}]} = -\frac{\iu}{2}\ev{[\hat{A}, \hat{B}]}.
\end{align}
Por lo tanto:
\begin{equation}
    |\ev{\Delta\hat{A}\Delta\hat{B}}|^2 = \frac{1}{4}\left|\ev{\{\Delta\hat{A}, \Delta\hat{B}\}}\right|^2 + \frac{1}{4}\left|\ev{[\hat{A}, \hat{B}]}\right|^2.
    \label{eq:modulus_expanded}
\end{equation}
Sustituyendo la Ec.~\eqref{eq:modulus_expanded} en la desigualdad de Cauchy--Schwarz \eqref{eq:cs_applied}, obtenemos el teorema fundamental:

\begin{teoremabox}[Principio Generalizado de Incertidumbre de Robertson--Schrödinger (1930)]
    Para dos observables hermíticos arbitrarios $\hat{A}$ y $\hat{B}$ medidos sobre un estado cuántico $\ket{\psi}$:
    \begin{equation}
        \boxed{(\Delta A)^2 (\Delta B)^2 \ge \frac{1}{4}\left|\ev{\{\Delta\hat{A}, \Delta\hat{B}\}}\right|^2 + \frac{1}{4}\left|\ev{[\hat{A}, \hat{B}]}\right|^2}.
        \label{eq:robertson_schrodinger}
    \end{equation}
\end{teoremabox}

\subsection{Reducción a la Relación de Incertidumbre de Heisenberg}

\begin{alertabox}[¿Por qué puede descartarse el término del anticonmutador?]
    El término que contiene el anticonmutador es el cuadrado de un número real:
    \begin{equation}
        \frac{1}{4}\left|\ev{\{\Delta\hat{A}, \Delta\hat{B}\}}\right|^2 \ge 0.
    \end{equation}
    Al tratarse de una cantidad estrictamente no negativa, eliminarla del miembro derecho de la desigualdad produce una cota inferior más débil, pero igualmente válida:
    \begin{equation}
        (\Delta A)^2 (\Delta B)^2 \ge \frac{1}{4}\left|\ev{[\hat{A}, \hat{B}]}\right|^2.
    \end{equation}
    Extrayendo la raíz cuadrada a ambos lados, obtenemos la famosísima \textbf{Relación de Incertidumbre de Heisenberg--Robertson (1929)}:
    \begin{equation}
        \boxed{\Delta A \, \Delta B \ge \frac{1}{2}\left|\ev{[\hat{A}, \hat{B}]}\right|}.
        \label{eq:heisenberg_robertson}
    \end{equation}
    \textbf{Advertencia Física:} Si bien la relación de Heisenberg \eqref{eq:heisenberg_robertson} es la más conocida en los textos, la desigualdad completa de Robertson--Schrödinger \eqref{eq:robertson_schrodinger} es más general y siempre al menos tan restrictiva como la de Heisenberg (siendo estrictamente más restrictiva siempre que la covarianza cuántica $\mathrm{Cov}(\hat{A}, \hat{B}) \neq 0$). Existen estados cuánticos correlacionados (como los \textit{estados comprimidos} o \textit{squeezed states} en óptica cuántica) donde la covarianza cuántica $\mathrm{Cov}(\hat{A}, \hat{B}) \neq 0$ juega un papel crucial, haciendo que el producto $\Delta A \Delta B$ sea estrictamente mayor que $\frac{1}{2}|\ev{[\hat{A}, \hat{B}]}|$ \cite{ballentine}.
\end{alertabox}

\subsection{Ejemplo Resuelto 2: Relación de Incertidumbre Posición--Momento y Paquete Gaussiano}

\begin{ejemplobox}[Ejemplo Resuelto 2: Incertidumbre Mínima y Paquete de Ondas Coherente]
    \textbf{Enunciado:}
    \begin{enumerate}
        \item Utilice la relación de Heisenberg--Robertson \eqref{eq:heisenberg_robertson} para hallar el límite inferior del producto de incertidumbres de la posición $\hat{x}$ y el momento lineal $\hat{p}_x$, sabiendo que el conmutador canónico es $[\hat{x}, \hat{p}_x] = \iu\hbar\hat{I}$.
        \item Considere una partícula libre descrita por una función de onda gaussiana real normalizada:
        \begin{equation}
            \psi(x) = \left(\frac{1}{\pi \sigma^2}\right)^{1/4} \exp\left(-\frac{x^2}{2\sigma^2}\right).
        \end{equation}
        Calcule $\ev{\hat{x}}$, $\ev{\hat{p}_x}$, $\Delta x$, $\Delta p_x$ y evalúe la covarianza cuántica $\mathrm{Cov}(\hat{x}, \hat{p}_x) = \frac{1}{2}\ev{\{\Delta\hat{x}, \Delta\hat{p}_x\}}$.
    \end{enumerate}

    \textbf{Solución Paso a Paso:}

    \textit{Parte 1: Cota Teórica:}
    Evaluamos el valor esperado del conmutador canónico:
    \begin{equation}
        \ev{[\hat{x}, \hat{p}_x]} = \ev{\iu\hbar\hat{I}} = \iu\hbar\braket{\psi}{\psi} = \iu\hbar.
    \end{equation}
    El módulo es $|\iu\hbar| = \hbar$. Aplicando directamente la Ec.~\eqref{eq:heisenberg_robertson}:
    \begin{equation}
        \Delta x \, \Delta p_x \ge \frac{1}{2}|\iu\hbar| = \frac{\hbar}{2}.
    \end{equation}
    Esta es la cota universal de Heisenberg para coordenadas cartesianas conjugadas.

    \textit{Parte 2: Cálculo Explícito en el Paquete Gaussiano:}
    Dado que $|\psi(x)|^2$ es una función par en $x \in (-\infty, \infty)$:
    \begin{equation}
        \ev{\hat{x}} = \int_{-\infty}^\infty x |\psi(x)|^2\,\dif x = 0.
    \end{equation}
    Calculamos $\ev{\hat{x}^2}$ mediante la integral gaussiana $\int_{-\infty}^\infty x^2 \eu^{-\alpha x^2}\,\dif x = \frac{1}{2}\sqrt{\frac{\pi}{\alpha^3}}$ con $\alpha = 1/\sigma^2$:
    \begin{equation}
        \ev{\hat{x}^2} = \left(\frac{1}{\pi\sigma^2}\right)^{1/2} \int_{-\infty}^\infty x^2 \eu^{-x^2/\sigma^2}\,\dif x = \frac{1}{\sqrt{\pi}\sigma} \cdot \frac{\sqrt{\pi}}{2}\sigma^3 = \frac{\sigma^2}{2}.
    \end{equation}
    Por ende:
    \begin{equation}
        \Delta x = \sqrt{\ev{\hat{x}^2} - \ev{\hat{x}}^2} = \sqrt{\frac{\sigma^2}{2}} = \frac{\sigma}{\sqrt{2}}.
    \end{equation}
    Para el momento lineal, actuamos con el operador diferencial $\hat{p}_x = -\iu\hbar\,\frac{\dif}{\dif x}$:
    \begin{equation}
        \hat{p}_x \psi(x) = -\iu\hbar \frac{\dif}{\dif x}\left[ C \eu^{-x^2/(2\sigma^2)} \right] = -\iu\hbar \left(-\frac{x}{\sigma^2}\right)\psi(x) = \iu\hbar\frac{x}{\sigma^2}\psi(x).
    \end{equation}
    El valor esperado del momento es nulo por antisimetría:
    \begin{equation}
        \ev{\hat{p}_x} = \int_{-\infty}^\infty \psi^*(x)\hat{p}_x\psi(x)\,\dif x = \frac{\iu\hbar}{\sigma^2}\int_{-\infty}^\infty x |\psi(x)|^2\,\dif x = 0.
    \end{equation}
    Para $\ev{\hat{p}_x^2}$, usamos la propiedad de que $\hat{p}_x$ es hermítico ($\braket{\psi}{\hat{p}_x^2\psi} = \braket{\hat{p}_x\psi}{\hat{p}_x\psi}$):
    \begin{align}
        \ev{\hat{p}_x^2} &= \int_{-\infty}^\infty |\hat{p}_x\psi(x)|^2\,\dif x = \int_{-\infty}^\infty \left|\iu\hbar\frac{x}{\sigma^2}\psi(x)\right|^2\,\dif x \nonumber\\
        &= \frac{\hbar^2}{\sigma^4}\int_{-\infty}^\infty x^2 |\psi(x)|^2\,\dif x = \frac{\hbar^2}{\sigma^4}\ev{\hat{x}^2} = \frac{\hbar^2}{\sigma^4}\left(\frac{\sigma^2}{2}\right) = \frac{\hbar^2}{2\sigma^2}.
    \end{align}
    Por ende:
    \begin{equation}
        \Delta p_x = \sqrt{\ev{\hat{p}_x^2} - \ev{\hat{p}_x}^2} = \frac{\hbar}{\sqrt{2}\sigma}.
    \end{equation}
    Multiplicando las dos incertidumbres:
    \begin{equation}
        \Delta x \, \Delta p_x = \left(\frac{\sigma}{\sqrt{2}}\right)\left(\frac{\hbar}{\sqrt{2}\sigma}\right) = \frac{\hbar}{2}.
    \end{equation}
    ¡El paquete gaussiano satura con precisión absoluta la igualdad de Heisenberg!

    Evaluemos finalmente la covarianza: como $\ev{\hat{x}} = \ev{\hat{p}_x} = 0$, se tiene $\Delta\hat{x} = \hat{x}$ y $\Delta\hat{p}_x = \hat{p}_x$. Calculando el valor esperado del producto ordenado:
    \begin{align}
        \ev{\hat{x}\hat{p}_x} &= \int_{-\infty}^\infty \psi^*(x)\hat{x}\left(\iu\hbar\frac{x}{\sigma^2}\psi(x)\right)\dif x = \frac{\iu\hbar}{\sigma^2}\int_{-\infty}^\infty x^2 |\psi(x)|^2\,\dif x \nonumber\\
        &= \frac{\iu\hbar}{\sigma^2}\ev{\hat{x}^2} = \frac{\iu\hbar}{\sigma^2}\left(\frac{\sigma^2}{2}\right) = \frac{\iu\hbar}{2}.
        \label{a2-eq:xp_exp}
    \end{align}
    Puesto que $\hat{x}$ y $\hat{p}_x$ son operadores hermíticos, $(\hat{x}\hat{p}_x)\adj = \hat{p}_x\hat{x}$, de modo que el valor esperado del producto inverso es su complejo conjugado:
    \begin{equation}
        \ev{\hat{p}_x\hat{x}} = \ev{\hat{x}\hat{p}_x}^* = \left(\frac{\iu\hbar}{2}\right)^* = -\frac{\iu\hbar}{2}.
        \label{a2-eq:px_exp}
    \end{equation}
    Por consiguiente, el valor esperado del anticonmutador se anula de forma exacta:
    \begin{equation}
        \ev{\{\hat{x}, \hat{p}_x\}} = \ev{\hat{x}\hat{p}_x + \hat{p}_x\hat{x}} = \ev{\hat{x}\hat{p}_x} + \ev{\hat{p}_x\hat{x}} = \frac{\iu\hbar}{2} - \frac{\iu\hbar}{2} = 0 \implies \mathrm{Cov}(\hat{x}, \hat{p}_x) = 0.
    \end{equation}
    Esto demuestra por cálculo explícito que para el paquete gaussiano mínimo la covarianza cuántica es estrictamente nula, de modo que las cotas de Heisenberg y Robertson--Schrödinger coinciden en $\hbar/2$.
\end{ejemplobox}

\begin{figure}[H]
    \centering
    % Fig. 1: paquete gaussiano en x y en p (hbar = 1)
\begin{tikzpicture}
\begin{groupplot}[
  group style={group size=2 by 1, horizontal sep=1.8cm},
  width=0.48\linewidth, height=5.2cm, domain=-4:4, samples=200,
  axis lines=left, xlabel near ticks, ylabel near ticks,
  legend style={font=\scriptsize, draw=none, fill=none},
  every axis plot/.append style={thick}]
\nextgroupplot[xlabel={$x$ (unidades arbitrarias de longitud, $\hbar=1$)}, ylabel={$|\psi(x)|^2$},
  xlabel style={font=\scriptsize}, title={\small Espacio de posiciones}, ymin=0, ymax=1.35,
  legend pos=north east]
  \addplot[uni-blue]{exp(-x^2/0.5^2)/(sqrt(pi)*0.5)};
  \addplot[alert-border]{exp(-x^2/1.5^2)/(sqrt(pi)*1.5)};
  \legend{$\sigma=0.5$, $\sigma=1.5$}
\nextgroupplot[xlabel={$p$ (unidades arbitrarias de momento, $\hbar=1$)}, ylabel={$|\tilde\psi(p)|^2$},
  xlabel style={font=\scriptsize}, title={\small Espacio de momentos}, ymin=0, ymax=1.35, legend pos=north east]
  \addplot[uni-blue]{0.5*exp(-0.5^2*x^2)/sqrt(pi)};
  \addplot[alert-border]{1.5*exp(-1.5^2*x^2)/sqrt(pi)};
  \legend{$\sigma=0.5$, $\sigma=1.5$}
\end{groupplot}
\end{tikzpicture}

    \caption{Paquete gaussiano del Ejemplo Resuelto 2 en posici\'on y en momento ($\hbar=1$). Un $\sigma$ peque\~no localiza la part\'icula en $x$ pero ensancha $|\tilde\psi(p)|^2$, y viceversa; en ambos casos $\Delta x\,\Delta p_x=\hbar/2$.}
    \label{fig:gauss_xp}
\end{figure}

%% ============================================================
%%  SECCIÓN 4 — Funciones de Operadores, Hadamard y Teorema BCH
%% ============================================================

\section{Funciones Analíticas de Operadores, Álgebra de Lie y Lemas BCH}
\label{sec:funciones_operadores}

\subsection{Definición de Funciones de Operadores Mediante Series de Potencias}

En mecánica cuántica avanzada, es sumamente frecuente encontrar expresiones donde un operador figura como argumento de una función trascendente, tal como la exponencial de un operador de evolución temporal $\eu^{-\iu\hat{H}t/\hbar}$ o de traslación espacial $\eu^{-\iu a\hat{p}/\hbar}$ \cite{dirac, sakurai}.

\begin{teoremabox}[Definición: Función Analítica de un Operador]
    Sea $f(z)$ una función analítica de variable compleja con desarrollo en serie de Taylor alrededor del origen con radio de convergencia $R$:
    \begin{equation}
        f(z) = \sum_{n=0}^\infty c_n z^n, \quad c_n = \frac{f^{(n)}(0)}{n!}.
    \end{equation}
    Para un operador lineal continuo $\hat{A}$ actuando sobre $\mathcal{H}$ cuyo radio espectral sea menor que $R$, la función del operador $f(\hat{A})$ se define formalmente por la serie de potencias operacionales:
    \begin{equation}
        f(\hat{A}) \equiv \sum_{n=0}^\infty c_n \hat{A}^n = c_0 \hat{I} + c_1 \hat{A} + c_2 \hat{A}^2 + \dots
        \label{eq:operator_function_def}
    \end{equation}
\end{teoremabox}

En particular, la \textbf{exponencial de un operador} es una función entera (convergente en norma para cualquier operador lineal acotado\footnote{En rigor del análisis funcional, observables cuánticos fundamentales como la posición $\hat{x}$, el momento $\hat{p}$ o el Hamiltoniano $\hat{H}$ son operadores \textit{no acotados} en $\mathcal{L}^2(\mathbb{R})$, de modo que sus series de potencias no convergen en norma sobre la totalidad del espacio. Para tales operadores autoadjuntos, la exponencial unitaria $\eu^{-\iu t\hat{A}/\hbar}$ (como $\eu^{-\iu a\hat{p}/\hbar}$ o $\eu^{-\iu t\hat{H}/\hbar}$) queda rigurosamente definida sobre todo el espacio de Hilbert $\mathcal{H}$ mediante el cálculo funcional espectral y el célebre \textbf{Teorema de Stone} sobre grupos uniparamétricos fuertemente continuos de operadores unitarios \cite{ballentine}. La expansión formal en serie de Taylor es válida en dominios densos de vectores analíticos. En este cuaderno mantendremos el tratamiento operacional estándar utilizado en la física teórica \cite{sakurai, zettili}.}):
\begin{equation}
    \boxed{\eu^{\hat{A}} \equiv \sum_{n=0}^\infty \frac{1}{n!} \hat{A}^n = \hat{I} + \hat{A} + \frac{1}{2!}\hat{A}^2 + \frac{1}{3!}\hat{A}^3 + \dots}
    \label{eq:operator_exponential}
\end{equation}

\begin{alertabox}[Propiedades Críticas de la Exponencial de Operadores]
    \begin{enumerate}
        \item Si $\hat{A}$ y $\hat{B}$ \textbf{conmutan} ($[\hat{A}, \hat{B}] = 0$), entonces se comportan como escalares y se cumple la propiedad exponencial usual:
        \begin{equation}
            \eu^{\hat{A}}\eu^{\hat{B}} = \eu^{\hat{A}+\hat{B}} = \eu^{\hat{B}}\eu^{\hat{A}}.
        \end{equation}
        En particular, $\eu^{\hat{A}}\eu^{-\hat{A}} = \eu^{\hat{A}-\hat{A}} = \eu^{\hat{0}} = \hat{I}$, lo que demuestra que:
        \begin{equation}
            (\eu^{\hat{A}})^{-1} = \eu^{-\hat{A}}.
        \end{equation}
        \item Si $\hat{A}$ y $\hat{B}$ \textbf{NO conmutan} ($[\hat{A}, \hat{B}] \neq 0$), entonces:
        \begin{equation}
            \boxed{\eu^{\hat{A}}\eu^{\hat{B}} \neq \eu^{\hat{A}+\hat{B}}}.
        \end{equation}
        El orden de los factores altera el resultado. Para calcular el producto en este caso general, es indispensable recurrir al Álgebra de Lie y al Teorema de Baker--Campbell--Hausdorff.
        \item Si $\hat{A}$ es hermítico ($\hat{A}\adj = \hat{A}$), el operador $\hat{U} = \eu^{\iu\hat{A}}$ es \textbf{unitario}:
        \begin{equation}
            \hat{U}\adj = \left(\sum_{n=0}^\infty \frac{\iu^n}{n!}\hat{A}^n\right)\adj = \sum_{n=0}^\infty \frac{(-\iu)^n}{n!}(\hat{A}\adj)^n = \sum_{n=0}^\infty \frac{(-\iu)^n}{n!}\hat{A}^n = \eu^{-\iu\hat{A}} = \hat{U}^{-1}.
        \end{equation}
    \end{enumerate}
\end{alertabox}

\subsection{El Adjunto de un Operador Función y la Conmutación Cuántica}
\label{subsec:adjunto_operador_funcion}

En el manejo algebraico de funciones de operadores, surgen dos interrogantes cruciales para la formulación cuántica: ¿cómo se calcula el conjugado hermítico de $F(\hat{A})$ y qué sucede cuando un operador $\hat{B}$ conmuta o no con $\hat{A}$?

\begin{teoremabox}[Propiedades del Adjunto de un Operador Función]
    Sea $F(z) = \sum_{n=0}^\infty c_n z^n$ una función analítica de variable compleja con coeficientes $c_n \in \mathbb{C}$, y sea $\hat{A}$ un operador lineal continuo. Entonces:
    \begin{enumerate}
        \item \textbf{Adjunto General:} El adjunto hermítico del operador función $F(\hat{A})$ está dado por:
        \begin{equation}
            \boxed{\left[ F(\hat{A}) \right]\adj = F^*(\hat{A}\adj) \equiv \sum_{n=0}^\infty c_n^* (\hat{A}\adj)^n}.
            \label{eq:adjunto_funcion_general}
        \end{equation}
        donde $F^*(z) \equiv \sum_{n=0}^\infty c_n^* z^n$ es la función con los coeficientes complejos conjugados.
        \item \textbf{Caso de Operador Hermítico:} Si $\hat{A}$ es autoadjunto ($\hat{A}\adj = \hat{A}$):
        \begin{equation}
            \left[ F(\hat{A}) \right]\adj = F^*(\hat{A}).
        \end{equation}
        En particular, si además los coeficientes $c_n$ son números reales ($F^*(z) = F(z)$), entonces $F(\hat{A})$ es un operador estrictamente \textbf{hermítico}:
        \begin{equation}
            \left[ F(\hat{A}) \right]\adj = F(\hat{A}).
        \end{equation}
        \item \textbf{Funciones Exponenciales Complejas:} Para cualquier parámetro escalar $\alpha \in \mathbb{C}$:
        \begin{equation}
            \boxed{\left( \eu^{\iu\alpha\hat{A}} \right)\adj = \eu^{-\iu\alpha^*\hat{A}\adj}}.
            \label{eq:adjunto_exponencial}
        \end{equation}
        Si $\alpha \in \mathbb{R}$ y $\hat{A}$ es hermítico ($\hat{A}\adj = \hat{A}$), entonces:
        \begin{equation}
            \left( \eu^{\iu\alpha\hat{A}} \right)\adj = \eu^{-\iu\alpha\hat{A}} = \left( \eu^{\iu\alpha\hat{A}} \right)^{-1},
        \end{equation}
        lo que demuestra rigurosamente que $\hat{U} \equiv \eu^{\iu\alpha\hat{A}}$ es un \textbf{operador unitario}.
    \end{enumerate}
\end{teoremabox}

\subsubsection{No Conmutatividad con Funciones de Operadores}

Si dos observables conmutan ($[\hat{A}, \hat{B}] = 0$), entonces $\hat{B}$ conmuta con cualquier potencia entera $\hat{A}^n$, y por ende conmuta con cualquier función analítica:
\begin{equation}
    [\hat{A}, \hat{B}] = 0 \implies [\hat{B}, F(\hat{A})] = 0.
\end{equation}
Por el contrario, si los operadores \textbf{no conmutan} ($[\hat{A}, \hat{B}] \neq 0$), entonces en general $\hat{B}$ \textbf{no conmuta con la función del operador}:
\begin{equation}
    [\hat{B}, F(\hat{A})] \neq 0 \quad (\text{en general}).
\end{equation}
(Cabe advertir que la implicación estricta $[\hat{A}, \hat{B}] \neq 0 \implies [\hat{B}, F(\hat{A})] \neq 0$ no es universal: por ejemplo, tomando las matrices de Pauli con $\hat{A} = \sigma_x$, $\hat{B} = \sigma_z$ y la función $F(z) = z^2$, se tiene $[\sigma_x, \sigma_z] = -2\iu\sigma_y \neq 0$, pero $F(\sigma_x) = \sigma_x^2 = \hat{I}$ conmuta idénticamente con $\sigma_z$).
Para cuantificar este conmutador de forma analítica, observemos la acción de $\hat{B}$ sobre una potencia monomial $\hat{A}^n$:
\begin{equation}
    [\hat{B}, \hat{A}^n] = \sum_{k=0}^{n-1} \hat{A}^k [\hat{B}, \hat{A}] \hat{A}^{n-1-k}.
\end{equation}
En el caso físicamente fundamental en que el conmutador básico conmuta con $\hat{A}$ (es decir, $[\hat{A}, [\hat{A}, \hat{B}]] = 0$, como ocurre para el par canónico de posición y momento $[\hat{x}, \hat{p}_x] = \iu\hbar\hat{I}$), el término $[\hat{B}, \hat{A}]$ sale como constante de la sumatoria:
\begin{equation}
    [\hat{B}, \hat{A}^n] = [\hat{B}, \hat{A}] \sum_{k=0}^{n-1} \hat{A}^{n-1} = n\hat{A}^{n-1}[\hat{B}, \hat{A}] = -n\hat{A}^{n-1}[\hat{A}, \hat{B}].
\end{equation}
Multiplicando por los coeficientes de Taylor $c_n$ y sumando sobre $n$:
\begin{equation}
    [\hat{B}, F(\hat{A})] = \sum_{n=1}^\infty c_n [\hat{B}, \hat{A}^n] = -[\hat{A}, \hat{B}] \sum_{n=1}^\infty n c_n \hat{A}^{n-1} = \boxed{-[\hat{A}, \hat{B}]\,F'(\hat{A})},
    \label{eq:comm_b_f_a}
\end{equation}
donde $F'(z) = \frac{\dif F}{\dif z}$ es la derivada compleja ordinaria de la función. Esta elegante identidad operativa permite calcular de un solo vistazo conmutadores que de otra manera requerirían cálculos matriciales interminables.

\subsection{El Lema de Hadamard (Fórmula de Trenzado de Operadores)}

Consideremos la transformación de un operador $\hat{B}$ mediante una conjugación unitaria o transformación de semejanza exponencial generada por $\hat{A}$:
\begin{equation}
    \hat{F}(\lambda) \equiv \eu^{\lambda\hat{A}}\hat{B}\eu^{-\lambda\hat{A}},
    \label{eq:hadamard_ansatz}
\end{equation}
donde $\lambda$ es un parámetro escalar continuo (real o complejo, independiente de las coordenadas espaciales).

\begin{teoremabox}[Lema de Hadamard]
    Para dos operadores lineales cualesquiera $\hat{A}$ y $\hat{B}$:
    \begin{align}
        \eu^{\lambda\hat{A}}\hat{B}\eu^{-\lambda\hat{A}} &= \hat{B} + \lambda[\hat{A}, \hat{B}] + \frac{\lambda^2}{2!}[\hat{A}, [\hat{A}, \hat{B}]] + \frac{\lambda^3}{3!}[\hat{A}, [\hat{A}, [\hat{A}, \hat{B}]]] + \dots \nonumber\\
        &= \sum_{n=0}^\infty \frac{\lambda^n}{n!} \mathrm{ad}_{\hat{A}}^n(\hat{B}),
        \label{eq:hadamard_lemma}
    \end{align}
    donde el operador de la representación adjunta $\mathrm{ad}_{\hat{A}}$ se define recursivamente por:
    \begin{equation}
        \mathrm{ad}_{\hat{A}}^0(\hat{B}) = \hat{B}, \quad \mathrm{ad}_{\hat{A}}^1(\hat{B}) = [\hat{A}, \hat{B}], \quad \mathrm{ad}_{\hat{A}}^n(\hat{B}) = [\hat{A}, \mathrm{ad}_{\hat{A}}^{n-1}(\hat{B})].
    \end{equation}
    En particular, evaluando en el valor unitario del parámetro $\lambda=1$:
    \begin{equation}
        \boxed{\eu^{\hat{A}}\hat{B}\eu^{-\hat{A}} = \hat{B} + [\hat{A}, \hat{B}] + \frac{1}{2!}[\hat{A}, [\hat{A}, \hat{B}]] + \frac{1}{3!}[\hat{A}, [\hat{A}, [\hat{A}, \hat{B}]]] + \dots}
    \end{equation}
\end{teoremabox}

\textbf{Demostración Rigurosa (Método de la Función Generatriz Paramétrica):}\\
Puesto que $\hat{F}(\lambda)$ depende analíticamente del escalar $\lambda$, podemos expandirla en una serie de Maclaurin ordinaria alrededor de $\lambda = 0$:
\begin{equation}
    \hat{F}(\lambda) = \hat{F}(0) + \lambda\left.\frac{\dif\hat{F}}{\dif \lambda}\right|_{\lambda=0} + \frac{\lambda^2}{2!}\left.\frac{\dif^2\hat{F}}{\dif \lambda^2}\right|_{\lambda=0} + \dots = \sum_{n=0}^\infty \frac{\lambda^n}{n!} \left.\frac{\dif^n\hat{F}}{\dif \lambda^n}\right|_{\lambda=0}.
    \label{eq:hadamard_maclaurin}
\end{equation}
Evaluemos el término de orden cero en $\lambda = 0$:
\begin{equation}
    \hat{F}(0) = \eu^{0\hat{A}}\hat{B}\eu^{-0\hat{A}} = \hat{I}\hat{B}\hat{I} = \hat{B}.
\end{equation}
Calculemos ahora la primera derivada de $\hat{F}(\lambda)$ respecto a $\lambda$ empleando la regla del producto de Leibniz para operadores:
\begin{align}
    \frac{\dif\hat{F}}{\dif \lambda} &= \frac{\dif}{\dif \lambda}\left( \eu^{\lambda\hat{A}}\hat{B}\eu^{-\lambda\hat{A}} \right) \nonumber\\
    &= \left(\frac{\dif}{\dif \lambda}\eu^{\lambda\hat{A}}\right)\hat{B}\eu^{-\lambda\hat{A}} + \eu^{\lambda\hat{A}}\hat{B}\left(\frac{\dif}{\dif \lambda}\eu^{-\lambda\hat{A}}\right).
\end{align}
Derivando término a término la serie de la exponencial escalar-operador:
\begin{equation}
    \frac{\dif}{\dif \lambda}\eu^{\lambda\hat{A}} = \frac{\dif}{\dif \lambda}\sum_{n=0}^\infty \frac{\lambda^n}{n!}\hat{A}^n = \sum_{n=1}^\infty \frac{n \lambda^{n-1}}{n!}\hat{A}^n = \hat{A}\sum_{m=0}^\infty \frac{\lambda^m}{m!}\hat{A}^m = \hat{A}\eu^{\lambda\hat{A}} = \eu^{\lambda\hat{A}}\hat{A}.
\end{equation}
Análogamente, para la exponencial negativa:
\begin{equation}
    \frac{\dif}{\dif \lambda}\eu^{-\lambda\hat{A}} = -\hat{A}\eu^{-\lambda\hat{A}} = -\eu^{-\lambda\hat{A}}\hat{A}.
\end{equation}
Sustituyendo estas derivadas:
\begin{align}
    \frac{\dif\hat{F}}{\dif \lambda} &= \hat{A}\left(\eu^{\lambda\hat{A}}\hat{B}\eu^{-\lambda\hat{A}}\right) - \left(\eu^{\lambda\hat{A}}\hat{B}\eu^{-\lambda\hat{A}}\right)\hat{A} \nonumber\\
    &= \hat{A}\hat{F}(\lambda) - \hat{F}(\lambda)\hat{A} = [\hat{A}, \hat{F}(\lambda)].
    \label{eq:hadamard_first_diff}
\end{align}
Evaluando la primera derivada en $\lambda = 0$:
\begin{equation}
    \left.\frac{\dif\hat{F}}{\dif \lambda}\right|_{\lambda=0} = [\hat{A}, \hat{F}(0)] = [\hat{A}, \hat{B}].
\end{equation}
Derivemos nuevamente respecto a $\lambda$ para obtener el término de segundo orden:
\begin{equation}
    \frac{\dif^2\hat{F}}{\dif \lambda^2} = \frac{\dif}{\dif \lambda}\left( \frac{\dif\hat{F}}{\dif \lambda} \right) = \frac{\dif}{\dif \lambda}[\hat{A}, \hat{F}(\lambda)] = \left[\hat{A}, \frac{\dif\hat{F}}{\dif \lambda}\right].
\end{equation}
Evaluando en $\lambda = 0$:
\begin{equation}
    \left.\frac{\dif^2\hat{F}}{\dif \lambda^2}\right|_{\lambda=0} = \left[\hat{A}, \left.\frac{\dif\hat{F}}{\dif \lambda}\right|_{\lambda=0}\right] = [\hat{A}, [\hat{A}, \hat{B}]].
\end{equation}
Por inducción matemática inmediata, la $n$-ésima derivada evaluada en el origen es:
\begin{equation}
    \left.\frac{\dif^n\hat{F}}{\dif \lambda^n}\right|_{\lambda=0} = \mathrm{ad}_{\hat{A}}^n(\hat{B}) = \underbrace{[\hat{A}, [\hat{A}, \dots, [\hat{A}, \hat{B}]] \dots ]}_{n \text{ conmutadores anidados}}.
\end{equation}
Sustituyendo estos coeficientes en la serie de Taylor \eqref{eq:hadamard_maclaurin}, queda demostrado formalmente el Lema de Hadamard \cite{zettili}.

\subsection{El Teorema de Baker--Campbell--Hausdorff (BCH) Simplificado}

¿Cómo podemos calcular el producto de dos exponenciales $\eu^{\hat{A}}\eu^{\hat{B}}$ cuando los operadores no conmutan? El teorema general de Baker, Campbell y Hausdorff establece que existe un operador $\hat{C}$ tal que $\eu^{\hat{A}}\eu^{\hat{B}} = \eu^{\hat{C}}$, donde $\hat{C} = \hat{A} + \hat{B} + \frac{1}{2}[\hat{A}, \hat{B}] + \frac{1}{12}[\hat{A}, [\hat{A}, \hat{B}]] - \frac{1}{12}[\hat{B}, [\hat{A}, \hat{B}]] + \dots$

En física cuántica, existe un caso extraordinariamente frecuente y de vital importancia (llamado el \textbf{caso c-número}): aquel donde el conmutador $[\hat{A}, \hat{B}]$ conmuta tanto con $\hat{A}$ como con $\hat{B}$, es decir:
\begin{equation}
    [[\hat{A}, \hat{B}], \hat{A}] = 0 \quad \text{y} \quad [[\hat{A}, \hat{B}], \hat{B}] = 0.
    \label{eq:bch_hypothesis}
\end{equation}
Esto ocurre siempre que $[\hat{A}, \hat{B}] = c\hat{I}$ es un escalar (por ejemplo, el conmutador canónico de posición y momento $[\hat{x}, \hat{p}_x] = \iu\hbar\hat{I}$, o los operadores de creación y aniquilación $[\hat{a}, \hat{a}\adj] = \hat{I}$).

\begin{teoremabox}[Teorema de Baker--Campbell--Hausdorff Simplificado (Fórmula de Weyl)]
    Si dos operadores lineales $\hat{A}$ y $\hat{B}$ satisfacen la condición \eqref{eq:bch_hypothesis}, entonces:
    \begin{equation}
        \boxed{\eu^{\hat{A}}\eu^{\hat{B}} = \eu^{\hat{A}+\hat{B}}\eu^{\frac{1}{2}[\hat{A}, \hat{B}]}}.
        \label{eq:bch_simplified}
    \end{equation}
\end{teoremabox}

\textbf{Demostración Paso a Paso:}\\
Para demostrar esta relación, introducimos un operador dependiente de un parámetro continuo $t \in [0, 1]$:
\begin{equation}
    \hat{G}(t) \equiv \eu^{t\hat{A}}\eu^{t\hat{B}}.
    \label{eq:g_ansatz}
\end{equation}
Notemos que para $t = 0$: $\hat{G}(0) = \eu^{\hat{0}}\eu^{\hat{0}} = \hat{I}$, y para $t = 1$: $\hat{G}(1) = \eu^{\hat{A}}\eu^{\hat{B}}$ (que es la magnitud deseada).

Derivemos $\hat{G}(t)$ respecto al parámetro escalar $t$:
\begin{equation}
    \frac{\dif\hat{G}}{\dif t} = \left(\frac{\dif}{\dif t}\eu^{t\hat{A}}\right)\eu^{t\hat{B}} + \eu^{t\hat{A}}\left(\frac{\dif}{\dif t}\eu^{t\hat{B}}\right) = \hat{A}\eu^{t\hat{A}}\eu^{t\hat{B}} + \eu^{t\hat{A}}\hat{B}\eu^{t\hat{B}}.
\end{equation}

\begin{alertabox}[¿Por qué insertamos la identidad en la forma $\eu^{-t\hat{A}}\eu^{t\hat{A}}$?]
    En el primer término, $\eu^{t\hat{A}}\eu^{t\hat{B}} = \hat{G}(t)$, de modo que queda $\hat{A}\hat{G}(t)$.\\
    Sin embargo, en el segundo término, el operador $\hat{B}$ se encuentra \enquote{atrapado} en medio de las exponenciales: $\eu^{t\hat{A}}\hat{B}\eu^{t\hat{B}}$.\\
    \textbf{¡NO podemos simplemente sacar $\hat{B}$ hacia la izquierda ni hacia la derecha porque $\hat{B}$ no conmuta con $\eu^{t\hat{A}}$!}
    
    Para poder factorizar $\hat{G}(t) = \eu^{t\hat{A}}\eu^{t\hat{B}}$ por la derecha de toda la ecuación, insertamos estratégicamente el operador identidad $\hat{I} = \eu^{-t\hat{A}}\eu^{t\hat{A}}$ justo entre $\hat{B}$ y $\eu^{t\hat{B}}$:
    \begin{equation}
        \eu^{t\hat{A}}\hat{B}\eu^{t\hat{B}} = \eu^{t\hat{A}}\hat{B}\underbrace{(\eu^{-t\hat{A}}\eu^{t\hat{A}})}_{\hat{I}}\eu^{t\hat{B}} = \left(\eu^{t\hat{A}}\hat{B}\eu^{-t\hat{A}}\right)\left(\eu^{t\hat{A}}\eu^{t\hat{B}}\right) = \left(\eu^{t\hat{A}}\hat{B}\eu^{-t\hat{A}}\right)\hat{G}(t).
    \end{equation}
    ¡Observemos la genialidad de este paso algebraico! La expresión entre paréntesis es exactamente la estructura del \textbf{Lema de Hadamard} con parámetro $t$.
\end{alertabox}

Aplicando el Lema de Hadamard \eqref{eq:hadamard_lemma} a la cantidad $\eu^{t\hat{A}}\hat{B}\eu^{-t\hat{A}}$:
\begin{equation}
    \eu^{t\hat{A}}\hat{B}\eu^{-t\hat{A}} = \hat{B} + t[\hat{A}, \hat{B}] + \frac{t^2}{2!}[\hat{A}, [\hat{A}, \hat{B}]] + \dots
\end{equation}
Pero según nuestra hipótesis fundamental de trabajo \eqref{eq:bch_hypothesis}, el segundo conmutador y todos los superiores se anulan idénticamente: $[\hat{A}, [\hat{A}, \hat{B}]] = 0$. Por lo tanto, la serie infinita \textbf{se trunca exactamente en el segundo término}:
\begin{equation}
    \eu^{t\hat{A}}\hat{B}\eu^{-t\hat{A}} = \hat{B} + t[\hat{A}, \hat{B}].
\end{equation}
Sustituyendo esto en la derivada de $\hat{G}(t)$:
\begin{align}
    \frac{\dif\hat{G}}{\dif t} &= \hat{A}\hat{G}(t) + (\hat{B} + t[\hat{A}, \hat{B}])\hat{G}(t) \nonumber\\
    &= \left( \hat{A} + \hat{B} + t[\hat{A}, \hat{B}] \right)\hat{G}(t).
    \label{eq:diff_eq_g}
\end{align}
Denotemos al operador entre paréntesis por $\hat{M}(t) \equiv \hat{A} + \hat{B} + t[\hat{A}, \hat{B}]$. Como $[\hat{A}+\hat{B}, [\hat{A}, \hat{B}]] = 0$, el operador $\hat{M}(t)$ conmuta consigo mismo a diferentes instantes $t$ y $t'$:
\begin{equation}
    [\hat{M}(t), \hat{M}(t')] = [\hat{A}+\hat{B} + t[\hat{A},\hat{B}], \, \hat{A}+\hat{B} + t'[\hat{A},\hat{B}]] = 0.
\end{equation}
Debido a esta conmutación, la ecuación diferencial de operadores \eqref{eq:diff_eq_g} se integra exactamente como una ecuación diferencial escalar por separación de variables:
\begin{equation}
    \hat{G}(t) = \exp\left( \int_0^t \hat{M}(t')\,\dif t' \right)\hat{G}(0).
\end{equation}
Evaluemos la integral en el argumento de la exponencial:
\begin{equation}
    \int_0^t \left( \hat{A} + \hat{B} + t'[\hat{A}, \hat{B}] \right)\,\dif t' = t(\hat{A} + \hat{B}) + \frac{t^2}{2}[\hat{A}, \hat{B}].
\end{equation}
Puesto que $\hat{G}(0) = \hat{I}$:
\begin{equation}
    \hat{G}(t) = \exp\left( t(\hat{A} + \hat{B}) + \frac{1}{2}t^2[\hat{A}, \hat{B}] \right).
\end{equation}
Nuevamente, como el operador $[\hat{A}, \hat{B}]$ conmuta con $(\hat{A}+\hat{B})$, podemos separar la exponencial de la suma en el producto de exponenciales:
\begin{equation}
    \hat{G}(t) = \eu^{t(\hat{A}+\hat{B})}\eu^{\frac{1}{2}t^2[\hat{A}, \hat{B}]}.
\end{equation}
Evaluando finalmente en $t = 1$:
\begin{equation}
    \hat{G}(1) = \eu^{\hat{A}}\eu^{\hat{B}} = \eu^{\hat{A}+\hat{B}}\eu^{\frac{1}{2}[\hat{A}, \hat{B}]}.
\end{equation}
La demostración queda enteramente completada y justificada.

\subsection{BCH Simplificado (c-Número) frente a la Serie BCH General de Momento Angular}
\label{subsec:bch_simplificado_vs_general}

Un error pedagógico extraordinariamente común en los primeros cursos de física cuántica es asumir que la fórmula:
\begin{equation}
    \eu^{\hat{A}}\eu^{\hat{B}} = \eu^{\hat{A}+\hat{B} + \frac{1}{2}[\hat{A}, \hat{B}]}
    \label{eq:bch_simplificada_recap}
\end{equation}
es válida universalmente para cualquier par de operadores cuánticos. \textbf{Esto es falso}: la Ec.~\eqref{eq:bch_simplificada_recap} es un caso particular que depende de una condición sumamente restrictiva.

\begin{alertabox}[¿Cuándo es válida la fórmula simplificada de Baker--Campbell--Hausdorff?]
    La fórmula cerrada \eqref{eq:bch_simplificada_recap} exige estrictamente que el conmutador $[\hat{A}, \hat{B}]$ sea un \textbf{c-número} (un múltiplo escalar de la identidad $\alpha\hat{I}$), o al menos que conmute con ambos operadores:
    \begin{equation}
        [\hat{A}, [\hat{A}, \hat{B}]] = 0 \quad \text{y} \quad [\hat{B}, [\hat{A}, \hat{B}]] = 0.
    \end{equation}
    \begin{itemize}
        \item \textbf{El caso canónico de posición y momento $(\hat{x}, \hat{p}_x)$:}
        Aquí $[\hat{x}, \hat{p}_x] = \iu\hbar\hat{I}$. Dado que $\iu\hbar\hat{I}$ es proporcional al operador identidad, conmuta absolutamente con cualquier operador del universo:
        \begin{equation}
            [\hat{x}, \, \iu\hbar\hat{I}] = 0, \qquad [\hat{p}_x, \, \iu\hbar\hat{I}] = 0.
        \end{equation}
        Por lo tanto, todos los conmutadores anidados de segundo orden en adelante se anulan idénticamente, y la fórmula simplificada \eqref{eq:bch_simplificada_recap} es \textbf{rigurosamente exacta}.
        \item \textbf{El caso de momento angular $(\hat{L}_x, \hat{L}_y)$:}
        Consideremos ahora dos generadores de rotaciones espaciales. Su conmutador básico es el álgebra de Lie $\mathfrak{so}(3)$:
        \begin{equation}
            [\hat{L}_x, \hat{L}_y] = \iu\hbar\hat{L}_z.
        \end{equation}
        ¡El resultado no es un número; \textbf{es otro operador cuántico hermítico ($\hat{L}_z$)}!\\
        Evaluemos el siguiente conmutador anidado de la serie:
        \begin{equation}
            [\hat{L}_x, [\hat{L}_x, \hat{L}_y]] = [\hat{L}_x, \, \iu\hbar\hat{L}_z] = \iu\hbar[\hat{L}_x, \hat{L}_z] = \iu\hbar(-\iu\hbar\hat{L}_y) = +\hbar^2\hat{L}_y \neq 0.
        \end{equation}
        Análogamente, para el segundo operador:
        \begin{equation}
            [\hat{L}_y, [\hat{L}_x, \hat{L}_y]] = [\hat{L}_y, \, \iu\hbar\hat{L}_z] = \iu\hbar[\hat{L}_y, \hat{L}_z] = \iu\hbar(\iu\hbar\hat{L}_x) = -\hbar^2\hat{L}_x \neq 0.
        \end{equation}
        ¡Los conmutadores anidados nunca se anulan; se reproducen infinitamente rotando entre $\hat{L}_x$, $\hat{L}_y$ y $\hat{L}_z$!\\
        Por consiguiente, \textbf{la fórmula exponencial simplificada falla por completo para el momento angular}:
        \begin{equation}
            \eu^{-\frac{\iu}{\hbar}\alpha\hat{L}_x}\eu^{-\frac{\iu}{\hbar}\beta\hat{L}_y} \neq \eu^{-\frac{\iu}{\hbar}(\alpha\hat{L}_x + \beta\hat{L}_y) + \frac{1}{2}\left[-\frac{\iu}{\hbar}\alpha\hat{L}_x, -\frac{\iu}{\hbar}\beta\hat{L}_y\right]}.
            \label{eq:bch_falla_L}
        \end{equation}
    \end{itemize}
\end{alertabox}

\begin{teoremabox}[La Serie General de Baker--Campbell--Hausdorff (BCH)]
    Para dos operadores genéricos $\hat{A}$ y $\hat{B}$ que no satisfacen la hipótesis de truncamiento a primer orden, el logaritmo del producto exponencial viene dado por la serie infinita completa de conmutadores anidados:
    \begin{equation}
        \ln\left( \eu^{\hat{A}}\eu^{\hat{B}} \right) = \hat{A} + \hat{B} + \frac{1}{2}[\hat{A}, \hat{B}] + \frac{1}{12}[\hat{A}, [\hat{A}, \hat{B}]] - \frac{1}{12}[\hat{B}, [\hat{A}, \hat{B}]] - \frac{1}{24}[\hat{B}, [\hat{A}, [\hat{A}, \hat{B}]]] + \dots
        \label{eq:bch_general_series}
    \end{equation}
    En la teoría de grupos de Lie, esta serie describe la composición de dos rotaciones espaciales en el grupo $\mathrm{SO}(3)$: una rotación de ángulo $\alpha$ alrededor de $x$ seguida de una rotación $\beta$ alrededor de $y$ no es una simple suma vectorial de rotaciones, sino una nueva rotación alrededor de un eje oblicuo resultante de la no conmutatividad geométrica de las rotaciones en $\mathbb{R}^3$.
\end{teoremabox}

\begin{figure}[H]
\centering
\begin{tikzpicture}[
    x={(-0.50cm,-0.36cm)}, y={(1cm,0cm)}, z={(0cm,1cm)},
    scale=1.6, >=Stealth, font=\small]

  % ---------- Panel (a): primero R_x, luego R_y ----------
  \begin{scope}
    \draw[grisUNI!40] (0,0,0)--(-1.2,0,0);
    \draw[grisUNI!40] (0,0,0)--(0,-1.6,0);
    \draw[grisUNI!40] (0,0,0)--(0,0,-1.5);
    \draw[->,grisUNI] (0,0,0)--(1.8,0,0) node[below left]{$x$};
    \draw[->,grisUNI] (0,0,0)--(0,1.8,0) node[right]{$y$};
    \draw[->,grisUNI] (0,0,0)--(0,0,1.8) node[above]{$z$};
    \draw[->,thick,grisUNI!70!black] (0,0,0)--(1,0,1)
          node[above left]{$\vec v=(1,0,1)$};
    \draw[->,thick,azulUNI,dashed] (0,0,0)--(1,-1,0)
          node[left]{$(1,-1,0)$};
    \draw[->,very thick,rojoUNI] (0,0,0)--(0,-1,-1)
          node[below]{$(0,-1,-1)$};
    \node at (0,0.4,2.3)
          {\textbf{(a)} $\hat R_y(\tfrac{\pi}{2})\,\hat R_x(\tfrac{\pi}{2})\,\vec v$};
  \end{scope}

  % ---------- Panel (b): primero R_y, luego R_x ----------
  \begin{scope}[shift={(canvas cs:x=0cm,y=0cm)}]
    \pgftransformxshift{4.9cm}
    \draw[grisUNI!40] (0,0,0)--(-1.2,0,0);
    \draw[grisUNI!40] (0,0,0)--(0,-1.6,0);
    \draw[grisUNI!40] (0,0,0)--(0,0,-1.5);
    \draw[->,grisUNI] (0,0,0)--(1.8,0,0) node[below left]{$x$};
    \draw[->,grisUNI] (0,0,0)--(0,1.8,0) node[right]{$y$};
    \draw[->,grisUNI] (0,0,0)--(0,0,1.8) node[above]{$z$};
    \draw[->,thick,grisUNI!70!black] (0,0,0)--(1,0,1)
          node[above left]{$\vec v=(1,0,1)$};
    \draw[->,thick,azulUNI,dashed] (0,0,0)--(1,0,-1)
          node[below left]{$(1,0,-1)$};
    \draw[->,very thick,rojoUNI] (0,0,0)--(1,1,0)
          node[below right]{$(1,1,0)$};
    \node at (0,0.4,2.3)
          {\textbf{(b)} $\hat R_x(\tfrac{\pi}{2})\,\hat R_y(\tfrac{\pi}{2})\,\vec v$};
  \end{scope}
\end{tikzpicture}
\caption{Las rotaciones finitas no conmutan. Un mismo vector $\vec v=(1,0,1)$,
girado $\pi/2$ alrededor de $x$ y de $y$ en órdenes distintos, termina en
posiciones distintas (rojo). Es la versión geométrica de
$\eu^{-\iu\alpha\hat L_x/\hbar}\eu^{-\iu\beta\hat L_y/\hbar}\neq
\eu^{-\iu\beta\hat L_y/\hbar}\eu^{-\iu\alpha\hat L_x/\hbar}$ (ec.~\eqref{eq:bch_falla_L}); línea
discontinua: posición intermedia.}
\label{fig:rot-no-conmutan}
\end{figure}


\subsection{Ejemplo Resuelto 3: El Operador de Traslación Espacial y Serie de Taylor}

\begin{ejemplobox}[Ejemplo Resuelto 3: Generador Infinitesimal de Traslaciones]
    \textbf{Enunciado:} El operador de traslación espacial finita por una distancia $a \in \mathbb{R}$ se define como:
    \begin{equation}
        \hat{T}(a) \equiv \exp\left(-\frac{\iu}{\hbar}a\hat{p}_x\right).
    \end{equation}
    \begin{enumerate}
        \item Utilice el Lema de Hadamard para demostrar que el operador posición transformado es $\hat{T}\adj(a)\hat{x}\hat{T}(a) = \hat{x} + a\hat{I}$.
        \item Aplique la definición en serie de potencias de $\hat{T}(a)$ sobre una función de onda arbitraria diferenciable $\psi(x)$ en la representación de coordenadas y demuestre que genera una traslación rígida: $\hat{T}(a)\psi(x) = \psi(x - a)$.
    \end{enumerate}

    \textbf{Solución Paso a Paso:}

    \textit{Paso 1: Aplicación del Lema de Hadamard:}
    Como $\hat{p}_x$ es hermítico, $\hat{T}\adj(a) = \exp\left(\frac{\iu}{\hbar}a\hat{p}_x\right)$. Queremos calcular:
    \begin{equation}
        \hat{x}' \equiv \hat{T}\adj(a)\hat{x}\hat{T}(a) = \exp\left(\frac{\iu}{\hbar}a\hat{p}_x\right)\hat{x}\exp\left(-\frac{\iu}{\hbar}a\hat{p}_x\right).
    \end{equation}
    Identificamos en la fórmula de Hadamard \eqref{eq:hadamard_lemma}: el operador conjugante es $\hat{A} = \frac{\iu}{\hbar}a\hat{p}_x$, el operador a transformar es $\hat{B} = \hat{x}$, y el parámetro auxiliar es $\lambda = 1$.
    Calculamos el primer conmutador:
    \begin{equation}
        [\hat{A}, \hat{B}] = \left[\frac{\iu}{\hbar}a\hat{p}_x, \, \hat{x}\right] = \frac{\iu a}{\hbar}[\hat{p}_x, \hat{x}].
    \end{equation}
    Recordando que el conmutador fundamental es $[\hat{x}, \hat{p}_x] = \iu\hbar\hat{I}$, su opuesto es $[\hat{p}_x, \hat{x}] = -\iu\hbar\hat{I}$. Por lo tanto:
    \begin{equation}
        [\hat{A}, \hat{B}] = \frac{\iu a}{\hbar}(-\iu\hbar\hat{I}) = -\iu^2 a \hat{I} = a\hat{I}.
    \end{equation}
    Calculamos el segundo conmutador de la serie de Hadamard:
    \begin{equation}
        [\hat{A}, [\hat{A}, \hat{B}]] = \left[\frac{\iu}{\hbar}a\hat{p}_x, \, a\hat{I}\right] = 0,
    \end{equation}
    puesto que $a\hat{I}$ es un múltiplo escalar de la identidad y conmuta con cualquier operador. Por consiguiente, \textbf{todos los conmutadores de orden dos y superior se anulan}. Sustituyendo en Hadamard:
    \begin{equation}
        \hat{T}\adj(a)\hat{x}\hat{T}(a) = \hat{x} + [\hat{A}, \hat{B}] + 0 = \hat{x} + a\hat{I}.
    \end{equation}
    Esto demuestra de manera puramente operacional que el operador momento $\hat{p}_x$ es el generador infinitesimal de traslaciones espaciales en la posición.

    \textit{Paso 2: Acción sobre la Función de Onda:}
    En la representación de posiciones, el operador momento actúa como la derivada espacial $\hat{p}_x \to -\iu\hbar\frac{\dif}{\dif x}$. Sustituyendo esto en el operador exponencial:
    \begin{equation}
        \hat{T}(a) = \exp\left(-\frac{\iu a}{\hbar}\left(-\iu\hbar\frac{\dif}{\dif x}\right)\right) = \exp\left(\iu^2 a\frac{\dif}{\dif x}\right) = \exp\left(-a\frac{\dif}{\dif x}\right).
    \end{equation}
    Expandiendo en su serie de Taylor operacional:
    \begin{equation}
        \hat{T}(a)\psi(x) = \sum_{n=0}^\infty \frac{1}{n!}\left(-a\frac{\dif}{\dif x}\right)^n \psi(x) = \sum_{n=0}^\infty \frac{(-a)^n}{n!}\frac{\dif^n\psi(x)}{\dif x^n}.
    \end{equation}
    Esta suma infinita es precisamente la \textbf{serie de Taylor} de la función $\psi(x)$ desplazada en una distancia $-a$:
    \begin{equation}
        \hat{T}(a)\psi(x) = \psi(x - a)\footnote{Esta expansión en serie de Taylor asume rigurosamente que la función de onda $\psi(x)$ es analítica (por ejemplo, funciones gaussianas multiplicadas por polinomios), garantizando que su serie de potencias converge a $\psi(x-a)$. Por el Teorema de Stone, el operador traslación $\hat{T}(a) = \eu^{-\iu a\hat{p}_x/\hbar}$ queda unívocamente definido como unitario sobre la totalidad de $\mathcal{L}^2(\mathbb{R})$, aplicándose válidamente incluso sobre funciones no analíticas.}.
    \end{equation}
    ¡El resultado conecta de manera deslumbrante el cálculo diferencial clásico con el álgebra de operadores cuánticos!
\end{ejemplobox}

\begin{figure}[H]
\centering
\begin{subfigure}[t]{0.47\textwidth}
\centering
\begin{tikzpicture}
\begin{axis}[width=\linewidth,height=5.8cm,xlabel={$x$},
  xmin=-4,xmax=6,ymin=0,ymax=1.3,
  legend style={font=\scriptsize,at={(0.5,-0.28)},anchor=north,legend columns=2},ytick=\empty,axis lines=left]
  \addplot[uniblue,very thick,domain=-4:6,samples=200] {exp(-x^2/2)};
  \addlegendentry{$\psi(x)$}
  \addplot[unired,very thick,domain=-4:6,samples=200] {exp(-(x-2)^2/2)};
  \addlegendentry{$\hat T(a)\psi=\psi(x-a)$}
  \draw[-{Stealth},thick] (axis cs:0,1.08) -- (axis cs:2,1.08)
        node[midway,above] {$a>0$};
\end{axis}
\end{tikzpicture}
\caption{Con $a=2$: la funci\'on se desplaza hacia la \textbf{derecha}.}
\end{subfigure}\hfill
\begin{subfigure}[t]{0.52\textwidth}
\centering
\begin{tikzpicture}
\begin{axis}[width=\linewidth,height=5.8cm,xlabel={$x$},
  xmin=-3,xmax=4,ymin=-0.9,ymax=1.5,
  legend style={font=\scriptsize,at={(0.5,-0.28)},anchor=north,legend columns=3},
  grid=both,minor tick num=1]
  \addplot[black,very thick,domain=-3:4,samples=200] {exp(-(x-1)^2/2)};
  \addlegendentry{exacta $\psi(x-1)$}
  \addplot[uniblue,dashed,thick,domain=-3:4,samples=200] {exp(-x^2/2)};
  \addlegendentry{$N=0$}
  \addplot[unired,thick,domain=-3:4,samples=200] {exp(-x^2/2)*(1+x)};
  \addlegendentry{$N=1$}
  \addplot[unigreen,thick,domain=-3:4,samples=200]
    {exp(-x^2/2)*(1+x+(x^2-1)/2)};
  \addlegendentry{$N=2$}
  \addplot[orange,thick,domain=-3:4,samples=200]
    {exp(-x^2/2)*(1+x+(x^2-1)/2+(x^3-3*x)/6)};
  \addlegendentry{$N=3$}
\end{axis}
\end{tikzpicture}
\caption{Sumas parciales de $\sum_{n}\frac{(-a)^n}{n!}\psi^{(n)}(x)$, con $a=1$ y $\psi=\eu^{-x^2/2}$.}
\end{subfigure}
\caption{Acci\'on de $\hat T(a)=\eu^{-a\,\dif/\dif x}$: el desarrollo de Taylor reconstruye $\psi(x-a)$.}
\label{fig:traslacion}
\end{figure}

\subsection{Simetrías Continuas, Leyes de Conservación y el Teorema de Noether Cuántico}
\label{subsec:noether_cuantico}

El concepto de operador exponencial desarrollado en esta sección constituye la piedra angular de la física teórica moderna: \textbf{conecta directamente las simetrías geométricas del espacio-tiempo con las leyes de conservación de magnitudes físicas} \cite{zettili, sakurai}.

\subsubsection{Transformaciones Unitarias Infinitesimales}
Consideremos una transformación continua parametrizada por un escalar infinitesimal $\epsilon \ll 1$. Dicha transformación se implementa mediante un operador unitario $\hat{U}_\epsilon$ cercano a la identidad:
\begin{equation}
    \hat{U}_\epsilon = \hat{I} - \frac{\iu}{\hbar}\epsilon\,\hat{G},
    \label{eq:infinitesimal_unitary}
\end{equation}
donde $\hat{G}$ es un operador \textbf{hermítico} ($\hat{G}\adj = \hat{G}$), denominado el \textbf{generador de la transformación}. El factor $-\iu/\hbar$ se introduce por convención para que $\hat{G}$ posea las dimensiones físicas de la magnitud observable correspondiente y para garantizar la unitariedad:
\begin{equation}
    \hat{U}_\epsilon\adj \hat{U}_\epsilon = \left(\hat{I} + \frac{\iu}{\hbar}\epsilon\,\hat{G}\right)\left(\hat{I} - \frac{\iu}{\hbar}\epsilon\,\hat{G}\right) = \hat{I} + \mathcal{O}(\epsilon^2).
\end{equation}

Bajo esta transformación, un estado $\ket{\psi}$ y un observable $\hat{A}$ se modifican infinitesimalmente de acuerdo a:
\begin{align}
    \ket{\psi'} &= \hat{U}_\epsilon \ket{\psi} = \ket{\psi} - \frac{\iu}{\hbar}\epsilon\,\hat{G}\ket{\psi}, \\
    \hat{A}' &= \hat{U}_\epsilon\adj \hat{A} \hat{U}_\epsilon = \left(\hat{I} + \frac{\iu}{\hbar}\epsilon\,\hat{G}\right)\hat{A}\left(\hat{I} - \frac{\iu}{\hbar}\epsilon\,\hat{G}\right) \nonumber\\
    &= \hat{A} + \frac{\iu}{\hbar}\epsilon\left(\hat{G}\hat{A} - \hat{A}\hat{G}\right) + \mathcal{O}(\epsilon^2) = \hat{A} + \frac{\iu}{\hbar}\epsilon\,[\hat{G}, \hat{A}] + \mathcal{O}(\epsilon^2).
\end{align}
La variación infinitesimal del observable $\delta\hat{A} \equiv \hat{A}' - \hat{A}$ viene dada estrictamente por:
\begin{equation}
    \boxed{\delta\hat{A} = +\frac{\iu}{\hbar}\epsilon\,[\hat{G}, \hat{A}]}.
    \label{eq:infinitesimal_variation}
\end{equation}
(Nótese que para una traslación espacial con $\hat{G} = \hat{p}_x$ y $\epsilon = a$, tenemos $[\hat{p}_x, \hat{x}] = -\iu\hbar\hat{I}$, lo cual conduce de forma plenamente consistente a $\delta\hat{x} = +\frac{\iu}{\hbar}a(-\iu\hbar\hat{I}) = +a\hat{I}$, recuperando $\hat{x}' = \hat{x} + a\hat{I}$).

\subsubsection{Imagen de Heisenberg, Invariancia del Hamiltoniano y Constantes de Movimiento}
Supongamos que la dinámica de un sistema aislado está gobernada por un Hamiltoniano $\hat{H}$ independiente del tiempo. Existen dos formas equivalentes de describir la evolución temporal cuántica:
\begin{itemize}
    \item \textbf{Imagen de Schrödinger:} Los estados evolucionan en el tiempo mediante el operador unitario $\hat{U}(t) = \eu^{-\iu\hat{H}t/\hbar}$ ($\ket{\psi_S(t)} = \hat{U}(t)\ket{\psi(0)}$), mientras que los operadores observables $\hat{A}_S$ permanecen fijos.
    \item \textbf{Imagen de Heisenberg:} Los estados permanecen fijos ($\ket{\psi_H} = \ket{\psi(0)}$) y la dependencia temporal se traslada a los operadores observables mediante la transformación de semejanza unitaria:
    \begin{equation}
        \hat{A}_H(t) \equiv \hat{U}\adj(t)\hat{A}_S\hat{U}(t) = \eu^{+\frac{\iu}{\hbar}\hat{H}t}\,\hat{A}_S\,\eu^{-\frac{\iu}{\hbar}\hat{H}t},
    \end{equation}
    de modo que el valor esperado físico es idéntico en ambas imágenes: $\matrixel{\psi_S(t)}{\hat{A}_S}{\psi_S(t)} = \matrixel{\psi(0)}{\hat{A}_H(t)}{\psi(0)} = \ev{\hat{A}(t)}$.
\end{itemize}
Observemos que $\hat{A}_H(t) = \eu^{+\frac{\iu}{\hbar}\hat{H}t}\hat{A}_S\eu^{-\frac{\iu}{\hbar}\hat{H}t}$ posee exactamente la estructura del \textbf{Lema de Hadamard} \eqref{eq:hadamard_ansatz} con $\lambda = t$ y $\hat{A} \to \frac{\iu}{\hbar}\hat{H}$. Derivando respecto al tiempo $t$ mediante la Ec.~\eqref{eq:hadamard_first_diff} (e incluyendo cualquier dependencia temporal explícita $\partial\hat{A}_S/\partial t$):
\begin{equation}
    \frac{\dif\hat{A}_H(t)}{\dif t} = \left[\frac{\iu}{\hbar}\hat{H}, \, \hat{A}_H(t)\right] + \left(\frac{\partial\hat{A}_S}{\partial t}\right)_H \implies \boxed{\frac{\dif\hat{A}_H}{\dif t} = \frac{1}{\iu\hbar}[\hat{A}_H, \hat{H}] + \left(\frac{\partial\hat{A}}{\partial t}\right)_H}.
    \label{eq:heisenberg_eq_motion}
\end{equation}
Esta es la célebre \textbf{Ecuación de Movimiento de Heisenberg} (que aplicaremos en el Problema Resuelto 4 para deducir la precesión de Larmor).

Se dice que una transformación unitaria $\hat{U}$ generada por $\hat{G}$ representa una \textbf{simetría del sistema} si el Hamiltoniano es invariante ante la transformación:
\begin{equation}
    \hat{H}' = \hat{U}\adj \hat{H} \hat{U} = \hat{H} \iff [\hat{G}, \hat{H}] = 0.
    \label{eq:hamiltonian_invariance}
\end{equation}
Tomando el valor esperado de la Ec.~\eqref{eq:heisenberg_eq_motion} para el generador $\hat{G}$ en un estado arbitrario $\ket{\psi}$, obtenemos el \textbf{Teorema de Ehrenfest}:
\begin{equation}
    \frac{\dif}{\dif t}\ev{\hat{G}} = \frac{1}{\iu\hbar}\ev{[\hat{G}, \hat{H}]} + \left\langle \frac{\partial\hat{G}}{\partial t} \right\rangle.
    \label{eq:ehrenfest_theorem}
\end{equation}
Si el generador $\hat{G}$ no depende explícitamente del tiempo ($\partial\hat{G}/\partial t = 0$) y conmuta con el Hamiltoniano ($[\hat{G}, \hat{H}] = 0$), se deduce de forma inmediata:
\begin{equation}
    \frac{\dif}{\dif t}\ev{\hat{G}} = 0 \implies \boxed{\ev{\hat{G}} = \text{Constante en el tiempo}}.
\end{equation}

\begin{teoremabox}[Teorema de Noether Cuántico]
    A toda simetría continua del Hamiltoniano le corresponde un observable físico hermítico cuya media estadística es estrictamente constante en el tiempo (una \textbf{constante de movimiento}).
\end{teoremabox}

\subsubsection{El Puente Clásico-Cuántico: De los Corchetes de Poisson a los Conmutadores}
En la formulación hamiltoniana de la mecánica clásica, el estado de una partícula se describe en el \textbf{espacio de fase} de coordenadas y momentos $(x_i, p_i)$ (formalizado geométricamente como una variedad simpléctica dotada de la 2-forma $\omega = \dif x_i \wedge \dif p_i$; no es necesario conocer geometría diferencial para seguir este curso, basta con recordar las derivadas parciales ordinarias respecto a $x_i$ y $p_i$). La evolución temporal de una magnitud clásica $A(x, p)$ está gobernada por los \textbf{Corchetes de Poisson}:
\begin{equation}
    \frac{\dif A}{\dif t} = \{A, H\}_{\text{clásico}} + \frac{\partial A}{\partial t}, \qquad \{A, B\}_{\text{clásico}} \equiv \sum_{i=1}^3 \left( \frac{\partial A}{\partial x_i}\frac{\partial B}{\partial p_i} - \frac{\partial A}{\partial p_i}\frac{\partial B}{\partial x_i} \right).
\end{equation}
El célebre \textbf{Postulado de Cuantización Canónica de Dirac} establece que la estructura geométrica simpléctica clásica se traduce en la estructura algebraica de conmutadores en el espacio de Hilbert mediante el mapeo universal:
\begin{equation}
    \boxed{\{A, B\}_{\text{clásico}} \quad \xrightarrow{\quad \text{Cuantización Canónica} \quad} \quad \frac{1}{\iu\hbar}[\hat{A}, \hat{B}]}.
    \label{eq:dirac_canonical_quantization}
\end{equation}
Toda magnitud conservada en la mecánica clásica ($\{G, H\} = 0$) se transforma en cuántica en un observable compatible con la energía ($[\hat{G}, \hat{H}] = 0$), tal como se ilustra sintéticamente en la Figura~\ref{fig:simetrias}.

\begin{figure}[H]
\centering
\begin{tikzpicture}[x=1cm,y=1cm]
  \node[font=\small\bfseries] at (1.0,0.9) {Simetr\'ia};
  \node[font=\small\bfseries] at (5.2,0.9) {Generador};
  \node[font=\small\bfseries] at (9.0,0.9) {Magnitud conservada};
  \draw[uniblue,thick] (-0.8,0.55) -- (11,0.55);
  \begin{scope}[shift={(0,-0.6)}]
    \draw[thick,uniblue] (0,0) circle (0.45);
    \draw[thick,uniblue] (0,0) -- (0,0.3);
    \draw[thick,uniblue] (0,0) -- (0.22,-0.1);
    \draw[-{Stealth},thick,unired] (0.7,0.1) -- (1.2,0.1) node[right,font=\small,text=black,xshift=2pt] {$t\to t+\Delta t$};
    \node at (5.2,0) {$\hat H$};
    \node at (9.0,0) {Energ\'ia $E$};
  \end{scope}
  \draw[gray!50] (-0.8,-1.3) -- (11,-1.3);
  \begin{scope}[shift={(0,-2.1)}]
    \draw[thick,uniblue,fill=uniblue!15] (-0.3,-0.2) rectangle (0.1,0.2);
    \draw[thick,uniblue,dashed] (0.5,-0.2) rectangle (0.9,0.2);
    \draw[-{Stealth},thick,unired] (0.15,0.38) -- (0.85,0.38);
    \node[font=\small,right] at (1.3,0) {$\vec r\to\vec r+\vec a$};
    \node at (5.2,0) {$\hat{\vec p}$};
    \node at (9.0,0) {Momento lineal $\vec p$};
  \end{scope}
  \draw[gray!50] (-0.8,-2.8) -- (11,-2.8);
  \begin{scope}[shift={(0,-3.6)}]
    \draw[thick,uniblue] (0,0) circle (0.4);
    \fill[uniblue] (0.4,0) circle (0.05);
    \draw[-{Stealth},thick,unired] (0.55,0.1) arc (10:110:0.55);
    \node[font=\small,right] at (1.3,0) {$\theta\to\theta+\Delta\theta$};
    \node at (5.2,0) {$\hat{\vec L}$};
    \node at (9.0,0) {Momento angular $\vec L$};
  \end{scope}
  \draw[uniblue,thick] (-0.8,-4.3) -- (11,-4.3);
  \foreach \y in {-0.6,-2.1,-3.6}{
    \draw[-{Stealth},gray] (3.3,\y) -- (4.3,\y);
    \draw[-{Stealth},gray] (6.2,\y) -- (7.0,\y);
  }
  \node[font=\scriptsize,unigray,align=center] at (5.1,-4.8)
    {Homogeneidad del tiempo $\;\cdot\;$ homogeneidad del espacio $\;\cdot\;$ isotrop\'ia del espacio};
\end{tikzpicture}
\caption{Simetr\'ias continuas, generadores herm\'iticos y magnitudes conservadas (Teorema de Noether cu\'antico).}
\label{fig:simetrias}
\end{figure}

%% ============================================================
%%  SECCIÓN 5 — Momento Angular Tensorial, Álgebra so(3) y CCOC
%% ============================================================

\section{Momento Angular Tensorial, Álgebra de Lie \texorpdfstring{$\mathfrak{so}(3)$}{so(3)} y el CCOC}
\label{sec:momento_angular}

\subsection{Herramientas Tensoriales en \texorpdfstring{$\mathbb{R}^3$}{R3}: Notación de Índices, Kronecker y Levi-Civita}

Para abordar el cálculo de conmutadores en tres dimensiones espaciales sin recurrir a cálculos algebraicos interminables, la física teórica moderna emplea la \textbf{notación de índices} y el convenio de suma de Einstein \cite{sakurai, zettili}:
\begin{itemize}
    \item Los índices latinos $i, j, k, l, m, \dots \in \{1, 2, 3\}$ denotan las componentes espaciales cartesianas $(x, y, z) \equiv (x_1, x_2, x_3)$.
    \item \textbf{Convenio de Suma de Einstein:} Todo índice repetido en un mismo término monomial implica una sumatoria sobre sus tres valores posibles:
    \begin{equation}
        A_i B_i \equiv \sum_{i=1}^3 A_i B_i = A_1 B_1 + A_2 B_2 + A_3 B_3 = \vec{A}\cdot\vec{B}.
    \end{equation}
\end{itemize}

\subsubsection{El Delta de Kronecker \texorpdfstring{$\delta_{ij}$}{delta\_ij}}
Se define como el tensor de segundo orden:
\begin{equation}
    \delta_{ij} \equiv \begin{cases} 1, & \text{si } i = j, \\ 0, & \text{si } i \neq j. \end{cases}
\end{equation}
Actúa como operador de sustitución de índices: $\delta_{ij} A_j = A_i$, y su traza en 3D es $\delta_{ii} = \delta_{11} + \delta_{22} + \delta_{33} = 3$.

\subsubsection{El Pseudotensor Totalmente Antisimétrico de Levi-Civita \texorpdfstring{$\varepsilon_{ijk}$}{epsilon\_ijk}}
Se define como el tensor de permutación en tres dimensiones:
\begin{equation}
    \varepsilon_{ijk} \equiv \begin{cases}
        +1, & \text{si } (i, j, k) \text{ es una permutación par de } (1, 2, 3), \\
        -1, & \text{si } (i, j, k) \text{ es una permutación impar de } (1, 2, 3), \\
        0,  & \text{si cualesquiera dos índices son iguales}.
    \end{cases}
\end{equation}
Permite expresar el producto vectorial euclídeo de dos vectores $\vec{A} \times \vec{B}$ mediante sus componentes:
\begin{equation}
    (\vec{A} \times \vec{B})_i = \varepsilon_{ijk} A_j B_k.
\end{equation}

\begin{teoremabox}[Identidad de Contracción Fundamental de Levi-Civita]
    El producto de dos tensores de Levi-Civita con un índice contraído (sumado) se expresa en términos de deltas de Kronecker como:
    \begin{equation}
        \boxed{\varepsilon_{ijk}\varepsilon_{lmk} = \delta_{il}\delta_{jm} - \delta_{im}\delta_{jl}}.
        \label{eq:levi_civita_contract_1}
    \end{equation}
    Si además contraemos un segundo índice ($m = j$):
    \begin{equation}
        \varepsilon_{ijk}\varepsilon_{ljk} = \delta_{il}\delta_{jj} - \delta_{ij}\delta_{jl} = 3\delta_{il} - \delta_{il} = \boxed{2\delta_{il}}.
        \label{eq:levi_civita_contract_2}
    \end{equation}
\end{teoremabox}

\begin{alertabox}[¿De dónde surge el factor 2 en la contracción de dos índices?]
    Para comprender físicamente la Ec.~\eqref{eq:levi_civita_contract_2}, consideremos fijar $i = 1$ y $l = 1$. La suma sobre $j$ y $k$ recorre:
    \begin{equation}
        \sum_{j,k=1}^3 \varepsilon_{1jk}\varepsilon_{1jk} = \varepsilon_{123}\varepsilon_{123} + \varepsilon_{132}\varepsilon_{132} = (+1)(+1) + (-1)(-1) = 1 + 1 = 2.
    \end{equation}
    Los únicos pares de índices que no se anulan son aquellos donde $j$ y $k$ son diferentes entre sí y distintos de $i=1$. Existen exactamente $2! = 2$ permutaciones posibles de dichos índices restantes, cada una contribuyendo con $(+1)^2 = 1$ o $(-1)^2 = 1$. Por ello, el factor combinatorio resultante es exactamente $2$.
\end{alertabox}

\subsection{Definición Cuántica del Operador Momento Angular}

En mecánica clásica, el momento angular de una partícula puntual respecto al origen es $\vec{L} = \vec{r} \times \vec{p}$. Por el principio de correspondencia de Dirac, reemplazamos las coordenadas y momentos canónicos por operadores hermíticos:
\begin{equation}
    [\hat{x}_i, \hat{p}_j] = \iu\hbar\delta_{ij}\hat{I}, \quad [\hat{x}_i, \hat{x}_j] = 0, \quad [\hat{p}_i, \hat{p}_j] = 0.
    \label{eq:canonical_comm_3d}
\end{equation}
El operador momento angular orbital cuántico se define en componentes tensoriales por:
\begin{equation}
    \boxed{\hat{L}_k \equiv \varepsilon_{kij}\hat{x}_i\hat{p}_j}.
    \label{eq:l_operator_def}
\end{equation}
Desarrollando explícitamente para cada componente cartesiana:
\begin{align}
    \hat{L}_1 &= \hat{L}_x = \varepsilon_{123}\hat{x}_2\hat{p}_3 + \varepsilon_{132}\hat{x}_3\hat{p}_2 = \hat{x}_2\hat{p}_3 - \hat{x}_3\hat{p}_2 = \hat{y}\hat{p}_z - \hat{z}\hat{p}_y, \\
    \hat{L}_2 &= \hat{L}_y = \varepsilon_{231}\hat{x}_3\hat{p}_1 + \varepsilon_{213}\hat{x}_1\hat{p}_3 = \hat{x}_3\hat{p}_1 - \hat{x}_1\hat{p}_3 = \hat{z}\hat{p}_x - \hat{x}\hat{p}_z, \\
    \hat{L}_3 &= \hat{L}_z = \varepsilon_{312}\hat{x}_1\hat{p}_2 + \varepsilon_{321}\hat{x}_2\hat{p}_1 = \hat{x}_1\hat{p}_2 - \hat{x}_2\hat{p}_1 = \hat{x}\hat{p}_y - \hat{y}\hat{p}_x.
\end{align}
Notemos que como $i \neq j$ en cada producto $\hat{x}_i\hat{p}_j$, las componentes de posición y momento involucradas conmutan entre sí ($[\hat{x}_i, \hat{p}_j] = \iu\hbar\delta_{ij} = 0$ para $i \neq j$), lo que garantiza que $\hat{L}_k$ es un operador estrictamente \textbf{hermítico}:
\begin{equation}
    \hat{L}_k\adj = (\varepsilon_{kij}\hat{x}_i\hat{p}_j)\adj = \varepsilon_{kij}\hat{p}_j\adj\hat{x}_i\adj = \varepsilon_{kij}\hat{p}_j\hat{x}_i = \varepsilon_{kij}\hat{x}_i\hat{p}_j = \hat{L}_k.
\end{equation}

\subsection{Demostración Rigurosa del Álgebra de Lie \texorpdfstring{$\mathfrak{so}(3)$}{so(3)}}

Deseamos calcular el conmutador fundamental entre dos componentes arbitrarias $[\hat{L}_i, \hat{L}_j]$ sin recurrir al cálculo manual de las tres componentes por separado.

\begin{teoremabox}[Álgebra de Lie del Momento Angular]
    Las componentes del operador momento angular cuántico satisfacen las relaciones de conmutación:
    \begin{equation}
        \boxed{[\hat{L}_i, \hat{L}_j] = \iu\hbar\,\varepsilon_{ijk}\hat{L}_k}.
        \label{eq:so3_algebra}
    \end{equation}
\end{teoremabox}

\subsubsection{Del Cálculo Cartesiano Manual a la Deducción Tensorial General}

\begin{alertabox}[{Calentamiento: Cálculo Cartesiano Manual de $[\hat{L}_x, \hat{L}_y]$ frente al Método Tensorial}]
    Antes de efectuar la demostración general con índices para las nueve combinaciones $[\hat{L}_i, \hat{L}_j]$, evaluemos manualmente un caso particular: $[\hat{L}_x, \hat{L}_y] = [\hat{y}\hat{p}_z - \hat{z}\hat{p}_y, \, \hat{z}\hat{p}_x - \hat{x}\hat{p}_z]$.\\
    Por bilinealidad, este conmutador se expande en cuatro términos:
    \begin{equation}
        [\hat{L}_x, \hat{L}_y] = [\hat{y}\hat{p}_z, \hat{z}\hat{p}_x] - \bcancel{[\hat{y}\hat{p}_z, \hat{x}\hat{p}_z]} - \bcancel{[\hat{z}\hat{p}_y, \hat{z}\hat{p}_x]} + [\hat{z}\hat{p}_y, \hat{x}\hat{p}_z].
    \end{equation}
    Los dos términos centrales se anulan porque todas las coordenadas y momentos que contienen pertenecen a ejes distintos y conmutan entre sí. En el primer y cuarto término, la única pareja que no conmuta es $\hat{z}$ con $\hat{p}_z$ ($[\hat{p}_z, \hat{z}] = -\iu\hbar\hat{I}$ y $[\hat{z}, \hat{p}_z] = +\iu\hbar\hat{I}$), mientras que $\hat{x}, \hat{p}_x, \hat{y}, \hat{p}_y$ salen como factores espectadores:
    \begin{equation}
        [\hat{L}_x, \hat{L}_y] = \hat{y}\hat{p}_x [\hat{p}_z, \hat{z}] + \hat{x}\hat{p}_y [\hat{z}, \hat{p}_z] = \hat{y}\hat{p}_x(-\iu\hbar\hat{I}) + \hat{x}\hat{p}_y(+\iu\hbar\hat{I}) = \iu\hbar(\hat{x}\hat{p}_y - \hat{y}\hat{p}_x) = \iu\hbar\hat{L}_z.
    \end{equation}
    \textbf{¿Qué hace el método tensorial que veremos a continuación?} Exactamente esta misma operación para los 9 pares $(i, j)$ simultáneamente. Para no confundir la doble sumatoria interna de $\hat{L}_i$ con la doble sumatoria interna de $\hat{L}_j$, bautizamos con \textbf{nombres de índices mudos distintos} a cada una: $(a, b)$ para $\hat{L}_i = \varepsilon_{iab}\hat{x}_a\hat{p}_b$ y $(c, d)$ para $\hat{L}_j = \varepsilon_{jcd}\hat{x}_c\hat{p}_d$.
\end{alertabox}

\textbf{Demostración Tensorial Paso a Paso:}\\
Insertamos la definición tensorial \eqref{eq:l_operator_def} empleando índices mudos independientes para no mezclar las sumas:
\begin{equation}
    [\hat{L}_i, \hat{L}_j] = \left[\varepsilon_{iab}\hat{x}_a\hat{p}_b, \, \varepsilon_{jcd}\hat{x}_c\hat{p}_d\right] = \varepsilon_{iab}\varepsilon_{jcd}\left[\hat{x}_a\hat{p}_b, \, \hat{x}_c\hat{p}_d\right].
    \label{eq:comm_ll_tensor}
\end{equation}
Aplicamos la identidad de Leibniz para el conmutador de dos productos de operadores:
\begin{equation}
    [\hat{A}\hat{B}, \hat{C}\hat{D}] = \hat{A}[\hat{B}, \hat{C}]\hat{D} + \hat{A}\hat{C}[\hat{B}, \hat{D}] + [\hat{A}, \hat{C}]\hat{D}\hat{B} + \hat{C}[\hat{A}, \hat{D}]\hat{B}.
\end{equation}
Evaluando para nuestras variables canónicas, y recordando que $[\hat{x}_a, \hat{x}_c] = 0$ y $[\hat{p}_b, \hat{p}_d] = 0$:
\begin{align}
    [\hat{x}_a\hat{p}_b, \, \hat{x}_c\hat{p}_d] &= \hat{x}_a [\hat{p}_b, \hat{x}_c]\hat{p}_d + \hat{x}_c [\hat{x}_a, \hat{p}_d]\hat{p}_b \nonumber\\
    &= \hat{x}_a(-\iu\hbar\delta_{bc})\hat{p}_d + \hat{x}_c(\iu\hbar\delta_{ad})\hat{p}_b \nonumber\\
    &= -\iu\hbar\delta_{bc}\hat{x}_a\hat{p}_d + \iu\hbar\delta_{ad}\hat{x}_c\hat{p}_b.
    \label{eq:leibniz_expansion}
\end{align}
Sustituyendo esta expresión dentro de la Ec.~\eqref{eq:comm_ll_tensor}, el conmutador se separa en dos términos: $[\hat{L}_i, \hat{L}_j] = \mathcal{T}_1 + \mathcal{T}_2$.

\textit{Evaluación del Primer Término $\mathcal{T}_1$:}
\begin{align}
    \mathcal{T}_1 &= \varepsilon_{iab}\varepsilon_{jcd}(-\iu\hbar\delta_{bc})\hat{x}_a\hat{p}_d = -\iu\hbar\,\varepsilon_{iab}\varepsilon_{jbd}\hat{x}_a\hat{p}_d.
\end{align}
Invertimos los dos últimos índices del segundo Levi-Civita para alinearlo con la identidad de contracción ($\varepsilon_{jbd} = -\varepsilon_{jdb}$):
\begin{equation}
    \mathcal{T}_1 = \iu\hbar\,\varepsilon_{iab}\varepsilon_{jdb}\hat{x}_a\hat{p}_d.
\end{equation}
Contraemos sobre el índice repetido $b$ aplicando la identidad \eqref{eq:levi_civita_contract_1}:
\begin{equation}
    \varepsilon_{iab}\varepsilon_{jdb} = \delta_{ij}\delta_{ad} - \delta_{id}\delta_{aj}.
\end{equation}
Por ende:
\begin{align}
    \mathcal{T}_1 &= \iu\hbar\left(\delta_{ij}\delta_{ad} - \delta_{id}\delta_{aj}\right)\hat{x}_a\hat{p}_d \nonumber\\
    &= \iu\hbar\,\delta_{ij}(\hat{x}_a\hat{p}_a) - \iu\hbar\,\hat{x}_j\hat{p}_i.
    \label{eq:t1_eval}
\end{align}

\textit{Evaluación del Segundo Término $\mathcal{T}_2$:}
\begin{align}
    \mathcal{T}_2 &= \varepsilon_{iab}\varepsilon_{jcd}(\iu\hbar\delta_{ad})\hat{x}_c\hat{p}_b = \iu\hbar\,\varepsilon_{iab}\varepsilon_{jca}\hat{x}_c\hat{p}_b.
\end{align}
Reordenamos cíclicamente los índices para contraer sobre $a$: $\varepsilon_{iab} = \varepsilon_{abi}$, y $\varepsilon_{jca} = \varepsilon_{ajc} = -\varepsilon_{acj}$:
\begin{equation}
    \mathcal{T}_2 = -\iu\hbar\,\varepsilon_{abi}\varepsilon_{acj}\hat{x}_c\hat{p}_b.
\end{equation}
Aplicando la identidad de contracción sobre $a$:
\begin{equation}
    \varepsilon_{abi}\varepsilon_{acj} = \delta_{bc}\delta_{ij} - \delta_{bj}\delta_{ic}.
\end{equation}
Por ende:
\begin{align}
    \mathcal{T}_2 &= -\iu\hbar\left(\delta_{bc}\delta_{ij} - \delta_{bj}\delta_{ic}\right)\hat{x}_c\hat{p}_b \nonumber\\
    &= -\iu\hbar\,\delta_{ij}(\hat{x}_b\hat{p}_b) + \iu\hbar\,\hat{x}_i\hat{p}_j.
    \label{eq:t2_eval}
\end{align}

\textit{Suma de Términos $\mathcal{T}_1 + \mathcal{T}_2$:}
Sumamos miembro a miembro las Ecs.~\eqref{eq:t1_eval} y \eqref{eq:t2_eval}:
\begin{equation}
    [\hat{L}_i, \hat{L}_j] = \iu\hbar\,\delta_{ij}(\hat{x}_a\hat{p}_a) - \iu\hbar\,\hat{x}_j\hat{p}_i - \iu\hbar\,\delta_{ij}(\hat{x}_b\hat{p}_b) + \iu\hbar\,\hat{x}_i\hat{p}_j.
\end{equation}
Como los índices mudos $a$ y $b$ representan la misma suma $\vec{x}\cdot\vec{p}$, los términos diagonales proporcionales a $\delta_{ij}$ \textbf{se cancelan de forma idéntica}:
\begin{equation}
    [\hat{L}_i, \hat{L}_j] = \iu\hbar\left(\hat{x}_i\hat{p}_j - \hat{x}_j\hat{p}_i\right).
    \label{eq:l_comm_step}
\end{equation}
Reconocemos la definición tensorial del momento angular: $\varepsilon_{ijk}\hat{L}_k = \varepsilon_{ijk}\varepsilon_{kab}\hat{x}_a\hat{p}_b = (\delta_{ia}\delta_{jb} - \delta_{ib}\delta_{ja})\hat{x}_a\hat{p}_b = \hat{x}_i\hat{p}_j - \hat{x}_j\hat{p}_i$.
Sustituyendo esto en la Ec.~\eqref{eq:l_comm_step}:
\begin{equation}
    [\hat{L}_i, \hat{L}_j] = \iu\hbar\,\varepsilon_{ijk}\hat{L}_k.
\end{equation}
La deducción es exacta, rigurosa y general para cualquier par de componentes.

\subsubsection{Conexión con Álgebra Lineal: El Grupo $\mathrm{SO}(3)$ y Representación Hermítica}

\begin{alertabox}[Conexión con Álgebra Lineal: El Grupo $\mathrm{SO}(3)$ y los Generadores de Rotación]
    En Álgebra Lineal II y Geometría Afín, toda rotación en $\mathbb{R}^3$ alrededor de un eje viene dada por una matriz ortogonal $R \in \mathrm{SO}(3)$ ($\det R = +1$, $R^T R = I$), la cual puede escribirse como una exponencial matricial $R = \eu^{\theta \Omega}$, donde $\Omega$ es una matriz real \textbf{antisimétrica} ($\Omega^T = -\Omega$).
    
    El espacio vectorial de matrices antisimétricas de $3\times 3$ conforma el álgebra de Lie $\mathfrak{so}(3)$, cuya base canónica está formada por las tres matrices generadoras $J_1, J_2, J_3$ definidas por componentes mediante el pseudotensor de Levi-Civita:
    \begin{equation}
        (J_i)_{jk} \equiv -\varepsilon_{ijk}.
    \end{equation}
    Un cálculo directo de conmutadores matriciales arroja el álgebra geométrica clásica:
    \begin{equation}
        [J_i, J_j] = \varepsilon_{ijk} J_k.
    \end{equation}
    En mecánica cuántica, los observables físicos deben ser operadores autoadjuntos (hermíticos, $\hat{L}_i\adj = \hat{L}_i$). Como las matrices geométricas $J_i$ son antisimétricas reales ($J_i\adj = J_i^T = -J_i$), multiplicarlas por $+\iu\hbar$ las convierte rigurosamente en operadores hermíticos:
    \begin{equation}
        \hat{L}_i \equiv +\iu\hbar\,J_i \implies \hat{L}_i\adj = (-\iu\hbar)(J_i\adj) = (-\iu\hbar)(-J_i) = +\iu\hbar\,J_i = \hat{L}_i.
    \end{equation}
    \textbf{¿Por qué multiplicamos exactamente por $+\iu\hbar$? (Intuición Física y Álgebra Lineal):}
    \begin{itemize}
        \item \textbf{El factor imaginario $\iu$:} Toda matriz antisimétrica real posee autovalores puramente imaginarios ($\lambda \in \iu\mathbb{R}$). Al multiplicar por $\iu$, dichos autovalores se convierten en números estrictamente \textbf{reales} ($\iu \cdot \iu\lambda = -\lambda \in \mathbb{R}$), satisfaciendo la exigencia cuántica de que toda medición experimental en el laboratorio arroje números reales.
        \item \textbf{La constante de Planck $\hbar$:} Las matrices geométricas $J_i$ son magnitudes adimensionales (ángulos puros en radianes), mientras que el momento angular físico posee dimensiones de acción ($\mathrm{J}\cdot\mathrm{s}$). La constante $\hbar$ dota al generador de las unidades dinámicas del universo cuántico.
    \end{itemize}
    Multiplicando el conmutador clásico $[J_i, J_j] = \varepsilon_{ijk}J_k$ por $(\iu\hbar)^2 = -\hbar^2$:
    \begin{align}
        [\hat{L}_i, \hat{L}_j] &= [\iu\hbar\,J_i, \, \iu\hbar\,J_j] = (\iu\hbar)^2 [J_i, J_j] = (-\hbar^2)\varepsilon_{ijk}J_k \nonumber\\
        &= \iu\hbar\,\varepsilon_{ijk}(\iu\hbar\,J_k) = \iu\hbar\,\varepsilon_{ijk}\hat{L}_k.
    \end{align}
    Esta identificación es además directamente consistente con el operador de rotación finita: si escribimos la rotación espacial como $R = \eu^{\theta J}$ con $J = \frac{1}{\iu\hbar}\hat{L} = -\frac{\iu}{\hbar}\hat{L}$, recuperamos el operador unitario cuántico de rotación $R(\theta) = \eu^{-\iu\theta\hat{L}/\hbar}$.\\
    \textbf{Representación Vectorial:} Cabe subrayar que estas matrices $3\times 3$ constituyen la \textbf{representación vectorial o adjunta} (correspondiente al subespacio irreductible de momento angular $\ell=1$) del álgebra de Lie $\mathfrak{so}(3)$, a diferencia de los operadores diferenciales actuando sobre el espacio funcional de dimensión infinita $\mathcal{L}^2(\mathbb{R}^3)$, que representan la misma álgebra abstracta.\\
    ¡El álgebra cuántica del momento angular no es un postulado misterioso! Es la \textbf{representación hermítica del álgebra de rotaciones de $\mathbb{R}^3$} que se estudia en álgebra lineal.
\end{alertabox}

\subsection{La Paradoja del Producto Cruz Cuántico: \texorpdfstring{$\hat{\vec{L}} \times \hat{\vec{L}} = \iu\hbar\hat{\vec{L}} \neq \vec{0}$}{L x L = ihbar L}}

En el álgebra vectorial elemental de $\mathbb{R}^3$, el producto vectorial de cualquier vector consigo mismo es idénticamente nulo: $\vec{v} \times \vec{v} = \vec{0}$, debido a que $\sin(0^\circ) = 0$ y a la antisimetría del determinante.

¿Qué ocurre en el universo cuántico cuando evaluamos el producto vectorial del vector de operadores $\hat{\vec{L}}$ consigo mismo?
\begin{equation}
    (\hat{\vec{L}} \times \hat{\vec{L}})_k \equiv \varepsilon_{kij}\hat{L}_i\hat{L}_j.
\end{equation}
Dado que el índice $k$ está fijo, los índices $i$ y $j$ se suman sobre las parejas posibles. Utilizando el truco algebraico de intercambiar índices mudos:
\begin{equation}
    \varepsilon_{kij}\hat{L}_i\hat{L}_j = \frac{1}{2}\varepsilon_{kij}\hat{L}_i\hat{L}_j + \frac{1}{2}\varepsilon_{kji}\hat{L}_j\hat{L}_i.
\end{equation}
Por la antisimetría de Levi-Civita, $\varepsilon_{kji} = -\varepsilon_{kij}$:
\begin{equation}
    (\hat{\vec{L}} \times \hat{\vec{L}})_k = \frac{1}{2}\varepsilon_{kij}(\hat{L}_i\hat{L}_j - \hat{L}_j\hat{L}_i) = \frac{1}{2}\varepsilon_{kij}[\hat{L}_i, \hat{L}_j].
\end{equation}
¡El producto vectorial se expresa directamente en términos del conmutador de operadores! Sustituyendo el álgebra de Lie recién demostrada $[\hat{L}_i, \hat{L}_j] = \iu\hbar\,\varepsilon_{ijm}\hat{L}_m$:
\begin{align}
    (\hat{\vec{L}} \times \hat{\vec{L}})_k &= \frac{1}{2}\varepsilon_{kij}\left( \iu\hbar\,\varepsilon_{ijm}\hat{L}_m \right) = \frac{\iu\hbar}{2}\left( \varepsilon_{kij}\varepsilon_{mij} \right)\hat{L}_m.
\end{align}
Aplicando la contracción de dos índices \eqref{eq:levi_civita_contract_2}, $\varepsilon_{kij}\varepsilon_{mij} = 2\delta_{km}$:
\begin{equation}
    (\hat{\vec{L}} \times \hat{\vec{L}})_k = \frac{\iu\hbar}{2}\left( 2\delta_{km} \right)\hat{L}_m = \iu\hbar\,\delta_{km}\hat{L}_m = \iu\hbar\,\hat{L}_k.
\end{equation}
Reuniendo las tres componentes en notación vectorial tridimensional:
\begin{equation}
    \boxed{\hat{\vec{L}} \times \hat{\vec{L}} = \iu\hbar\,\hat{\vec{L}} \neq \vec{0}}.
    \label{eq:quantum_cross_product}
\end{equation}
Este resultado es una de las firmas más bellas y distintivas del carácter no conmutativo de la mecánica cuántica: el vector momento angular \enquote{no puede alinearse consigo mismo} en el sentido clásico, pues sus componentes cartesianas son incompatibles entre sí \cite{sakurai}.

\subsection{El Operador Casimir \texorpdfstring{$\hat{L}^2$}{L\textasciicircum 2}\ e Invariancia Rotacional}

Definimos el operador escalar del momento angular al cuadrado (denominado formalmente el \textbf{operador Casimir cuadrático} del grupo de rotaciones $\mathrm{SO}(3)$):
\begin{equation}
    \hat{L}^2 \equiv \hat{L}_x^2 + \hat{L}_y^2 + \hat{L}_z^2 = \hat{L}_j\hat{L}_j.
    \label{eq:casimir_def}
\end{equation}

\begin{teoremabox}[Conmutación del Operador Casimir]
    El operador $\hat{L}^2$ conmuta simultáneamente con cada una de las componentes del momento angular:
    \begin{equation}
        \boxed{[\hat{L}^2, \hat{L}_i] = 0, \quad \forall\,i \in \{1, 2, 3\}}.
        \label{eq:l2_comm}
    \end{equation}
\end{teoremabox}

\subsubsection{Conmutación del Operador Casimir y Derivación de Leibniz}

\textbf{Demostración:}\\
Calculemos el conmutador aplicando la regla de Leibniz para productos:
\begin{align}
    [\hat{L}^2, \hat{L}_i] &= [\hat{L}_j\hat{L}_j, \hat{L}_i] = \hat{L}_j[\hat{L}_j, \hat{L}_i] + [\hat{L}_j, \hat{L}_i]\hat{L}_j.
\end{align}
Sustituyendo $[\hat{L}_j, \hat{L}_i] = \iu\hbar\,\varepsilon_{jik}\hat{L}_k = -\iu\hbar\,\varepsilon_{ijk}\hat{L}_k$:
\begin{align}
    [\hat{L}^2, \hat{L}_i] &= \hat{L}_j(-\iu\hbar\,\varepsilon_{ijk}\hat{L}_k) + (-\iu\hbar\,\varepsilon_{ijk}\hat{L}_k)\hat{L}_j \nonumber\\
    &= -\iu\hbar\,\varepsilon_{ijk}\left( \hat{L}_j\hat{L}_k + \hat{L}_k\hat{L}_j \right) = -\iu\hbar\,\varepsilon_{ijk}\{\hat{L}_j, \hat{L}_k\}.
    \label{eq:casimir_step}
\end{align}

\subsubsection{Teorema de Anulación: Contracción Simétrico-Antisimétrico}

\begin{alertabox}[¿Por qué la contracción de un tensor simétrico con uno antisimétrico se anula idénticamente?]
    Examinemos con lupa el término resultante $\varepsilon_{ijk}\{\hat{L}_j, \hat{L}_k\}$ sumado sobre los índices mudos $j$ y $k$:\\
    1. El anticonmutador es una cantidad \textbf{estrictamente simétrica} ante el intercambio de índices:
    \begin{equation}
        S_{jk} \equiv \{\hat{L}_j, \hat{L}_k\} = \hat{L}_j\hat{L}_k + \hat{L}_k\hat{L}_j = \{\hat{L}_k, \hat{L}_j\} = S_{kj}.
    \end{equation}
    2. El tensor de Levi-Civita $\varepsilon_{ijk}$ es \textbf{estrictamente antisimétrico} ante el intercambio de $j$ y $k$:
    \begin{equation}
        \varepsilon_{ikj} = -\varepsilon_{ijk}.
    \end{equation}
    Sea la suma escalar $T_i \equiv \varepsilon_{ijk} S_{jk}$. Como $j$ y $k$ son índices mudos de sumatoria, podemos cambiarles el nombre sin alterar el valor de la suma ($j \to k$ y $k \to j$):
    \begin{equation}
        T_i = \varepsilon_{ikj} S_{kj}.
    \end{equation}
    Utilizando la simetría de $S$ y la antisimetría de $\varepsilon$:
    \begin{equation}
        T_i = (-\varepsilon_{ijk}) S_{jk} = -(\varepsilon_{ijk} S_{jk}) = -T_i.
    \end{equation}
    La única cantidad matemática que es igual a su propio negativo es el cero:
    \begin{equation}
        T_i = -T_i \implies 2T_i = 0 \implies \boxed{T_i = \varepsilon_{ijk}\{\hat{L}_j, \hat{L}_k\} \equiv 0}.
    \end{equation}
    Sustituyendo esto en la Ec.~\eqref{eq:casimir_step}, se concluye irrefutablemente que $[\hat{L}^2, \hat{L}_i] = 0$.
\end{alertabox}

\subsection{Conjunto Completo de Observables Compatibles (CCOC)}

\subsubsection{Incompatibilidad Cuántica y Relaciones de Incertidumbre Transversales}

En mecánica clásica, es posible medir simultáneamente las tres componentes del vector momento angular $(L_x, L_y, L_z)$ con precisión infinita. En mecánica cuántica, sin embargo:
\begin{equation}
    [\hat{L}_x, \hat{L}_y] = \iu\hbar\hat{L}_z \neq 0.
\end{equation}
Por el teorema general de incompatibilidad cuántica, para cualquier estado físico con momento angular no nulo ($\ell \ge 1$), \textbf{es imposible medir simultáneamente más de una componente del momento angular con precisión absoluta}. Si preparamos un autoestado de $\hat{L}_z$ con autovalor $m\hbar$, las componentes transversales $\hat{L}_x$ y $\hat{L}_y$ poseen incertidumbres intrínsecas gobernadas por la relación de Heisenberg:
\begin{equation}
    \Delta L_x \Delta L_y \ge \frac{\hbar}{2}|\ev{\hat{L}_z}| = \frac{\hbar^2 |m|}{2}.
\end{equation}
Nótese que para autoestados con $m = 0$ pertenecientes a multipletes con $\ell \ge 1$, si bien la cota inferior directa en la desigualdad se anula ($\ev{\hat{L}_z}=0$), las incertidumbres cuánticas transversales reales siguen siendo estrictamente positivas: $\Delta L_x = \Delta L_y = \hbar\sqrt{\ell(\ell+1)/2} > 0$ (como se demuestra rigurosamente en el Problema Resuelto~3). Solo para el estado $\ell = 0$ (donde $\hat{\vec{L}}\ket{0, 0} = \vec{0}$) las tres componentes son simultáneamente nítidas con $\Delta L_x = \Delta L_y = \Delta L_z = 0$.

\begin{figure}[H]
    \centering
    \begin{tikzpicture}[>=Stealth, font=\small, scale=1.0]
  \def\L{2.4495}  % sqrt(6)
  % esfera de radio |L|
  \draw[grisUNI!60,dashed] (0,0) circle (\L);
  % eje z
  \draw[->,grisUNI] (0,-3.1)--(0,3.1) node[above]{$L_z$};
  \foreach \m in {-2,-1,0,1,2}{
    \pgfmathsetmacro{\r}{sqrt(6-\m*\m)}
    % cono: circunferencia de precesión (vista en perspectiva)
    \draw[azulUNI!70,dashed] (0,\m) ellipse ({\r} and {0.22*\r});
    % vector L con L_z = m hbar
    \draw[->,very thick,azulUNI] (0,0)--(\r,\m);
    \fill[azulUNI] (\r,\m) circle (1.6pt);
    % marca en el eje
    \draw (-0.07,\m)--(0.07,\m);
    \node[anchor=east,font=\footnotesize] at (-0.1,\m) {$%
      \ifnum\m=0 0\else\ifnum\m=1 \hbar\else\ifnum\m=-1 -\hbar\else \m\hbar\fi\fi\fi$};
  }
  % anotaciones
  \node[anchor=west,font=\footnotesize,align=left] at (3.0,2.05)
       {$|\vec L|=\hbar\sqrt{\ell(\ell+1)}=\hbar\sqrt6\approx2{,}45\,\hbar$};
  \draw[->,grisUNI] (3.0,1.85)--(1.95,1.55);
  \node[anchor=west,font=\footnotesize,align=left] at (3.0,-2.0)
       {$L_x,\,L_y$ indeterminados:\\ cada vector gira sobre su cono};
  \draw[<->,rojoUNI] (-2.75,-2.0)--(-2.75,2.0);
  \node[rojoUNI,font=\footnotesize,rotate=90] at (-3.05,0)
       {$|L_z|\le 2\hbar < |\vec L|$};
\end{tikzpicture}

    \caption{Modelo vectorial para $\ell=2$: los cinco estados $|2,m\rangle$. La proyección $L_z=m\hbar$ está definida, mientras que $L_x$ y $L_y$ quedan indeterminados (el vector yace en algún punto de su cono). Como $\sqrt{\ell(\ell+1)}>\ell$, ningún estado alinea $\vec L$ con el eje $z$.}
    \label{fig:cono_L}
    \label{fig:conos}
\end{figure}


Sin embargo, como demostramos en la Ec.~\eqref{eq:l2_comm}, el operador magnitud al cuadrado $\hat{L}^2$ conmuta con \textit{todas} las componentes individuales:
\begin{equation}
    [\hat{L}^2, \hat{L}_x] = 0, \quad [\hat{L}^2, \hat{L}_y] = 0, \quad [\hat{L}^2, \hat{L}_z] = 0.
\end{equation}
Esto nos permite elegir $\hat{L}^2$ y \textbf{una única componente espacial} (por convención universal, el eje polar $\hat{L}_z$) para formar un conjunto compatible:
\begin{equation}
    [\hat{L}^2, \hat{L}_z] = 0.
\end{equation}

\subsubsection{Base Ortonormal Común $\{\ket{\ell, m}\}$ y Operadores Compatibles}

El teorema fundamental de observables que conmutan garantiza que existe una base ortonormal simultánea común de autoestados para ambos operadores, denotada por $\{\ket{\ell, m}\}$:
\begin{align}
    \hat{L}^2 \ket{\ell, m} &= \hbar^2 \ell(\ell+1)\ket{\ell, m}, \quad \ell \in \{0, 1, 2, \dots\}, \label{eq:l2_eigenvalues} \\
    \hat{L}_z \ket{\ell, m} &= \hbar m \ket{\ell, m}, \quad m \in \{-\ell, -\ell+1, \dots, \ell-1, \ell\}. \label{eq:lz_eigenvalues}
\end{align}

\begin{teoremabox}[Definición: Conjunto Completo de Observables Compatibles (CCOC)]
    Un conjunto de operadores hermíticos $\{\hat{A}, \hat{B}, \hat{C}, \dots\}$ constituye un \textbf{CCOC} (o CSCO por sus siglas en inglés, \textit{Complete Set of Compatible Observables}) si:
    \begin{enumerate}
        \item Todos los operadores conmutan mutuamente entre sí por pares ($[\hat{A}, \hat{B}] = 0$, etc.).
        \item Existe una base ortonormal común única donde cada combinación de autovalores $\{a, b, c, \dots\}$ especifica \textbf{un único estado físico}, eliminando por completo cualquier degeneración espectral residual.
    \end{enumerate}
\end{teoremabox}

\subsubsection{Jerarquía del CCOC: Del Potencial Central al Espín del Electrón}

En un potencial esféricamente simétrico central $V(r)$ (como el átomo de hidrógeno), el hamiltoniano $\hat{H} = \frac{\hat{p}^2}{2m} + V(r)$ es invariante ante rotaciones espaciales, lo que implica:
\begin{equation}
    [\hat{H}, \hat{L}^2] = 0 \quad \text{y} \quad [\hat{H}, \hat{L}_z] = 0.
\end{equation}
Para los grados de libertad orbitales (espaciales), el \textbf{CCOC del potencial central} viene dado por:
\begin{equation}
    \boxed{\text{CCOC}_{\text{orbital}} = \{\hat{H}, \hat{L}^2, \hat{L}_z\}},
\end{equation}
cuyos autoestados estacionarios comunes están biunívocamente identificados por los números cuánticos orbitales $\{n, \ell, m_\ell\}$ mediante el estado $\ket{n, \ell, m_\ell}$ \cite{cohen, zettili}.\\
Cuando se incorpora el grado de libertad del \textbf{espín del electrón} ($s = 1/2$), como los operadores orbitales conmutan con los operadores de espín ($[\hat{\vec{L}}, \hat{\vec{S}}] = 0$), el CCOC completo del átomo se expande naturalmente a:
\begin{equation}
    \boxed{\text{CCOC}_{\text{completo}} = \{\hat{H}, \hat{L}^2, \hat{L}_z, \hat{S}^2, \hat{S}_z\}},
\end{equation}
especificando sin ambigüedad la base ortonormal completa $\ket{n, \ell, m_\ell, s, m_s}$, donde $m_s = \pm 1/2$. (Cabe notar que para un electrón individual con espín intrínseco fijo $s=1/2$, el operador Casimir $\hat{S}^2 = s(s+1)\hbar^2\hat{I} = \frac{3}{4}\hbar^2\hat{I}$ es un múltiplo escalar constante de la identidad, por lo que el observable dinámico que levanta la degeneración de espín es $\hat{S}_z$).

\subsection{Operadores Escalera de Momento Angular \texorpdfstring{$\hat{L}_\pm$}{L+-} y Cuantización Espectral}
\label{subsec:operadores_escalera}

Para determinar algebraicamente los autovalores y autoestados del momento angular sin resolver ecuaciones diferenciales en derivadas parciales, introducimos los \textbf{operadores escalera} (o de subida y bajada) \cite{sakurai, zettili}:
\begin{equation}
    \boxed{\hat{L}_+ \equiv \hat{L}_x + \iu\hat{L}_y, \qquad \hat{L}_- \equiv \hat{L}_x - \iu\hat{L}_y}.
    \label{eq:ladder_def}
\end{equation}
Como $\hat{L}_x$ y $\hat{L}_y$ son hermíticos, cada operador escalera es el adjunto hermítico del otro:
\begin{equation}
    (\hat{L}_+)\adj = (\hat{L}_x + \iu\hat{L}_y)\adj = \hat{L}_x\adj - \iu\hat{L}_y\adj = \hat{L}_x - \iu\hat{L}_y = \hat{L}_-.
\end{equation}
Inversamente, podemos recuperar las componentes cartesianas mediante las combinaciones:
\begin{equation}
    \hat{L}_x = \frac{1}{2}\left(\hat{L}_+ + \hat{L}_-\right), \qquad \hat{L}_y = \frac{1}{2\iu}\left(\hat{L}_+ - \hat{L}_-\right).
\end{equation}

\subsubsection{Álgebra Fundamental de los Operadores Escalera}

Utilizando el álgebra de Lie $[\hat{L}_i, \hat{L}_j] = \iu\hbar\varepsilon_{ijk}\hat{L}_k$, evaluamos los conmutadores de $\hat{L}_\pm$ con los operadores del CCOC:
\begin{enumerate}
    \item \textbf{Conmutador con la componente polar $\hat{L}_z$:}
    \begin{align}
        [\hat{L}_z, \hat{L}_\pm] &= [\hat{L}_z, \hat{L}_x \pm \iu\hat{L}_y] = [\hat{L}_z, \hat{L}_x] \pm \iu[\hat{L}_z, \hat{L}_y] \nonumber\\
        &= (\iu\hbar\hat{L}_y) \pm \iu(-\iu\hbar\hat{L}_x) = \iu\hbar\hat{L}_y \pm \hbar\hat{L}_x = \pm\hbar(\hat{L}_x \pm \iu\hat{L}_y) \nonumber\\
        &= \boxed{\pm\hbar\hat{L}_\pm}.
        \label{eq:comm_lz_lpm}
    \end{align}
    \item \textbf{Conmutador con el operador Casimir $\hat{L}^2$:}
    Como $[\hat{L}^2, \hat{L}_x] = 0$ y $[\hat{L}^2, \hat{L}_y] = 0$, se deduce inmediatamente por linealidad:
    \begin{equation}
        \boxed{[\hat{L}^2, \hat{L}_\pm] = 0}.
        \label{eq:comm_l2_lpm}
    \end{equation}
    \item \textbf{Conmutador mutuo entre subida y bajada:}
    \begin{align}
        [\hat{L}_+, \hat{L}_-] &= [\hat{L}_x + \iu\hat{L}_y, \, \hat{L}_x - \iu\hat{L}_y] = -\iu[\hat{L}_x, \hat{L}_y] + \iu[\hat{L}_y, \hat{L}_x] \nonumber\\
        &= -\iu(\iu\hbar\hat{L}_z) + \iu(-\iu\hbar\hat{L}_z) = \hbar\hat{L}_z + \hbar\hat{L}_z = \boxed{2\hbar\hat{L}_z}.
        \label{eq:comm_lplus_lminus}
    \end{align}
\end{enumerate}

\subsubsection{Descomposición de $\hat{L}^2$ y Acción sobre la Base $\{\ket{\ell, m}\}$}

Multiplicando los operadores escalera en ambos órdenes:
\begin{align}
    \hat{L}_-\hat{L}_+ &= (\hat{L}_x - \iu\hat{L}_y)(\hat{L}_x + \iu\hat{L}_y) = \hat{L}_x^2 + \hat{L}_y^2 + \iu[\hat{L}_x, \hat{L}_y] = \hat{L}_x^2 + \hat{L}_y^2 - \hbar\hat{L}_z, \\
    \hat{L}_+\hat{L}_- &= (\hat{L}_x + \iu\hat{L}_y)(\hat{L}_x - \iu\hat{L}_y) = \hat{L}_x^2 + \hat{L}_y^2 - \iu[\hat{L}_x, \hat{L}_y] = \hat{L}_x^2 + \hat{L}_y^2 + \hbar\hat{L}_z.
\end{align}
Sumando $\hat{L}_z^2$ a ambos lados y reconociendo $\hat{L}^2 = \hat{L}_x^2 + \hat{L}_y^2 + \hat{L}_z^2$, obtenemos las dos identidades maestras:
\begin{equation}
    \boxed{\hat{L}^2 = \hat{L}_\mp\hat{L}_\pm + \hat{L}_z^2 \pm \hbar\hat{L}_z = \frac{1}{2}(\hat{L}_+\hat{L}_- + \hat{L}_-\hat{L}_+) + \hat{L}_z^2}.
    \label{eq:l2_ladder_decomp}
\end{equation}
Examinemos ahora cómo actúa $\hat{L}_\pm$ sobre un autoestado común $\ket{\ell, m}$. Utilizando la Ec.~\eqref{eq:comm_lz_lpm}:
\begin{equation}
    \hat{L}_z (\hat{L}_\pm\ket{\ell, m}) = (\hat{L}_\pm\hat{L}_z \pm \hbar\hat{L}_\pm)\ket{\ell, m} = (\hbar m \pm \hbar)\hat{L}_\pm\ket{\ell, m} = \hbar(m \pm 1)(\hat{L}_\pm\ket{\ell, m}).
\end{equation}
¡El operador $\hat{L}_+$ incrementa el número cuántico magnético $m$ en una unidad (+1), mientras que $\hat{L}_-$ lo decrementa en una unidad (-1), preservando intacto el valor de $\ell$ debido a que $[\hat{L}^2, \hat{L}_\pm] = 0$!\\
Para calcular la constante de proporcionalidad (normalización), evaluamos la norma al cuadrado:
\begin{align}
    \norm{\hat{L}_\pm\ket{\ell, m}}^2 &= \bra{\ell, m}\hat{L}_\pm\adj\hat{L}_\pm\ket{\ell, m} = \bra{\ell, m}\hat{L}_\mp\hat{L}_\pm\ket{\ell, m} \nonumber\\
    &= \matrixel{\ell, m}{\hat{L}^2 - \hat{L}_z^2 \mp \hbar\hat{L}_z}{\ell, m} \nonumber\\
    &= \hbar^2\left[\ell(\ell+1) - m^2 \mp m\right] = \hbar^2\left[\ell(\ell+1) - m(m \pm 1)\right].
\end{align}
Eligiendo la convención de fase estándar de Condon--Shortley (coeficientes reales y positivos):
\begin{equation}
    \boxed{\hat{L}_\pm\ket{\ell, m} = \hbar\sqrt{\ell(\ell+1) - m(m \pm 1)}\,\ket{\ell, m \pm 1}}.
    \label{eq:ladder_action}
\end{equation}
En particular, para el valor máximo $m = +\ell$ y el mínimo $m = -\ell$:
\begin{equation}
    \hat{L}_+\ket{\ell, \ell} = 0, \qquad \hat{L}_-\ket{\ell, -\ell} = 0,
\end{equation}
lo que demuestra que el espectro de proyecciones está acotado en el intervalo finito $m \in \{-\ell, -\ell+1, \dots, \ell-1, \ell\}$, conteniendo exactamente $(2\ell+1)$ estados independientes.

\begin{figure}[H]
\centering
\begin{tikzpicture}[>=Stealth, font=\small, lev/.style={thick}]
  \def\h{1.15}   % separación vertical entre niveles

  % ---------------- l = 1 ----------------
  \begin{scope}[xshift=0cm]
    \node[font=\normalsize] at (0,3.75) {$\ell=1$};
    \foreach \m in {-1,0,1}{
      \draw[lev] (-1.05,\m*\h)--(1.05,\m*\h);
      \node[anchor=east] at (-1.1,\m*\h) {$m=\m$};
    }
    \foreach \m/\c in {-1/{\hbar\sqrt2},0/{\hbar\sqrt2}}{
      \draw[->,thick,azulUNI]   (-0.4,\m*\h+0.06)--(-0.4,\m*\h+\h-0.06);
      \draw[->,thick,naranjaUNI](0.4,\m*\h+\h-0.06)--(0.4,\m*\h+0.06);
      \node[azulUNI,anchor=east,font=\scriptsize] at (-0.48,\m*\h+0.5*\h) {$\c$};
    }
    \node[azulUNI,font=\scriptsize,align=center] at (0,1*\h+0.75)
         {$\hat L_+|1,1\rangle=0$};
    \node[naranjaUNI,font=\scriptsize,align=center] at (0,-1*\h-0.75)
         {$\hat L_-|1,-1\rangle=0$};
  \end{scope}

  % ---------------- l = 2 ----------------
  \begin{scope}[xshift=6.8cm]
    \node[font=\normalsize] at (0,3.75) {$\ell=2$};
    \foreach \m in {-2,-1,0,1,2}{
      \draw[lev] (-1.05,\m*\h)--(1.05,\m*\h);
      \node[anchor=east] at (-1.1,\m*\h) {$m=\m$};
    }
    \foreach \m/\c in {-2/{2\hbar},-1/{\hbar\sqrt6},0/{\hbar\sqrt6},1/{2\hbar}}{
      \draw[->,thick,azulUNI]   (-0.4,\m*\h+0.06)--(-0.4,\m*\h+\h-0.06);
      \draw[->,thick,naranjaUNI](0.4,\m*\h+\h-0.06)--(0.4,\m*\h+0.06);
      \node[azulUNI,anchor=east,font=\scriptsize] at (-0.48,\m*\h+0.5*\h) {$\c$};
    }
    \node[azulUNI,font=\scriptsize] at (0,2*\h+0.35)
         {$\hat L_+|2,2\rangle=0$};
    \node[naranjaUNI,font=\scriptsize] at (0,-2*\h-0.35)
         {$\hat L_-|2,-2\rangle=0$};
  \end{scope}

  % leyenda
  \draw[->,thick,azulUNI]    (2.7,-3.55)--(3.3,-3.55);
  \node[anchor=west,font=\scriptsize] at (3.35,-3.55) {$\hat L_+$ (sube)};
  \draw[->,thick,naranjaUNI] (5.1,-3.55)--(5.7,-3.55);
  \node[anchor=west,font=\scriptsize] at (5.75,-3.55) {$\hat L_-$ (baja)};
\end{tikzpicture}
\caption{Operadores escalera. Cada flecha sube o baja una unidad de $m$ sin
cambiar $\ell$. El coeficiente sobre la flecha es
$\hbar\sqrt{\ell(\ell+1)-m(m\pm1)}$ y es el mismo para $\hat L_+$ (de $m$ a
$m+1$) y para $\hat L_-$ (de $m+1$ a $m$). Para $\ell=1$ son las entradas
$\hbar\sqrt2$ de las matrices~\eqref{eq:matrix_ladder_1}.}
\label{fig:escalera}
\end{figure}


\subsection{Representación Matricial Explícita para \texorpdfstring{$\ell = 1$}{l = 1} (Espacio 3D)}
\label{subsec:matrices_l1}

Para el valor $\ell = 1$, la dimensión del subespacio es $2\ell + 1 = 3$. La base ortonormal estándar, ordenada según autovalores decrecientes de $\hat{L}_z$, es:
\begin{equation}
    \mathcal{B} = \left\{ \ket{1, 1} \equiv \begin{pmatrix} 1 \\ 0 \\ 0 \end{pmatrix}, \quad \ket{1, 0} \equiv \begin{pmatrix} 0 \\ 1 \\ 0 \end{pmatrix}, \quad \ket{1, -1} \equiv \begin{pmatrix} 0 \\ 0 \\ 1 \end{pmatrix} \right\}.
\end{equation}
Calculamos elemento a elemento las matrices representantes de cada operador en esta base:
\begin{enumerate}
    \item \textbf{Matriz de $\hat{L}_z$:} Puesto que la base es propia de $\hat{L}_z$ ($\hat{L}_z\ket{1, m} = m\hbar\ket{1, m}$), su matriz es puramente diagonal:
    \begin{equation}
        \hat{L}_z = \hbar \begin{pmatrix} 1 & 0 & 0 \\ 0 & 0 & 0 \\ 0 & 0 & -1 \end{pmatrix}.
        \label{eq:matrix_lz_1}
    \end{equation}
    \item \textbf{Matrices Escalera $\hat{L}_+$ y $\hat{L}_-$:} Aplicando la Ec.~\eqref{eq:ladder_action}:
    \begin{align}
        \hat{L}_+\ket{1, 1} &= 0, \\
        \hat{L}_+\ket{1, 0} &= \hbar\sqrt{1(2) - 0(1)}\ket{1, 1} = \hbar\sqrt{2}\ket{1, 1}, \\
        \hat{L}_+\ket{1, -1} &= \hbar\sqrt{1(2) - (-1)(0)}\ket{1, 0} = \hbar\sqrt{2}\ket{1, 0}.
    \end{align}
    Por ende, las matrices escalera de $3\times 3$ resultan ser:
    \begin{equation}
        \hat{L}_+ = \hbar\sqrt{2}\begin{pmatrix} 0 & 1 & 0 \\ 0 & 0 & 1 \\ 0 & 0 & 0 \end{pmatrix}, \qquad \hat{L}_- = (\hat{L}_+)\adj = \hbar\sqrt{2}\begin{pmatrix} 0 & 0 & 0 \\ 1 & 0 & 0 \\ 0 & 1 & 0 \end{pmatrix}.
        \label{eq:matrix_ladder_1}
    \end{equation}
    \item \textbf{Matrices Cartesianas $\hat{L}_x$ y $\hat{L}_y$:}
    \begin{align}
        \hat{L}_x &= \frac{1}{2}(\hat{L}_+ + \hat{L}_-) = \frac{\hbar}{\sqrt{2}}\begin{pmatrix} 0 & 1 & 0 \\ 1 & 0 & 1 \\ 0 & 1 & 0 \end{pmatrix}, \label{eq:matrix_lx_1} \\
        \hat{L}_y &= \frac{1}{2\iu}(\hat{L}_+ - \hat{L}_-) = \frac{\hbar}{\sqrt{2}}\begin{pmatrix} 0 & -\iu & 0 \\ \iu & 0 & -\iu \\ 0 & \iu & 0 \end{pmatrix}. \label{eq:matrix_ly_1}
    \end{align}
    \item \textbf{Matriz del Casimir $\hat{L}^2$:}
    Puesto que $\hat{L}^2\ket{1, m} = \hbar^2 1(1+1)\ket{1, m} = 2\hbar^2\ket{1, m}$, es proporcional a la identidad $3\times 3$:
    \begin{equation}
        \hat{L}^2 = 2\hbar^2 \begin{pmatrix} 1 & 0 & 0 \\ 0 & 1 & 0 \\ 0 & 0 & 1 \end{pmatrix} = 2\hbar^2 \hat{I}_3.
    \end{equation}
\end{enumerate}

\begin{alertabox}[Verificación Matricial Directa del Álgebra de Conmutadores]
    Comprobemos matricialmente el conmutador $[\hat{L}_x, \hat{L}_y]$:
    \begin{align}
        \hat{L}_x\hat{L}_y &= \frac{\hbar^2}{2}\begin{pmatrix} 0 & 1 & 0 \\ 1 & 0 & 1 \\ 0 & 1 & 0 \end{pmatrix}\begin{pmatrix} 0 & -\iu & 0 \\ \iu & 0 & -\iu \\ 0 & \iu & 0 \end{pmatrix} = \frac{\hbar^2}{2}\begin{pmatrix} \iu & 0 & -\iu \\ 0 & 0 & 0 \\ \iu & 0 & -\iu \end{pmatrix}, \\
        \hat{L}_y\hat{L}_x &= \frac{\hbar^2}{2}\begin{pmatrix} 0 & -\iu & 0 \\ \iu & 0 & -\iu \\ 0 & \iu & 0 \end{pmatrix}\begin{pmatrix} 0 & 1 & 0 \\ 1 & 0 & 1 \\ 0 & 1 & 0 \end{pmatrix} = \frac{\hbar^2}{2}\begin{pmatrix} -\iu & 0 & -\iu \\ 0 & 0 & 0 \\ \iu & 0 & \iu \end{pmatrix}.
    \end{align}
    Restando ambas matrices:
    \begin{equation}
        [\hat{L}_x, \hat{L}_y] = \frac{\hbar^2}{2}\begin{pmatrix} 2\iu & 0 & 0 \\ 0 & 0 & 0 \\ 0 & 0 & -2\iu \end{pmatrix} = \iu\hbar\left[ \hbar\begin{pmatrix} 1 & 0 & 0 \\ 0 & 0 & 0 \\ 0 & 0 & -1 \end{pmatrix} \right] = \iu\hbar\hat{L}_z.
    \end{equation}
    ¡El álgebra de Lie se satisface con total precisión aritmética en la representación matricial tridimensional!
\end{alertabox}

\subsection{El Hamiltoniano Central, el Laplaciano Esférico y la Separación de Variables}
\label{subsec:hamiltoniano_central_laplaciano}

En coordenadas esféricas $(r, \theta, \phi)$, el operador laplaciano $\nabla^2$ se descompone naturalmente en una parte radial y una parte angular:
\begin{equation}
    \nabla^2 = \frac{1}{r^2}\frac{\partial}{\partial r}\left(r^2\frac{\partial}{\partial r}\right) + \frac{1}{r^2}\left[ \frac{1}{\sin\theta}\frac{\partial}{\partial\theta}\left(\sin\theta\frac{\partial}{\partial\theta}\right) + \frac{1}{\sin^2\theta}\frac{\partial^2}{\partial\phi^2} \right].
    \label{eq:laplaciano_esferico}
\end{equation}
Reconociendo la representación diferencial del operador Casimir $\hat{L}^2$ en coordenadas esféricas:
\begin{equation}
    \hat{L}^2 = -\hbar^2\left[ \frac{1}{\sin\theta}\frac{\partial}{\partial\theta}\left(\sin\theta\frac{\partial}{\partial\theta}\right) + \frac{1}{\sin^2\theta}\frac{\partial^2}{\partial\phi^2} \right], \qquad \hat{L}_z = -\iu\hbar\frac{\partial}{\partial\phi}.
\end{equation}
el Laplaciano se reescribe de forma sumamente compacta y elegante como:
\begin{equation}
    \boxed{\nabla^2 = \frac{1}{r^2}\frac{\partial}{\partial r}\left(r^2\frac{\partial}{\partial r}\right) - \frac{\hat{L}^2}{\hbar^2 r^2}}.
    \label{eq:laplaciano_con_L2}
\end{equation}
Sustituyendo esto en el Hamiltoniano de una partícula de masa reducida $\mu$ sometida a un potencial esférico central $V(r)$:
\begin{equation}
    \boxed{\hat{H} = -\frac{\hbar^2}{2\mu}\frac{1}{r^2}\frac{\partial}{\partial r}\left(r^2\frac{\partial}{\partial r}\right) + \frac{\hat{L}^2}{2\mu r^2} + V(r)}.
    \label{eq:hamiltoniano_central_expandido}
\end{equation}

\begin{teoremabox}[Invariancia y Conmutación en Potenciales Centrales]
    Para cualquier potencial con simetría esférica $V(r)$, el Hamiltoniano conmuta con $\hat{L}^2$ y con todas las componentes $\hat{L}_i$:
    \begin{equation}
        [\hat{H}, \hat{L}^2] = 0, \qquad [\hat{H}, \hat{L}_z] = 0.
    \end{equation}
\end{teoremabox}

\textbf{Demostración Didáctica:}\\
1. Los operadores radiales $\frac{\partial}{\partial r}$ y las funciones del radio $r$ y $V(r)$ actúan exclusivamente sobre la coordenada $r$, por lo que conmutan idénticamente con los operadores diferenciales angulares $\hat{L}^2$ y $\hat{L}_z$ que sólo actúan sobre $(\theta, \phi)$.\\
2. Por otro lado, la única dependencia angular del Hamiltoniano $\hat{H}$ en la Ec.~\eqref{eq:hamiltoniano_central_expandido} está contenida en el término $\frac{\hat{L}^2}{2\mu r^2}$. Como $[\hat{L}^2, \hat{L}^2] = 0$ y $[\hat{L}^2, \hat{L}_z] = 0$ (demostrado en la Ec.~\eqref{eq:l2_comm}), concluimos de inmediato que:
\begin{equation}
    [\hat{H}, \hat{L}^2] = 0 \quad \text{y} \quad [\hat{H}, \hat{L}_z] = 0.
\end{equation}
El Teorema 2.3 garantiza la existencia de una base ortonormal común de autoestados simultáneos para el CCOC $\{\hat{H}, \hat{L}^2, \hat{L}_z\}$. Debido a la estructura diferencial del laplaciano en coordenadas esféricas (donde la parte angular coincide con $-\hat{L}^2/(\hbar^2 r^2)$), el método de separación de variables conduce de forma natural a la factorización espacial de las funciones de onda:
\begin{equation}
    \psi_{n\ell m}(r, \theta, \phi) = R_{n\ell}(r)\,Y_\ell^m(\theta, \phi),
\end{equation}
donde los \textbf{armónicos esféricos} $Y_\ell^m(\theta, \phi)$ son los autoestados comunes de $\hat{L}^2$ y $\hat{L}_z$, y la función radial $R_{n\ell}(r)$ satisface la ecuación diferencial 1D:
\begin{equation}
    \left[ -\frac{\hbar^2}{2\mu}\frac{1}{r^2}\frac{\dif}{\dif r}\left(r^2\frac{\dif}{\dif r}\right) + \frac{\hbar^2\ell(\ell+1)}{2\mu r^2} + V(r) \right] R_{n\ell}(r) = E_{n\ell} R_{n\ell}(r).
\end{equation}
El término $\frac{\hbar^2\ell(\ell+1)}{2\mu r^2}$ representa la \textbf{barrera centrífuga cuántica}: una fuerza repulsiva de origen cinemático que aleja a la partícula del origen para todo estado con momento angular no nulo ($\ell \ge 1$).

\subsection{El Producto Tensorial \texorpdfstring{$\mathcal{H}_1 \otimes \mathcal{H}_2$}{H1 x H2} y el Acoplamiento de Momentos Angulares}
\label{subsec:producto_tensorial_acoplamiento}

Cuando consideramos sistemas físicos compuestos por múltiples partículas (por ejemplo, dos electrones) o grados de libertad independientes de una misma partícula (como el movimiento espacial orbital acoplado al espín intrínseco), el espacio de estados global se construye mediante el \textbf{producto tensorial de espacios de Hilbert} \cite{sakurai, cohen}.

\begin{teoremabox}[Definición Formal: Espacio Producto Tensorial $\mathcal{H}_1 \otimes \mathcal{H}_2$]
    Sean $\mathcal{H}_1$ y $\mathcal{H}_2$ dos espacios de Hilbert de dimensiones $N_1$ y $N_2$ con bases ortonormales $\{\ket{u_i}\}_{i=1}^{N_1}$ y $\{\ket{v_j}\}_{j=1}^{N_2}$ respectivamente. El espacio producto tensorial $\mathcal{H} = \mathcal{H}_1 \otimes \mathcal{H}_2$ satisface:
    \begin{enumerate}
        \item \textbf{Bilinealidad:} Para todo $\alpha, \beta \in \mathbb{C}$:
        \begin{equation}
            (\alpha\ket{u_1} + \beta\ket{u_2}) \otimes \ket{v} = \alpha(\ket{u_1} \otimes \ket{v}) + \beta(\ket{u_2} \otimes \ket{v}).
        \end{equation}
        \item \textbf{Producto Interno Compuesto:} Para estados producto $\ket{\Phi} = \ket{\phi_1} \otimes \ket{\phi_2}$ y $\ket{\Psi} = \ket{\psi_1} \otimes \ket{\psi_2}$:
        \begin{equation}
            \braket{\Phi}{\Psi} = \braket{\phi_1}{\psi_1}_{\mathcal{H}_1} \cdot \braket{\phi_2}{\psi_2}_{\mathcal{H}_2}.
        \end{equation}
        \item \textbf{Dimensión Total:} La dimensión es el producto de las dimensiones:
        \begin{equation}
            \dim(\mathcal{H}_1 \otimes \mathcal{H}_2) = \dim(\mathcal{H}_1) \cdot \dim(\mathcal{H}_2) = N_1 \cdot N_2.
        \end{equation}
        \item \textbf{Base Producto Desacoplada:} El conjunto de pares ordenados $\{\ket{u_i} \otimes \ket{v_j}\} \equiv \{\ket{u_i, v_j}\}$ conforma una base ortonormal completa de $\mathcal{H}_1 \otimes \mathcal{H}_2$, satisfaciendo la relación de clausura:
        \begin{equation}
            \sum_{i=1}^{N_1} \sum_{j=1}^{N_2} \ket{u_i, v_j}\bra{u_i, v_j} = \hat{I}_1 \otimes \hat{I}_2 = \hat{I}.
        \end{equation}
        \item \textbf{Producto de Operadores Compuestos (Propiedad Fundamental):}
        \begin{equation}
            \boxed{(\hat{A}_1 \otimes \hat{B}_1)(\hat{A}_2 \otimes \hat{B}_2) = (\hat{A}_1\hat{A}_2) \otimes (\hat{B}_1\hat{B}_2)}.
            \label{eq:operator_tensor_product_rule}
        \end{equation}
    \end{enumerate}
\end{teoremabox}

\subsubsection{Motivación y Necesidad Física: El Acoplamiento Espín-Órbita ($LS$)}
\label{subsec:acoplamiento_espin_orbita}

¿Por qué resulta indispensable en la física cuántica real sumar momentos angulares y abandonar la descripción individual de cada subsistema?

En un átomo idealizado sin interacciones magnéticas relativistas internas, el Hamiltoniano no perturbado $\hat{H}_0$ conmuta independientemente tanto con el momento angular orbital $\hat{\vec{L}}$ como con el espín intrínseco $\hat{\vec{S}}$:
\begin{equation}
    [\hat{H}_0, \hat{\vec{L}}] = 0 \quad \text{y} \quad [\hat{H}_0, \hat{\vec{S}}] = 0.
\end{equation}
En dicho régimen simplificado, las componentes individuales $\hat{L}_z$ y $\hat{S}_z$ se conservan de forma separada en el tiempo.

Sin embargo, en el mundo real, un electrón que orbita a velocidades relativistas alrededor del núcleo experimenta en su propio sistema de referencia en reposo un intenso campo magnético efectivo generado por el movimiento aparente del núcleo cargado. Dicho campo magnético interno interactúa con el momento magnético dipolar intrínseco del electrón, generando el término de \textbf{acoplamiento espín-órbita}:
\begin{equation}
    \hat{H}_{SO} = \xi(r)\,\hat{\vec{L}}\cdot\hat{\vec{S}} = \xi(r)\left( \hat{L}_x\hat{S}_x + \hat{L}_y\hat{S}_y + \hat{L}_z\hat{S}_z \right), \quad \xi(r) > 0.
    \label{eq:h_so_def}
\end{equation}
Al evaluar la conmutación de este término con las componentes orbitales individuales:
\begin{align}
    [\hat{H}_{SO}, \hat{L}_i] &= \xi(r)[\hat{L}_j\hat{S}_j, \, \hat{L}_i] = \xi(r)[\hat{L}_j, \hat{L}_i]\hat{S}_j = \xi(r)(-\iu\hbar\,\varepsilon_{jik}\hat{L}_k)\hat{S}_j \nonumber\\
    &= \iu\hbar\,\xi(r)\,\varepsilon_{ijk}\hat{L}_j\hat{S}_k \neq 0.
\end{align}
Análogamente, para el espín: $[\hat{H}_{SO}, \hat{S}_i] = -\iu\hbar\,\xi(r)\,\varepsilon_{ijk}\hat{L}_j\hat{S}_k \neq 0$.

\begin{alertabox}[Consecuencia Física Fundamental: Ruptura de Conservación Aislada]
    Debido a la interacción espín-órbita ($\propto \hat{\vec{L}}\cdot\hat{\vec{S}}$), el momento angular orbital $\hat{\vec{L}}$ y el espín $\hat{\vec{S}}$ \textbf{dejan de conmutar individualmente con el Hamiltoniano total}:
    \begin{equation}
        [\hat{H}_{\text{total}}, \hat{\vec{L}}] \neq 0 \quad \text{y} \quad [\hat{H}_{\text{total}}, \hat{\vec{S}}] \neq 0.
    \end{equation}
    Por consiguiente, por el Teorema de Ehrenfest:
    \begin{equation}
        \frac{\dif}{\dif t}\ev{\hat{\vec{L}}} \neq 0, \qquad \frac{\dif}{\dif t}\ev{\hat{\vec{S}}} \neq 0.
    \end{equation}
    \textbf{¡Ni el momento orbital ni el espín se conservan de manera aislada!} Un torque magnético interno transfiere momento angular de forma continua entre el movimiento orbital y el espín.
\end{alertabox}

\subsubsection{Suma de Momentos Angulares y el Operador Momento Angular Total $\hat{\vec{J}}$}

Sin embargo, al sumar ambas variaciones temporales:
\begin{equation}
    [\hat{H}_{SO}, \, \hat{L}_i + \hat{S}_i] = \iu\hbar\,\xi(r)\,\varepsilon_{ijk}\hat{L}_j\hat{S}_k - \iu\hbar\,\xi(r)\,\varepsilon_{ijk}\hat{L}_j\hat{S}_k = 0.
\end{equation}
¡La suma total conmuta exactamente con el Hamiltoniano! Esto motiva la definición del \textbf{operador momento angular total}:
\begin{equation}
    \boxed{\hat{\vec{J}} \equiv \hat{\vec{L}} + \hat{\vec{S}} = \hat{\vec{L}} \otimes \hat{I}_S + \hat{I}_L \otimes \hat{\vec{S}}}, \qquad [\hat{H}_{\text{total}}, \hat{\vec{J}}] = 0.
\end{equation}
En un contexto más general donde se acoplan dos subsistemas cualesquiera (dos electrones, o dos grados de libertad independientes) con operadores $\hat{\vec{J}}_1$ y $\hat{\vec{J}}_2$:
\begin{equation}
    [\hat{J}_{1i}, \hat{J}_{1j}] = \iu\hbar\,\varepsilon_{ijk}\hat{J}_{1k}, \qquad [\hat{J}_{2i}, \hat{J}_{2j}] = \iu\hbar\,\varepsilon_{ijk}\hat{J}_{2k}, \qquad [\hat{J}_{1i} \otimes \hat{I}_2, \, \hat{I}_1 \otimes \hat{J}_{2j}] = 0.
\end{equation}
Definimos el \textbf{operador momento angular total} como:
\begin{equation}
    \boxed{\hat{\vec{J}} \equiv \hat{\vec{J}}_1 \otimes \hat{I}_2 + \hat{I}_1 \otimes \hat{\vec{J}}_2 \equiv \hat{\vec{J}}_1 + \hat{\vec{J}}_2}.
\end{equation}
Comprobemos que el operador suma satisface el álgebra de Lie del momento angular empleando la regla del producto tensorial \eqref{eq:operator_tensor_product_rule}:
\begin{align}
    [\hat{J}_i, \hat{J}_j] &= [\hat{J}_{1i} + \hat{J}_{2i}, \, \hat{J}_{1j} + \hat{J}_{2j}] = [\hat{J}_{1i}, \hat{J}_{1j}] + [\hat{J}_{2i}, \hat{J}_{2j}] + \underbrace{[\hat{J}_{1i}, \hat{J}_{2j}]}_{0} + \underbrace{[\hat{J}_{2i}, \hat{J}_{1j}]}_{0} \nonumber\\
    &= \iu\hbar\,\varepsilon_{ijk}\hat{J}_{1k} + \iu\hbar\,\varepsilon_{ijk}\hat{J}_{2k} = \iu\hbar\,\varepsilon_{ijk}(\hat{J}_{1k} + \hat{J}_{2k}) = \boxed{\iu\hbar\,\varepsilon_{ijk}\hat{J}_k}.
\end{align}
¡El operador total $\hat{\vec{J}}$ es un momento angular genuino y riguroso!

\subsubsection{Incompatibilidad entre la Base Desacoplada y la Base Acoplada}

Consideremos el operador Casimir total al cuadrado:
\begin{equation}
    \hat{J}^2 = (\hat{\vec{J}}_1 + \hat{\vec{J}}_2)^2 = \hat{J}_1^2 + \hat{J}_2^2 + 2\hat{\vec{J}}_1\cdot\hat{\vec{J}}_2 = \hat{J}_1^2 + \hat{J}_2^2 + 2\hat{J}_{1z}\hat{J}_{2z} + \hat{J}_{1+}\hat{J}_{2-} + \hat{J}_{1-}\hat{J}_{2+}.
    \label{eq:j2_expansion_coupled}
\end{equation}
Evaluemos las relaciones de compatibilidad de este operador con los operadores individuales:
\begin{itemize}
    \item Conmuta con las magnitudes individuales al cuadrado: $[\hat{J}^2, \hat{J}_1^2] = 0$ y $[\hat{J}^2, \hat{J}_2^2] = 0$.
    \item Conmuta con la componente polar total: $[\hat{J}^2, \hat{J}_z] = 0$, donde $\hat{J}_z = \hat{J}_{1z} + \hat{J}_{2z}$.
    \item \textbf{¡NO conmuta con las componentes polares individuales $\hat{J}_{1z}$ y $\hat{J}_{2z}$!}\\
    En efecto, calculando el conmutador con $\hat{J}_{1z}$:
    \begin{align}
        [\hat{J}^2, \hat{J}_{1z}] &= [\hat{J}_1^2 + \hat{J}_2^2 + 2(\hat{J}_{1x}\hat{J}_{2x} + \hat{J}_{1y}\hat{J}_{2y} + \hat{J}_{1z}\hat{J}_{2z}), \, \hat{J}_{1z}] \nonumber\\
        &= 2[\hat{J}_{1x}, \hat{J}_{1z}]\hat{J}_{2x} + 2[\hat{J}_{1y}, \hat{J}_{1z}]\hat{J}_{2y} = 2(-\iu\hbar\hat{J}_{1y})\hat{J}_{2x} + 2(\iu\hbar\hat{J}_{1x})\hat{J}_{2y} \nonumber\\
        &= \boxed{2\iu\hbar\left(\hat{J}_{1x}\hat{J}_{2y} - \hat{J}_{1y}\hat{J}_{2x}\right) \neq 0 \quad (\text{en general})}.
    \end{align}
\end{itemize}
Esta no conmutatividad es el origen de los dos esquemas de descripción cuántica de un sistema compuesto:
\begin{enumerate}
    \item \textbf{Base Desacoplada $\{\ket{j_1, m_1; j_2, m_2}\}$:} Es la base simultánea del CCOC desacoplado $\{\hat{J}_1^2, \hat{J}_2^2, \hat{J}_{1z}, \hat{J}_{2z}\}$. En esta base, las proyecciones individuales $m_1$ y $m_2$ están perfectamente definidas, pero la magnitud total del momento angular $J$ está completamente indeterminada.
    \item \textbf{Base Acoplada $\{\ket{j, M; j_1, j_2}\}$:} Es la base simultánea del CCOC acoplado $\{\hat{J}_1^2, \hat{J}_2^2, \hat{J}^2, \hat{J}_z\}$. En esta base, la magnitud total $J$ y la proyección total $M = m_1 + m_2$ son nítidas, pero las componentes individuales $m_1$ y $m_2$ no tienen valor definido.
\end{enumerate}

\subsubsection{Regla de Suma de Momentos Angulares y Coeficientes de Clebsch--Gordan}

Para dos momentos angulares dados $j_1$ y $j_2$, los valores permitidos para el momento angular total $j$ recorren el intervalo dado por la \textbf{desigualdad triangular cuántica}:
\begin{equation}
    \boxed{|j_1 - j_2| \le j \le j_1 + j_2}, \quad \text{con paso entero } \Delta j = 1.
\end{equation}
El espacio producto total se descompone en suma directa de subespacios irreductibles:
\begin{equation}
    \mathcal{H}_{j_1} \otimes \mathcal{H}_{j_2} = \bigoplus_{j=|j_1-j_2|}^{j_1+j_2} \mathcal{H}_j.
\end{equation}
Comprobación de la conservación de la dimensión total:
\begin{equation}
    \sum_{j=|j_1-j_2|}^{j_1+j_2} (2j + 1) = (2j_1 + 1)(2j_2 + 1).
\end{equation}
El cambio de base unitario entre la base desacoplada y la acoplada se realiza mediante los \textbf{Coeficientes de Clebsch--Gordan} $C_{m_1, m_2; M}^{j_1, j_2; j} \equiv \braket{j_1, m_1; j_2, m_2}{j, M}$:
\begin{equation}
    \boxed{\ket{j, M} = \sum_{m_1=-j_1}^{j_1} \sum_{m_2=-j_2}^{j_2} \braket{j_1, m_1; j_2, m_2}{j, M} \ket{j_1, m_1; j_2, m_2}}.
\end{equation}
Un coeficiente de Clebsch--Gordan es estrictamente nulo a menos que satisfaga la regla de conservación de la proyección magnética:
\begin{equation}
    \braket{j_1, m_1; j_2, m_2}{j, M} = 0 \quad \text{si } M \neq m_1 + m_2.
\end{equation}

\subsubsection{Ejemplo Canónico 1: Acoplamiento de Dos Espines \texorpdfstring{$s_1 = 1/2$}{s1=1/2} y \texorpdfstring{$s_2 = 1/2$}{s2=1/2}}

La base desacoplada $\{\ket{++}, \ket{+-}, \ket{-+}, \ket{--}\}$ tiene dimensión $2\times 2 = 4$. Los valores posibles de $S$ son:
\begin{equation}
    S \in \{|1/2 - 1/2|, \, 1/2 + 1/2\} = \{0, 1\}.
\end{equation}
\begin{itemize}
    \item \textbf{Estado Triplete ($S = 1$, multiplicidad $2S+1 = 3$):}
    \begin{align}
        \ket{S=1, M=1} &= \ket{++}, \\
        \ket{S=1, M=0} &= \frac{1}{\sqrt{2}}\left( \ket{+-} + \ket{-+} \right), \\
        \ket{S=1, M=-1} &= \ket{--},
    \end{align}
    donde el estado con $M = 0$ se obtiene aplicando el operador de descenso $\hat{S}_- = \hat{S}_{1-} + \hat{S}_{2-}$ sobre el estado de máxima proyección $\ket{1, 1}$.
    \item \textbf{Estado Singlete ($S = 0$, multiplicidad $2S+1 = 1$):}
    Por ortogonalidad con el estado triplete $\ket{1, 0}$ en el mismo subespacio $M = 0$:
    \begin{equation}
        \ket{S=0, M=0} = \frac{1}{\sqrt{2}}\left( \ket{+-} - \ket{-+} \right).
    \end{equation}
\end{itemize}
¡El estado singlete es estrictamente antisimétrico ante el intercambio de las dos partículas y representa el célebre estado de entrelazamiento cuántico máximo (estado de Bell)!

\subsubsection{Ejemplo Canónico 2: Acoplamiento \texorpdfstring{$\ell_1 = 1 \otimes \ell_2 = 1$}{l1=1 x l2=1} y Matriz del Bloque \texorpdfstring{$M = 0$}{M=0}}

Consideremos dos momentos angulares de valor $\ell_1 = 1$ y $\ell_2 = 1$. La dimensión total del espacio es $(2\cdot 1 + 1)(2\cdot 1 + 1) = 3 \times 3 = 9$. Los multipletes totales posibles son:
\begin{equation}
    j \in \{0, 1, 2\}, \qquad \dim = 5 + 3 + 1 = 9.
\end{equation}
Para verificar algebraicamente que el operador Casimir $\hat{J}^2$ no es diagonal en la base desacoplada $\{\ket{m_1, m_2}\}$, utilicemos su expresión operativa \eqref{eq:j2_expansion_coupled}:
\begin{equation}
    \hat{J}^2 = \hat{L}_1^2 + \hat{L}_2^2 + 2\hat{L}_{1z}\hat{L}_{2z} + \hat{L}_{1+}\hat{L}_{2-} + \hat{L}_{1-}\hat{L}_{2+}.
\end{equation}
Evaluemos el bloque con proyección magnética nula $M = m_1 + m_2 = 0$. En la base desacoplada ordenada $\{\ket{1, -1}, \ket{0, 0}, \ket{-1, 1}\}$:
\begin{align}
    \hat{J}^2\ket{1, -1} &= \left(2\hbar^2 + 2\hbar^2 + 2\hbar^2(1)(-1)\right)\ket{1, -1} + \hat{L}_{1-}\hat{L}_{2+}\ket{1, -1} \nonumber\\
    &= 2\hbar^2\ket{1, -1} + (\hat{L}_{1-}\ket{1}_1) \otimes (\hat{L}_{2+}\ket{-1}_2) \nonumber\\
    &= 2\hbar^2\ket{1, -1} + (\hbar\sqrt{2}\ket{0}_1) \otimes (\hbar\sqrt{2}\ket{0}_2) = 2\hbar^2\ket{1, -1} + 2\hbar^2\ket{0, 0}.
\end{align}
¡Aparece un elemento fuera de la diagonal proporcional a $2\hbar^2$ que acopla $\ket{1, -1}$ con $\ket{0, 0}$!\\
Efectuando el mismo cálculo sobre los otros dos vectores, la matriz representante de $\hat{J}^2$ en el subespacio $M = 0$ resulta ser exactamente:
\begin{equation}
    \hat{J}^2|_{M=0} = \hbar^2 \begin{pmatrix} 2 & 2 & 0 \\ 2 & 4 & 2 \\ 0 & 2 & 2 \end{pmatrix}.
    \label{eq:matrix_j2_block_zero}
\end{equation}
Calculando los autovalores de esta matriz de $3\times 3$:
\begin{equation}
    \det\left( \hat{J}^2|_{M=0} - \lambda I \right) = -\lambda(\lambda - 2\hbar^2)(\lambda - 6\hbar^2) = 0 \implies \lambda \in \{0, \, 2\hbar^2, \, 6\hbar^2\}.
\end{equation}
Reconocemos de inmediato los valores de $\hbar^2 j(j+1)$:
\begin{itemize}
    \item Para $j = 0$: $\hbar^2 0(1) = 0$.
    \item Para $j = 1$: $\hbar^2 1(2) = 2\hbar^2$.
    \item Para $j = 2$: $\hbar^2 2(3) = 6\hbar^2$.
\end{itemize}
La traza de la matriz es $2 + 4 + 2 = 8 = 0 + 2 + 6$, conservada de forma idéntica. ¡En la base acoplada $\{\ket{j, M}\}$, la matriz $\hat{J}^2$ es estrictamente diagonal!

\begin{figure}[H]
\centering
\begin{tikzpicture}[x=2.2cm,y=1.6cm,>=Stealth,font=\small,
    ket/.style={rounded corners=3pt,draw=grisUNI!70,fill=white,
                minimum width=1.25cm,minimum height=0.7cm,inner sep=2pt},
    cero/.style={ket,draw=azulUNI,fill=azulUNI!15,very thick}]

  % rótulos de ejes
  \node[font=\footnotesize] at (0,-1.95) {$m_1\ (\text{partícula }1)$};
  \node[font=\footnotesize,rotate=90] at (-1.75,0) {$m_2\ (\text{partícula }2)$};

  % diagonales M = m1 + m2 = const
  \draw[dashed,grisUNI!70] (0.65,1.35)--(1.35,0.65)   node[right,font=\footnotesize]{$M=2$};
  \draw[dashed,grisUNI!70] (-0.35,1.35)--(1.35,-0.35) node[right,font=\footnotesize]{$M=1$};
  \draw[dashed,azulUNI]    (-1.35,1.35)--(1.35,-1.35) node[right,font=\footnotesize,azulUNI]{$M=0$};
  \draw[dashed,grisUNI!70] (-1.35,0.35)--(0.35,-1.35) node[right,font=\footnotesize]{$M=-1$};
  \draw[dashed,grisUNI!70] (-1.35,-0.65)--(-0.65,-1.35) node[right,font=\footnotesize]{$M=-2$};

  % estados |m1,m2>
  \foreach \a in {-1,0,1}{
    \foreach \b in {-1,0,1}{
      \pgfmathtruncatemacro{\s}{\a+\b}
      \ifnum\s=0
        \node[cero] at (\a,\b) {$|{\a},{\b}\rangle$};
      \else
        \node[ket]  at (\a,\b) {$|{\a},{\b}\rangle$};
      \fi
    }
  }

  % acoplamientos de J^2 dentro de M = 0
  \draw[<->,very thick,rojoUNI] (0.77,-0.77)--(0.23,-0.23)
        node[midway,above right,font=\footnotesize]{$2\hbar^2$};
  \draw[<->,very thick,rojoUNI] (-0.23,0.23)--(-0.77,0.77)
        node[midway,above right,font=\footnotesize]{$2\hbar^2$};

  % matriz del bloque M = 0
  \node[anchor=west,align=left] at (canvas cs:x=4.2cm,y=0.0cm) {%
    $\hat J^{2}\big|_{M=0}=\hbar^2
     \begin{pmatrix}2&2&0\\2&4&2\\0&2&2\end{pmatrix}$\\[8pt]
    \footnotesize autovalores $\{0,\,2\hbar^2,\,6\hbar^2\}$\\
    \footnotesize $\Rightarrow j=0,1,2$\\[6pt]
    \footnotesize $\dim=5+3+1=9$};
\end{tikzpicture}
\caption{Acoplamiento $\ell_1=1\otimes\ell_2=1$ en la base desacoplada
$|m_1,m_2\rangle$. $\hat J^2$ conserva $M=m_1+m_2$, por lo que sólo mezcla
estados sobre la misma diagonal discontinua. En el bloque $M=0$ (azul) los
términos $\hat L_{1\pm}\hat L_{2\mp}$ generan los elementos fuera de la
diagonal $2\hbar^2$ de la ec.~\eqref{eq:matrix_j2_block_zero}.}
\label{fig:cuadricula-m1m2}
\end{figure}


%% ============================================================
%%  SECCIÓN 6 — Banco de Problemas Resueltos y Ejercicios Propuestos
%% ============================================================

\section{Banco de Problemas Resueltos y Ejercicios Propuestos}
\label{sec:ejercicios}

\subsection{Problemas Resueltos Paso a Paso}

\subsubsection[Problema 1: Rotaciones finitas vía Hadamard]{Problema Resuelto 1: Rotaciones Espaciales Cuánticas Finitas vía Hadamard}
\begin{ejemplobox}[Problema Resuelto 1: Rotaciones Espaciales Cuánticas Finitas vía Hadamard]
    \textbf{Enunciado:} El operador de rotación finita por un ángulo $\theta \in \mathbb{R}$ alrededor del eje polar $z$ está dado por:
    \begin{equation}
        \hat{R}_z(\theta) \equiv \exp\left(-\frac{\iu}{\hbar}\theta\hat{L}_z\right).
    \end{equation}
    \begin{enumerate}
        \item Calcule los conmutadores fundamentales $[\hat{L}_z, \hat{x}]$ y $[\hat{L}_z, \hat{y}]$.
        \item Utilice el Lema de Hadamard para deducir la forma del operador transformado $\hat{x}' \equiv \hat{R}_z\adj(\theta)\hat{x}\hat{R}_z(\theta)$ y verifique que reproduce la matriz de rotación geométrica estándar en el plano $xy$:
        \begin{equation}
            \hat{x}' = \hat{x}\cos\theta - \hat{y}\sin\theta.
        \end{equation}
    \end{enumerate}

    \textbf{Solución Paso a Paso:}

    \textit{Paso 1: Cálculo de Conmutadores de Posición con $\hat{L}_z$:}
    Escribimos $\hat{L}_z = \hat{x}\hat{p}_y - \hat{y}\hat{p}_x$. Evaluamos el conmutador con la coordenada $\hat{x}$:
    \begin{align}
        [\hat{L}_z, \hat{x}] &= [\hat{x}\hat{p}_y - \hat{y}\hat{p}_x, \, \hat{x}] = [\hat{x}\hat{p}_y, \hat{x}] - [\hat{y}\hat{p}_x, \hat{x}].
    \end{align}
    Como $\hat{x}$ conmuta con $\hat{p}_y$ y consigo mismo, $[\hat{x}\hat{p}_y, \hat{x}] = 0$. Por ende:
    \begin{equation}
        [\hat{L}_z, \hat{x}] = -[\hat{y}\hat{p}_x, \hat{x}] = -\hat{y}[\hat{p}_x, \hat{x}] - [\hat{y}, \hat{x}]\hat{p}_x.
    \end{equation}
    Sabiendo que $[\hat{p}_x, \hat{x}] = -\iu\hbar\hat{I}$ y $[\hat{y}, \hat{x}] = 0$:
    \begin{equation}
        [\hat{L}_z, \hat{x}] = -\hat{y}(-\iu\hbar\hat{I}) = \boxed{\iu\hbar\hat{y}}.
        \label{eq:lz_x_comm}
    \end{equation}
    Análogamente, para la coordenada $\hat{y}$:
    \begin{equation}
        [\hat{L}_z, \hat{y}] = [\hat{x}\hat{p}_y - \hat{y}\hat{p}_x, \, \hat{y}] = [\hat{x}\hat{p}_y, \hat{y}] = \hat{x}[\hat{p}_y, \hat{y}] = \hat{x}(-\iu\hbar\hat{I}) = \boxed{-\iu\hbar\hat{x}}.
        \label{eq:lz_y_comm}
    \end{equation}

    \textit{Paso 2: Expansión de Hadamard:}
    Como $\hat{L}_z$ es hermítico, el operador conjugante es $\hat{R}_z\adj(\theta) = \exp\left(\frac{\iu}{\hbar}\theta\hat{L}_z\right)$. Aplicamos el Lema de Hadamard definiendo el operador $\hat{A} \equiv \frac{\iu\theta}{\hbar}\hat{L}_z$ y $\hat{B} \equiv \hat{x}$:
    \begin{equation}
        \hat{x}' = \eu^{\hat{A}}\hat{x}\eu^{-\hat{A}} = \hat{x} + [\hat{A}, \hat{x}] + \frac{1}{2!}[\hat{A}, [\hat{A}, \hat{x}]] + \frac{1}{3!}[\hat{A}, [\hat{A}, [\hat{A}, \hat{x}]]] + \dots
    \end{equation}
    Evaluemos término a término los conmutadores anidados:
    \begin{itemize}
        \item Primer orden ($n = 1$):
        \begin{equation}
            [\hat{A}, \hat{x}] = \left[\frac{\iu\theta}{\hbar}\hat{L}_z, \, \hat{x}\right] = \frac{\iu\theta}{\hbar}[\hat{L}_z, \hat{x}] = \frac{\iu\theta}{\hbar}(\iu\hbar\hat{y}) = \iu^2 \theta \hat{y} = -\theta\hat{y}.
        \end{equation}
        \item Segundo orden ($n = 2$):
        \begin{equation}
            [\hat{A}, [\hat{A}, \hat{x}]] = [\hat{A}, -\theta\hat{y}] = -\theta\left[\frac{\iu\theta}{\hbar}\hat{L}_z, \, \hat{y}\right] = -\frac{\iu\theta^2}{\hbar}(-\iu\hbar\hat{x}) = \iu^2 \theta^2 \hat{x} = -\theta^2 \hat{x}.
        \end{equation}
        \item Tercer orden ($n = 3$):
        \begin{equation}
            [\hat{A}, [\hat{A}, [\hat{A}, \hat{x}]]] = [\hat{A}, -\theta^2\hat{x}] = -\theta^2 [\hat{A}, \hat{x}] = -\theta^2(-\theta\hat{y}) = +\theta^3\hat{y}.
        \end{equation}
        \item Cuarto orden ($n = 4$):
        \begin{equation}
            [\hat{A}, [\hat{A}, [\hat{A}, [\hat{A}, \hat{x}]]]] = [\hat{A}, \theta^3\hat{y}] = \theta^3(+\theta\hat{x}) = +\theta^4\hat{x}.
        \end{equation}
    \end{itemize}
    Sustituyendo estos resultados en la serie infinita y agrupando por las componentes de los operadores $\hat{x}$ y $\hat{y}$:
    \begin{align}
        \hat{x}' &= \hat{x}\left(1 - \frac{\theta^2}{2!} + \frac{\theta^4}{4!} - \dots\right) - \hat{y}\left(\theta - \frac{\theta^3}{3!} + \frac{\theta^5}{5!} - \dots\right).
    \end{align}
    Reconocemos de inmediato las expansiones de Taylor clásicas de las funciones coseno y seno:
    \begin{equation}
        \boxed{\hat{x}' = \hat{x}\cos\theta - \hat{y}\sin\theta}.
    \end{equation}
    De manera idéntica se demuestra que $\hat{y}' = \hat{x}\sin\theta + \hat{y}\cos\theta$. ¡El Lema de Hadamard genera la rotación geométrica exacta en el plano euclídeo a partir del álgebra pura de conmutadores!
\end{ejemplobox}

\begin{figure}[H]
\centering
\begin{subfigure}[t]{0.45\textwidth}
\centering
\begin{tikzpicture}[scale=2.2]
  \draw[->,gray] (-0.3,0) -- (1.35,0) node[right] {$\hat x$};
  \draw[->,gray] (0,-0.9) -- (0,1.2) node[above] {$\hat y$};
  \draw[dashed,gray!60] (0,0) circle (1);
  \draw[-{Stealth},very thick,uniblue] (0,0) -- (1,0);
  \draw[-{Stealth},very thick,unired] (0,0) -- ({cos(40)},{-sin(40)})
        node[below right] {$\hat x'$};
  \draw[unigreen,thick] (0.45,0) arc (0:-40:0.45);
  \node[unigreen] at (0.62,-0.13) {$\theta$};
  \draw[dotted,unired] ({cos(40)},{-sin(40)}) -- ({cos(40)},0);
  \draw[dotted,unired] ({cos(40)},{-sin(40)}) -- (0,{-sin(40)});
  \node[font=\scriptsize,unired,above] at ({cos(40)},0) {$\cos\theta$};
  \node[font=\scriptsize,unired,left] at (0,{-sin(40)}) {$-\sin\theta$};
\end{tikzpicture}
\caption{$\hat x'=\hat x\cos\theta-\hat y\sin\theta$ (componentes sobre $\hat x,\hat y$).}
\end{subfigure}\hfill
\begin{subfigure}[t]{0.52\textwidth}
\centering
\begin{tikzpicture}
\begin{axis}[width=\linewidth,height=6.2cm,xlabel={$\theta$},
  xmin=0,xmax=7,ymin=-2.2,ymax=2.2,
  legend style={font=\scriptsize,at={(0.5,-0.28)},anchor=north,legend columns=3},
  grid=both,minor tick num=1,
  xtick={0,1.5708,3.1416,4.7124,6.2832},
  xticklabels={$0$,$\pi/2$,$\pi$,$3\pi/2$,$2\pi$}]
  \addplot[black,very thick,domain=0:7,samples=200] {cos(deg(x))};
  \addlegendentry{$\cos\theta$}
  \addplot[uniblue,dashed,thick,domain=0:7,samples=100] {1};
  \addlegendentry{orden 0}
  \addplot[unired,thick,domain=0:7,samples=100] {1-x^2/2};
  \addlegendentry{hasta orden 2}
  \addplot[unigreen,thick,domain=0:7,samples=100] {1-x^2/2+x^4/24};
  \addlegendentry{hasta orden 4}
  \addplot[orange,thick,domain=0:7,samples=100] {1-x^2/2+x^4/24-x^6/720};
  \addlegendentry{hasta orden 6}
\end{axis}
\end{tikzpicture}
\caption{Sumas parciales de la serie de Hadamard (coeficiente de $\hat x$).}
\end{subfigure}
\caption{Rotaci\'on finita generada por $\hat L_z$: cada orden de la serie de Hadamard a\~nade un t\'ermino de la serie de Taylor de $\cos\theta$ y $\sin\theta$.}
\label{fig:rot_hadamard}
\end{figure}

\subsubsection[Problema 2: Operador vectorial y tensores]{Problema Resuelto 2: El Operador Vectorial Cuántico y Tensores Cartesianos}
\begin{ejemplobox}[Problema Resuelto 2: El Operador Vectorial Cuántico]
    \textbf{Enunciado:} En mecánica cuántica, se define formalmente como \textbf{operador vectorial} a cualquier triplete de observables $\hat{\vec{V}} = (\hat{V}_1, \hat{V}_2, \hat{V}_3)$ que bajo una rotación espacial se transforme como un vector geométrico, lo cual equivale a la relación de conmutación:
    \begin{equation}
        [\hat{L}_i, \hat{V}_j] = \iu\hbar\,\varepsilon_{ijk}\hat{V}_k.
    \end{equation}
    Demuestre que el vector de posición $\hat{\vec{r}} = (\hat{x}_1, \hat{x}_2, \hat{x}_3)$ y el vector de momento lineal $\hat{\vec{p}} = (\hat{p}_1, \hat{p}_2, \hat{p}_3)$ son operadores vectoriales genuinos empleando cálculo tensorial.

    \textbf{Solución Paso a Paso:}
    
    Escribimos la componente $\hat{L}_i$ en notación tensorial: $\hat{L}_i = \varepsilon_{iab}\hat{x}_a\hat{p}_b$. Evaluamos el conmutador con la coordenada $\hat{x}_j$:
    \begin{align}
        [\hat{L}_i, \hat{x}_j] &= [\varepsilon_{iab}\hat{x}_a\hat{p}_b, \, \hat{x}_j] = \varepsilon_{iab}[\hat{x}_a\hat{p}_b, \, \hat{x}_j].
    \end{align}
    Aplicando la regla del conmutador de un producto $[\hat{A}\hat{B}, \hat{C}] = \hat{A}[\hat{B}, \hat{C}] + [\hat{A}, \hat{C}]\hat{B}$:
    \begin{equation}
        [\hat{x}_a\hat{p}_b, \, \hat{x}_j] = \hat{x}_a[\hat{p}_b, \hat{x}_j] + [\hat{x}_a, \hat{x}_j]\hat{p}_b.
    \end{equation}
    Dado que $[\hat{x}_a, \hat{x}_j] = 0$ y $[\hat{p}_b, \hat{x}_j] = -\iu\hbar\delta_{bj}$:
    \begin{equation}
        [\hat{x}_a\hat{p}_b, \, \hat{x}_j] = \hat{x}_a(-\iu\hbar\delta_{bj}) = -\iu\hbar\delta_{bj}\hat{x}_a.
    \end{equation}
    Sustituyendo esto dentro de la suma tensorial:
    \begin{equation}
        [\hat{L}_i, \hat{x}_j] = \varepsilon_{iab}(-\iu\hbar\delta_{bj}\hat{x}_a) = -\iu\hbar\,\varepsilon_{iaj}\hat{x}_a.
    \end{equation}
    Intercambiamos los dos últimos índices del tensor de Levi-Civita para invertir el signo: $\varepsilon_{iaj} = -\varepsilon_{ija}$:
    \begin{equation}
        [\hat{L}_i, \hat{x}_j] = \iu\hbar\,\varepsilon_{ija}\hat{x}_a = \boxed{\iu\hbar\,\varepsilon_{ijk}\hat{x}_k}.
    \end{equation}
    Para el momento $\hat{p}_j$, evaluamos de la misma manera:
    \begin{equation}
        [\hat{L}_i, \hat{p}_j] = \varepsilon_{iab}[\hat{x}_a\hat{p}_b, \, \hat{p}_j] = \varepsilon_{iab}[\hat{x}_a, \hat{p}_j]\hat{p}_b = \varepsilon_{iab}(\iu\hbar\delta_{aj})\hat{p}_b = \iu\hbar\,\varepsilon_{ijb}\hat{p}_b = \boxed{\iu\hbar\,\varepsilon_{ijk}\hat{p}_k}.
    \end{equation}
    Ambos operadores satisfacen la definición exacta de operador vectorial cuántico.
\end{ejemplobox}

\subsubsection[Problema 3: Incertidumbre de momento angular en $\ket{\ell, m}$]{Problema Resuelto 3: Incertidumbre de Momento Angular en $\ket{\ell, m}$ y Operadores Escalera}
\begin{ejemplobox}[Problema Resuelto 3: Incertidumbre de Momento Angular en el Autoestado $\ket{\ell, m}$]
    \textbf{Enunciado:} Considere una partícula en un autoestado común $\ket{\ell, m}$ de $\hat{L}^2$ y $\hat{L}_z$.
    \begin{enumerate}
        \item Calcule los valores esperados $\ev{\hat{L}_x}$ y $\ev{\hat{L}_y}$.
        \item Calcule $\ev{\hat{L}_x^2}$ y $\ev{\hat{L}_y^2}$ utilizando los operadores escalera $\hat{L}_\pm = \hat{L}_x \pm \iu\hat{L}_y$.
        \item Verifique que el producto de incertidumbres $\Delta L_x \Delta L_y$ cumple estrictamente la relación de Heisenberg.
    \end{enumerate}

    \textbf{Solución Paso a Paso:}

    \textit{Paso 1: Operadores Escalera y Valores Esperados:}
    Los operadores escalera se definen como $\hat{L}_+ = \hat{L}_x + \iu\hat{L}_y$ y $\hat{L}_- = \hat{L}_x - \iu\hat{L}_y$. Despejando las componentes cartesianas:
    \begin{equation}
        \hat{L}_x = \frac{1}{2}(\hat{L}_+ + \hat{L}_-), \qquad \hat{L}_y = \frac{1}{2\iu}(\hat{L}_+ - \hat{L}_-).
    \end{equation}
    La acción de los operadores escalera sobre la base ortonormal es:
    \begin{equation}
        \hat{L}_\pm \ket{\ell, m} = \hbar\sqrt{\ell(\ell+1) - m(m\pm 1)}\ket{\ell, m\pm 1}.
    \end{equation}
    Al tomar el producto interno con $\bra{\ell, m}$, como $\braket{\ell, m}{\ell, m\pm 1} = 0$ por ortogonalidad:
    \begin{equation}
        \ev{\hat{L}_+} = \matrixel{\ell, m}{\hat{L}_+}{\ell, m} = 0, \qquad \ev{\hat{L}_-} = \matrixel{\ell, m}{\hat{L}_-}{\ell, m} = 0.
    \end{equation}
    Por lo tanto:
    \begin{equation}
        \ev{\hat{L}_x} = 0, \qquad \ev{\hat{L}_y} = 0.
    \end{equation}

    \textit{Paso 2: Varianza Cuántica y Demostración Algebraica de $\ev{\hat{L}_x^2} = \ev{\hat{L}_y^2}$:}
    Para evaluar rigurosamente $\ev{\hat{L}_x^2}$ y $\ev{\hat{L}_y^2}$ sin recurrir a argumentos heurísticos de simetría, expresemos los cuadrados directamente en términos de los operadores escalera:
    \begin{align}
        \hat{L}_x^2 &= \frac{1}{4}\left(\hat{L}_+ + \hat{L}_-\right)^2 = \frac{1}{4}\left( \hat{L}_+^2 + \hat{L}_-^2 + \hat{L}_+\hat{L}_- + \hat{L}_-\hat{L}_+ \right), \\
        \hat{L}_y^2 &= -\frac{1}{4}\left(\hat{L}_+ - \hat{L}_-\right)^2 = -\frac{1}{4}\left( \hat{L}_+^2 + \hat{L}_-^2 - \hat{L}_+\hat{L}_- - \hat{L}_-\hat{L}_+ \right).
    \end{align}
    Al tomar el valor esperado en el estado $\ket{\ell, m}$, los términos $\hat{L}_\pm^2$ desplazan el número cuántico magnético en dos unidades ($\hat{L}_\pm^2\ket{\ell, m} \propto \ket{\ell, m \pm 2}$). Como la base es ortogonal ($\braket{\ell, m}{\ell, m \pm 2} = 0$), sus valores esperados se anulan idénticamente:
    \begin{equation}
        \ev{\hat{L}_+^2} = \matrixel{\ell, m}{\hat{L}_+^2}{\ell, m} = 0, \qquad \ev{\hat{L}_-^2} = \matrixel{\ell, m}{\hat{L}_-^2}{\ell, m} = 0.
    \end{equation}
    En consecuencia, ambos valores esperados se reducen de forma exacta a la misma expresión simétrica:
    \begin{equation}
        \ev{\hat{L}_x^2} = \ev{\hat{L}_y^2} = \frac{1}{4}\ev{\hat{L}_+\hat{L}_- + \hat{L}_-\hat{L}_+}.
    \end{equation}
    Por otro lado, la suma de las tres componentes al cuadrado es el operador Casimir $\hat{L}^2$:
    \begin{equation}
        \hat{L}^2 = \hat{L}_x^2 + \hat{L}_y^2 + \hat{L}_z^2 \implies \ev{\hat{L}_x^2} + \ev{\hat{L}_y^2} = \ev{\hat{L}^2} - \ev{\hat{L}_z^2} = \hbar^2[\ell(\ell+1) - m^2].
    \end{equation}
    Como ambos valores esperados son iguales, se deduce algebraicamente:
    \begin{equation}
        \ev{\hat{L}_x^2} = \ev{\hat{L}_y^2} = \frac{\hbar^2}{2}[\ell(\ell+1) - m^2].
    \end{equation}
    Dado que las medias son nulas ($\ev{\hat{L}_x} = \ev{\hat{L}_y} = 0$), las desviaciones estándar son:
    \begin{equation}
        \Delta L_x = \Delta L_y = \hbar\sqrt{\frac{\ell(\ell+1) - m^2}{2}}.
    \end{equation}

    \textit{Paso 3: Verificación de la Relación de Incertidumbre:}
    El producto de incertidumbres resulta ser:
    \begin{equation}
        \Delta L_x \Delta L_y = \frac{\hbar^2}{2}[\ell(\ell+1) - m^2].
    \end{equation}
    Por otro lado, la cota inferior dictada por la relación de Heisenberg para $[\hat{L}_x, \hat{L}_y] = \iu\hbar\hat{L}_z$ es:
    \begin{equation}
        \frac{1}{2}|\ev{[\hat{L}_x, \hat{L}_y]}| = \frac{\hbar}{2}|\ev{\hat{L}_z}| = \frac{\hbar^2}{2}|m|.
    \end{equation}
    Comprobemos que este valor estadístico satisface rigurosamente el principio de incertidumbre de Heisenberg generalizado:
    \begin{equation}
        \frac{\hbar^2}{2}[\ell(\ell+1) - m^2] \ge \frac{\hbar^2}{2}|m|.
    \end{equation}
    Esto equivale algebraicamente a:
    \begin{equation}
        \ell(\ell+1) - m^2 \ge |m| \iff \ell(\ell+1) \ge m^2 + |m| = |m|(|m|+1).
    \end{equation}
    Dado que por la teoría general de momentos angulares $|m| \le \ell$, la función $f(x) = x(x+1)$ es monótonamente creciente para $x \ge 0$. Por consiguiente, $\ell(\ell+1) \ge |m|(|m|+1)$ se satisface para todo $m \in \{-\ell, \dots, \ell\}$. La igualdad se satura si y solo si $|m| = \ell$ (estados de máximo alineamiento angular).
\end{ejemplobox}

\begin{figure}[H]
    \centering
    % Fig. 4: producto de incertidumbres vs cota, \ell = 3 (unidades hbar^2)
\begin{tikzpicture}
\begin{axis}[width=0.8\linewidth, height=6cm, domain=-3:3, samples=100,
  xlabel={$m$}, ylabel={$\Delta L_x\,\Delta L_y\;/\;\hbar^2$},
  axis lines=left, xtick={-3,-2,-1,0,1,2,3}, ymin=0, ymax=9.5,
  legend style={font=\scriptsize, draw=none, fill=white, at={(0.5,0.97)}, anchor=north},
  every axis plot/.append style={thick}]
  \addplot[uni-blue]{0.5*(12-x^2)};
  \addplot[alert-border]{0.5*abs(x)};
  \addplot[only marks, mark=*, uni-blue, mark size=2pt, samples at={-3,-2,...,3}]{0.5*(12-x^2)};
  \addplot[only marks, mark=square*, example-border, mark size=2.5pt, samples at={-3,3}]{0.5*abs(x)};
  \legend{$\tfrac12[\ell(\ell+1)-m^2]$, cota $\tfrac12|m|$, valores permitidos, saturaci\'on ($|m|=\ell$)}
\end{axis}
\end{tikzpicture}

    \caption{Producto $\Delta L_x\,\Delta L_y$ en $\ket{\ell,m}$ frente a la cota de Heisenberg, para $\ell=3$. La igualdad solo se alcanza en $|m|=\ell$.}
    \label{fig:incert_Lxy}
\end{figure}

\subsubsection[Problema 4: Efecto Zeeman y precesión de Larmor]{Problema Resuelto 4: Efecto Zeeman Orbital, Ecuaciones de Heisenberg y Precesión de Larmor}
\begin{ejemplobox}[Problema Resuelto 4: Efecto Zeeman Orbital y Precesión de Larmor Cuántica]
    \textbf{Enunciado:} Un electrón de masa $m_e$ y carga $-e$ en un átomo se encuentra inmerso en un campo magnético estático y homogéneo $\vec{B} = B_0\hat{z}$. El momento dipolar magnético orbital está acoplado al momento angular orbital mediante $\hat{\vec{\mu}}_L = -\frac{e}{2m_e}\hat{\vec{L}}$. El Hamiltoniano total del sistema (efecto Zeeman normal) viene dado por:
    \begin{equation}
        \hat{H} = \hat{H}_0 - \hat{\vec{\mu}}_L \cdot \vec{B} = \hat{H}_0 + \omega_L \hat{L}_z,
    \end{equation}
    donde $\omega_L \equiv \frac{e B_0}{2m_e}$ es la frecuencia clásica de Larmor y $[\hat{H}_0, \hat{\vec{L}}] = 0$.
    \begin{enumerate}
        \item Calcule los conmutadores fundamentales $[\hat{H}, \hat{L}_x]$, $[\hat{H}, \hat{L}_y]$ y $[\hat{H}, \hat{L}_z]$.
        \item En la representación de Heisenberg, la evolución temporal de cualquier observable satisface $\frac{\dif\hat{A}_H}{\dif t} = \frac{1}{\iu\hbar}[\hat{A}_H, \hat{H}]$. Deduzca el sistema de ecuaciones diferenciales acopladas para $\frac{\dif\hat{L}_x}{\dif t}$, $\frac{\dif\hat{L}_y}{\dif t}$ y $\frac{\dif\hat{L}_z}{\dif t}$.
        \item Desacople el sistema definiendo los operadores escalera $\hat{L}_\pm = \hat{L}_x \pm \iu\hat{L}_y$ y obtenga las soluciones analíticas para $\hat{L}_x(t)$ y $\hat{L}_y(t)$ en función de sus valores iniciales en $t = 0$.
        \item Demuestre el Teorema de Ehrenfest para este sistema y describa la trayectoria cuántica de los valores esperados $\ev{\hat{\vec{L}}(t)}$.
    \end{enumerate}

    \textbf{Solución Paso a Paso:}

    \textit{Paso 1: Conmutadores con el Hamiltoniano:}
    Puesto que $[\hat{H}_0, \hat{\vec{L}}] = 0$, los conmutadores se reducen al acoplamiento con $\hat{L}_z$:
    \begin{align}
        [\hat{H}, \hat{L}_x] &= [\hat{H}_0 + \omega_L\hat{L}_z, \, \hat{L}_x] = \omega_L [\hat{L}_z, \hat{L}_x] = \omega_L(\iu\hbar\hat{L}_y) = \iu\hbar\omega_L\hat{L}_y, \\
        [\hat{H}, \hat{L}_y] &= [\hat{H}_0 + \omega_L\hat{L}_z, \, \hat{L}_y] = \omega_L [\hat{L}_z, \hat{L}_y] = \omega_L(-\iu\hbar\hat{L}_x) = -\iu\hbar\omega_L\hat{L}_x, \\
        [\hat{H}, \hat{L}_z] &= \omega_L [\hat{L}_z, \hat{L}_z] = 0.
    \end{align}

    \textit{Paso 2: Ecuaciones de Movimiento de Heisenberg:}
    Sustituyendo los conmutadores en la ecuación de Heisenberg (notando que $[\hat{L}_i, \hat{H}] = -[\hat{H}, \hat{L}_i]$):
    \begin{align}
        \frac{\dif\hat{L}_x}{\dif t} &= \frac{1}{\iu\hbar}[\hat{L}_x, \hat{H}] = -\frac{1}{\iu\hbar}[\hat{H}, \hat{L}_x] = -\frac{1}{\iu\hbar}(\iu\hbar\omega_L\hat{L}_y) = -\omega_L\hat{L}_y, \\
        \frac{\dif\hat{L}_y}{\dif t} &= \frac{1}{\iu\hbar}[\hat{L}_y, \hat{H}] = -\frac{1}{\iu\hbar}[\hat{H}, \hat{L}_y] = -\frac{1}{\iu\hbar}(-\iu\hbar\omega_L\hat{L}_x) = +\omega_L\hat{L}_x, \\
        \frac{\dif\hat{L}_z}{\dif t} &= \frac{1}{\iu\hbar}[\hat{L}_z, \hat{H}] = 0 \implies \hat{L}_z(t) = \hat{L}_z(0) = \text{Constante de movimiento}.
    \end{align}

    \textit{Paso 3: Desacoplamiento mediante Operadores Escalera:}
    Combinando linealmente las dos primeras ecuaciones:
    \begin{equation}
        \frac{\dif\hat{L}_\pm}{\dif t} = \frac{\dif\hat{L}_x}{\dif t} \pm \iu\frac{\dif\hat{L}_y}{\dif t} = -\omega_L\hat{L}_y \pm \iu\omega_L\hat{L}_x = \pm\iu\omega_L(\hat{L}_x \pm \iu\hat{L}_y) = \pm\iu\omega_L\hat{L}_\pm.
    \end{equation}
    ¡Las ecuaciones para $\hat{L}_+$ y $\hat{L}_-$ están completamente desacopladas! Integrando directamente respecto al tiempo:
    \begin{equation}
        \hat{L}_\pm(t) = \hat{L}_\pm(0)\eu^{\pm\iu\omega_L t}.
    \end{equation}
    Recuperamos las componentes cartesianas usando $\hat{L}_x = \frac{1}{2}(\hat{L}_+ + \hat{L}_-)$ y $\hat{L}_y = \frac{1}{2\iu}(\hat{L}_+ - \hat{L}_-)$:
    \begin{align}
        \hat{L}_x(t) &= \frac{1}{2}\left( \hat{L}_+(0)\eu^{\iu\omega_L t} + \hat{L}_-(0)\eu^{-\iu\omega_L t} \right) \nonumber\\
        &= \frac{1}{2}\left( (\hat{L}_x(0) + \iu\hat{L}_y(0))\eu^{\iu\omega_L t} + (\hat{L}_x(0) - \iu\hat{L}_y(0))\eu^{-\iu\omega_L t} \right) \nonumber\\
        &= \hat{L}_x(0)\left(\frac{\eu^{\iu\omega_L t} + \eu^{-\iu\omega_L t}}{2}\right) - \hat{L}_y(0)\left(\frac{\eu^{\iu\omega_L t} - \eu^{-\iu\omega_L t}}{2\iu}\right) \nonumber\\
        &= \boxed{\hat{L}_x(0)\cos(\omega_L t) - \hat{L}_y(0)\sin(\omega_L t)}.
    \end{align}
    Análogamente, para la componente transversal $y$:
    \begin{equation}
        \boxed{\hat{L}_y(t) = \hat{L}_x(0)\sin(\omega_L t) + \hat{L}_y(0)\cos(\omega_L t)}.
    \end{equation}

    \textit{Paso 4: Teorema de Ehrenfest y Precesión de Larmor:}
    Tomando el valor esperado en un estado cuántico arbitrario $\ket{\psi}$:
    \begin{equation}
        \frac{\dif}{\dif t}\ev{\hat{\vec{L}}} = \vec{\omega}_L \times \ev{\hat{\vec{L}}}, \quad \text{con } \vec{\omega}_L = \omega_L\hat{z}.
    \end{equation}
    El vector valor esperado $\ev{\hat{\vec{L}}(t)}$ describe una rotación circular uniforme en el plano $xy$ con velocidad angular constante $\omega_L$, manteniendo constante su ángulo de inclinación respecto al eje $z$. ¡La mecánica cuántica reproduce con total fidelidad la célebre precesión de Larmor del giróscopo magnético clásico!
\end{ejemplobox}

\subsubsection[Problema 5: Oscilador 2D y modos circulares]{Problema Resuelto 5: Oscilador Armónico 2D Isótropo y Modos Circulares \texorpdfstring{$\ket{R}$}{R} y \texorpdfstring{$\ket{L}$}{L}}
\begin{ejemplobox}[Problema Resuelto 5: Oscilador 2D Isótropo y Modos Circulares Cuánticos]
    \textbf{Enunciado:} Considere una partícula de masa $M$ en un potencial armónico bidimensional isótropo:
    \begin{equation}
        \hat{H} = \frac{\hat{p}_x^2 + \hat{p}_y^2}{2M} + \frac{1}{2}M\omega^2(\hat{x}^2 + \hat{y}^2) = \hbar\omega\left( \hat{a}_x\adj\hat{a}_x + \hat{a}_y\adj\hat{a}_y + 1 \right),
    \end{equation}
    y el operador momento angular orbital perpendicular al plano de oscilación:
    \begin{equation}
        \hat{L}_z = \hat{x}\hat{p}_y - \hat{y}\hat{p}_x = -\iu\hbar\left(\hat{a}_x\adj\hat{a}_y - \hat{a}_y\adj\hat{a}_x\right).
    \end{equation}
    \begin{enumerate}
        \item Demuestre que $[\hat{H}, \hat{L}_z] = 0$. ¿Qué garantiza este resultado sobre los autoestados de energía?
        \item Para el primer nivel de energía excitado $E_1 = 2\hbar\omega$, la base cartesiana desacoplada $\{\ket{n_x, n_y}\}$ es bidimensional: $\{\ket{1, 0} = \hat{a}_x\adj\ket{0, 0}, \, \ket{0, 1} = \hat{a}_y\adj\ket{0, 0}\}$. Construya la matriz de $2\times 2$ representante de $\hat{L}_z$ restringida a este subespacio degenerado $E_1$.
        \item Diagonalice la matriz de $\hat{L}_z$ en $E_1$ y halle los autovalores y autoestados simultáneos comunes a $\hat{H}$ y $\hat{L}_z$.
        \item Demuestre que estos autoestados corresponden a modos circulares cuánticos dextrógiro $\ket{R}$ y levógiro $\ket{L}$ con autovalores $L_z = \pm\hbar$.
    \end{enumerate}

    \textbf{Solución Paso a Paso:}

    \textit{Paso 1: Conmutación $[\hat{H}, \hat{L}_z] = 0$:}
    Escribimos $\hat{H} = \hbar\omega(\hat{N} + 1)$, donde $\hat{N} = \hat{a}_x\adj\hat{a}_x + \hat{a}_y\adj\hat{a}_y$. Evaluamos el conmutador con $\hat{L}_z$:
    \begin{align}
        [\hat{N}, \hat{a}_x\adj\hat{a}_y] &= [\hat{a}_x\adj\hat{a}_x, \hat{a}_x\adj\hat{a}_y] + [\hat{a}_y\adj\hat{a}_y, \hat{a}_x\adj\hat{a}_y] \nonumber\\
        &= \hat{a}_x\adj[\hat{a}_x, \hat{a}_x\adj]\hat{a}_y + \hat{a}_x\adj[\hat{a}_y\adj\hat{a}_y, \hat{a}_y] \nonumber\\
        &= \hat{a}_x\adj(1)\hat{a}_y + \hat{a}_x\adj(-\hat{a}_y) = \hat{a}_x\adj\hat{a}_y - \hat{a}_x\adj\hat{a}_y = 0.
    \end{align}
    Análogamente, $[\hat{N}, \hat{a}_y\adj\hat{a}_x] = 0$. Por linealidad:
    \begin{equation}
        [\hat{H}, \hat{L}_z] = -\iu\hbar^2\omega [\hat{N}, \, \hat{a}_x\adj\hat{a}_y - \hat{a}_y\adj\hat{a}_x] = 0.
    \end{equation}
    Por el Teorema 2.3, $\hat{H}$ y $\hat{L}_z$ son compatibles y admiten una base ortonormal común.

    \textit{Paso 2: Matriz Representante en el Subespacio Degenerado $E_1$:}
    La base ordenada de $E_1$ es $\{\ket{1, 0}, \ket{0, 1}\}$. Evaluemos la acción de $\hat{L}_z = -\iu\hbar(\hat{a}_x\adj\hat{a}_y - \hat{a}_y\adj\hat{a}_x)$:
    \begin{align}
        \hat{L}_z\ket{1, 0} &= -\iu\hbar\left( \hat{a}_x\adj\hat{a}_y\ket{1, 0} - \hat{a}_y\adj\hat{a}_x\ket{1, 0} \right) = -\iu\hbar(0 - \hat{a}_y\adj\ket{0, 0}) = +\iu\hbar\ket{0, 1}, \\
        \hat{L}_z\ket{0, 1} &= -\iu\hbar\left( \hat{a}_x\adj\hat{a}_y\ket{0, 1} - \hat{a}_y\adj\hat{a}_x\ket{0, 1} \right) = -\iu\hbar(\hat{a}_x\adj\ket{0, 0} - 0) = -\iu\hbar\ket{1, 0}.
    \end{align}
    Calculamos los elementos de matriz en la base $\{\ket{1, 0}, \ket{0, 1}\}$:
    \begin{equation}
        \hat{L}_z|_{E_1} = \begin{pmatrix}
            \matrixel{1, 0}{\hat{L}_z}{1, 0} & \matrixel{1, 0}{\hat{L}_z}{0, 1} \\
            \matrixel{0, 1}{\hat{L}_z}{1, 0} & \matrixel{0, 1}{\hat{L}_z}{0, 1}
        \end{pmatrix} = \hbar \begin{pmatrix} 0 & -\iu \\ \iu & 0 \end{pmatrix} = \hbar\sigma_y.
    \end{equation}
    ¡La matriz representante en este subespacio es exactamente la matriz de Pauli $\sigma_y$ multiplicada por $\hbar$!

    \textit{Paso 3 y 4: Diagonalización y Modos Circulares $\ket{R}$ y $\ket{L}$:}
    El polinomio característico es:
    \begin{equation}
        \det(\hat{L}_z|_{E_1} - \lambda I) = \begin{vmatrix} -\lambda & -\iu\hbar \\ \iu\hbar & -\lambda \end{vmatrix} = \lambda^2 - \hbar^2 = 0 \implies \lambda = \pm\hbar.
    \end{equation}
    Determinemos los autovectores normalizados para cada autovalor:
    \begin{itemize}
        \item Para $\lambda_1 = +\hbar$ (modo circular dextrógiro $\ket{R}$):
        \begin{equation}
            \begin{pmatrix} -\hbar & -\iu\hbar \\ \iu\hbar & -\hbar \end{pmatrix}\begin{pmatrix} c_1 \\ c_2 \end{pmatrix} = 0 \implies c_1 = -\iu c_2 \implies \boxed{\ket{R} = \frac{1}{\sqrt{2}}\left( \ket{1, 0} + \iu\ket{0, 1} \right)}.
        \end{equation}
        \item Para $\lambda_2 = -\hbar$ (modo circular levógiro $\ket{L}$):
        \begin{equation}
            \begin{pmatrix} \hbar & -\iu\hbar \\ \iu\hbar & \hbar \end{pmatrix}\begin{pmatrix} c_1 \\ c_2 \end{pmatrix} = 0 \implies c_1 = \iu c_2 \implies \boxed{\ket{L} = \frac{1}{\sqrt{2}}\left( \ket{1, 0} - \iu\ket{0, 1} \right)}.
        \end{equation}
    \end{itemize}
    Notemos la física detrás del resultado: el estado $\ket{1, 0}$ representa una oscilación lineal pura a lo largo del eje $x$, mientras que $\ket{0, 1}$ oscila a lo largo de $y$. La combinación $\ket{R} \propto \ket{1, 0} + \iu\ket{0, 1}$ introduce un desfase cuántico de $\pi/2$ ($\eu^{\iu\pi/2} = \iu$), exactamente idéntico a cómo dos ondas clásicas desfasadas en cuadratura generan polarización circular dextrógira y levógira. ¡El Teorema 2.3 (diagonalización simultánea) en bloques degenerados resuelve el problema con elegancia absoluta!
\end{ejemplobox}

\subsection{Ejercicios Propuestos para Autoestudio (Nivel Avanzado / Zettili)}

\subsubsection[Ejercicio 1: Conmutador de la exponencial]{Ejercicio Propuesto 1: Conmutador de la Exponencial de un Operador}
\begin{ejerciciobox}[Ejercicio Propuesto 1: Conmutador de la Exponencial de un Operador]
    \textbf{Nivel: Intermedio}\\
    Sean $\hat{A}$ y $\hat{B}$ dos operadores que satisfacen la condición $[[\hat{A}, \hat{B}], \hat{A}] = 0$.
    \begin{enumerate}
        \item Demuestre que $[\hat{A}^n, \hat{B}] = n\hat{A}^{n-1}[\hat{A}, \hat{B}]$ para todo entero positivo $n \ge 1$.
        \item Utilice la expansión en serie de potencias de la exponencial para demostrar formalmente que:
        \begin{equation}
            [\eu^{\hat{A}}, \hat{B}] = [\hat{A}, \hat{B}]\eu^{\hat{A}}.
        \end{equation}
    \end{enumerate}
    \textit{Pista:} Utilice inducción matemática para la parte (1) y derive término a término en la parte (2).
\end{ejerciciobox}

\subsubsection[Ejercicio 2: Generador de dilataciones $\hat{D}$]{Ejercicio Propuesto 2: Invariancia Rotacional del Generador de Dilataciones $\hat{D}$}
\begin{ejerciciobox}[Ejercicio Propuesto 2: Invariancia Rotacional del Generador de Dilataciones $\hat{D}$]
    \textbf{Nivel: Avanzado}\\
    El generador hermítico de las transformaciones de escala (dilataciones espaciales) en tres dimensiones se define mediante la combinación simetrizada de posición y momento:
    \begin{equation}
        \hat{D} \equiv \frac{1}{2}\left(\hat{\vec{r}}\cdot\hat{\vec{p}} + \hat{\vec{p}}\cdot\hat{\vec{r}}\right) = \hat{\vec{r}}\cdot\hat{\vec{p}} - \frac{3}{2}\iu\hbar\hat{I}.
    \end{equation}
    \begin{enumerate}
        \item Verifique que $\hat{D}$ es estrictamente autoadjunto ($\hat{D}\adj = \hat{D}$) y demuestre que conmuta con cada componente del momento angular: $[\hat{L}_i, \hat{D}] = 0$ para todo $i \in \{1, 2, 3\}$.
        \item Demuestre que el operador Casimir $\hat{L}^2$ conmuta con $\hat{D}$:
        \begin{equation}
            [\hat{L}^2, \, \hat{D}] = 0.
        \end{equation}
        ¿Cuál es el significado físico de este resultado en términos de simetrías continuas y transformaciones de escala?
    \end{enumerate}
    \textit{Intuición Didáctica y Pista:} En la física clásica, la cantidad escalar que mide dilataciones de escala es $\vec{r}\cdot\vec{p}$. No obstante, en mecánica cuántica, como $\hat{x}_k$ y $\hat{p}_k$ no conmutan, el producto simple no es hermítico: $(\hat{\vec{r}}\cdot\hat{\vec{p}})\adj = \hat{\vec{r}}\cdot\hat{\vec{p}} - 3\iu\hbar\hat{I}$. La \textbf{regla de oro} de la cuantización canónica para construir observables físicos a partir de productos clásicos es promediar simétricamente ambos órdenes: $\frac{1}{2}(\hat{A}\hat{B} + \hat{B}\hat{A})$ (conocido como ordenamiento de Weyl). Como la constante $-\frac{3}{2}\iu\hbar\hat{I}$ conmuta con todo operador, concéntrese en evaluar $[\hat{L}_i, \hat{x}_k\hat{p}_k] = [\hat{L}_i, \hat{x}_k]\hat{p}_k + \hat{x}_k[\hat{L}_i, \hat{p}_k]$ y contraiga el tensor simétrico resultante con el Levi-Civita $\varepsilon_{ikm}$.
\end{ejerciciobox}

\subsubsection[Ejercicio 3: Robertson--Schrödinger y chirp]{Ejercicio Propuesto 3: Relación de Robertson--Schrödinger en Paquetes con Chirp}
\begin{ejerciciobox}[Ejercicio Propuesto 3: Relación de Robertson--Schrödinger en Estados Coherentes Comprimidos]
    \textbf{Nivel: Avanzado (Física Cuántica Contemporánea)}\\
    Considere una partícula unidimensional en un estado caracterizado por una función de onda con chirp o modulación de fase cuadrática:
    \begin{equation}
        \psi(x) = \left(\frac{1}{\pi\sigma^2}\right)^{1/4}\exp\left(-\frac{x^2}{2\sigma^2} + \iu\,\frac{\alpha}{2}x^2\right), \quad \alpha \in \mathbb{R}.
    \end{equation}
    \begin{enumerate}
        \item Demuestre que la incertidumbre espacial sigue siendo $\Delta x = \frac{\sigma}{\sqrt{2}}$.
        \item Demuestre que la incertidumbre en el momento se incrementa a:
        \begin{equation}
            \Delta p_x = \frac{\hbar}{\sqrt{2}\sigma}\sqrt{1 + \alpha^2\sigma^4}.
        \end{equation}
        \item Calcule la covarianza cuántica $\mathrm{Cov}(\hat{x}, \hat{p}_x) = \frac{1}{2}\ev{\{\hat{x}, \hat{p}_x\}}$ y demuestre que satura con igualdad exacta la relación generalizada de \textbf{Robertson--Schrödinger} \eqref{eq:robertson_schrodinger}, pero es estrictamente mayor que la cota simple de Heisenberg $\hbar/2$.
    \end{enumerate}
\end{ejerciciobox}

\begin{figure}[H]
    \centering
    % Fig. 2: gaussiana con chirp (sigma = 1, hbar = 1)
\begin{tikzpicture}
\begin{groupplot}[
  group style={group size=2 by 1, horizontal sep=2.2cm},
  width=0.5\linewidth, height=5.2cm,
  axis lines=left, every axis plot/.append style={thick}]
\nextgroupplot[domain=-3.5:3.5, samples=400, xlabel={$x/\sigma$}, ylabel={$\mathrm{Re}\,\psi(x)$},
  title={\small Parte real de $\psi$}, ymin=-0.85, ymax=0.85,
  legend style={font=\scriptsize, draw=none, fill=none, at={(0.98,0.98)}, anchor=north east}]
  \addplot[uni-blue]{exp(-x^2/2)*cos(deg(0*x^2/2))*0.7511};
  \addplot[alert-border]{exp(-x^2/2)*cos(deg(3*x^2/2))*0.7511};
  \addplot[gray, dashed, thin, forget plot]{exp(-x^2/2)*0.7511};
  \addplot[gray, dashed, thin, forget plot]{-exp(-x^2/2)*0.7511};
  \legend{$\alpha\sigma^2=0$, $\alpha\sigma^2=3$}
\nextgroupplot[xlabel={$x/\sigma$}, ylabel={$\sigma \cdot p/\hbar$},
  title={\small Elipse de incertidumbre en $(x,p)$},
  xmin=-1.3, xmax=1.3, ymin=-3.2, ymax=3.2, axis lines=left,
  legend style={font=\scriptsize, draw=none, fill=none, at={(0.02,0.98)}, anchor=north west}]
  \addplot[uni-blue, domain=0:360, samples=100]
    ({0.7071*cos(x)}, {0.7071*sin(x)});
  \addplot[alert-border, domain=0:360, samples=100]
    ({0.7071*cos(x)}, {1.4142*cos(x)+0.7071*sin(x)});
  \legend{$\alpha\sigma^2=0$, $\alpha\sigma^2=2$}
\end{groupplot}
\end{tikzpicture}

    \caption{Gaussiana con chirp (Ejercicio Propuesto 3). Izquierda: la envolvente es la misma, pero las oscilaciones se aceleran en los bordes. Derecha: la elipse de incertidumbre en el plano $(x,p)$ se inclina (covarianza $\hbar\alpha\sigma^2/2\neq0$) y $\Delta p_x$ crece, mientras $\Delta x$ no cambia.}
    \label{fig:chirp}
\end{figure}

\subsubsection[Ejercicio 4: Acoplamiento $LS$ y centro de gravedad]{Ejercicio Propuesto 4: Acoplamiento Espín-Órbita $\hat{H}_{SO}$ y Regla del Centro de Gravedad}
\begin{ejerciciobox}[Ejercicio Propuesto 4: Estructura Fina y Desdoblamiento Espín-Órbita]
    \textbf{Nivel: Avanzado}\\
    En la estructura fina de un átomo hidrogenoide, el acoplamiento magnético entre el espín del electrón ($s = 1/2$) y su movimiento orbital ($\ell = 1$, nivel $2p$) genera el término de interacción:
    \begin{equation}
        \hat{H}_{SO} = \xi(\hat{\vec{L}}\cdot\hat{\vec{S}}), \quad \xi > 0,
    \end{equation}
    donde la constante $\xi \equiv \ev{\xi(r)}_{2p}$ es el valor esperado radial de la función de acoplamiento espín-órbita $\xi(r)$ evaluada en el subespacio del nivel $2p$ (consistente con la Sección~\ref{subsec:acoplamiento_espin_orbita}).
    \begin{enumerate}
        \item A partir de la definición del momento angular total $\hat{\vec{J}} = \hat{\vec{L}} + \hat{\vec{S}}$, demuestre la identidad operacional:
        \begin{equation}
            \hat{\vec{L}}\cdot\hat{\vec{S}} = \frac{1}{2}\left( \hat{J}^2 - \hat{L}^2 - \hat{S}^2 \right).
        \end{equation}
        \item Aplicando las reglas de adición de momentos angulares ($|\ell - s| \le j \le \ell + s$), determine los dos valores permitidos de $j$, su notación espectroscópica estándar ($^2P_{1/2}$ y $^2P_{3/2}$) y sus respectivas degeneraciones $g_j = 2j + 1$.
        \item Calcule el corrimiento de energía $\Delta E(j)$ inducido por $\hat{H}_{SO}$ sobre cada uno de los dos multipletes y demuestre que la separación energética (desdoblamiento de estructura fina) es:
        \begin{equation}
            \Delta E_{\text{fina}} = E(^2P_{3/2}) - E(^2P_{1/2}) = \frac{3}{2}\xi\hbar^2.
        \end{equation}
        \item Demuestre que se satisface la \textbf{regla de suma del centro de gravedad}:
        \begin{equation}
            \sum_j (2j + 1)\Delta E(j) = 0,
        \end{equation}
        y explique su significado físico en términos de la invariancia de la traza del operador perturbativo.
    \end{enumerate}
\end{ejerciciobox}

\begin{figure}[H]
\centering
\begin{tikzpicture}[>=Stealth,font=\small]
  \def\sc{2.0}   % cm por unidad de  xi*hbar^2

  % eje de energía
  \draw[->,grisUNI] (-0.8,-2.9)--(-0.8,2.0) node[above]{$E$};

  % nivel sin acoplamiento
  \draw[very thick] (0,0)--(2.2,0);
  \node[anchor=east] at (-0.05,0) {$2p$};
  \node[above,font=\footnotesize] at (1.1,0.05) {$g=6$};

  % niveles desdoblados
  \draw[very thick,rojoUNI] (5,{0.5*\sc})--(7.2,{0.5*\sc});
  \draw[very thick,azulUNI] (5,{-1.0*\sc})--(7.2,{-1.0*\sc});
  \draw[dotted,grisUNI] (2.2,0)--(5,{0.5*\sc});
  \draw[dotted,grisUNI] (2.2,0)--(5,{-1.0*\sc});

  \node[rojoUNI,anchor=south west,font=\small] at (5,{0.5*\sc+0.05})
       {${}^{2}\!P_{3/2}$\; ($g=4$)};
  \node[rojoUNI,anchor=north west,font=\footnotesize] at (5,{0.5*\sc-0.05})
       {$\Delta E=+\tfrac12\,\xi\hbar^2$};
  \node[azulUNI,anchor=south west,font=\small] at (5,{-1.0*\sc+0.05})
       {${}^{2}\!P_{1/2}$\; ($g=2$)};
  \node[azulUNI,anchor=north west,font=\footnotesize] at (5,{-1.0*\sc-0.05})
       {$\Delta E=-\,\xi\hbar^2$};

  % centro de gravedad
  \draw[dashed,naranjaUNI] (2.2,0)--(7.5,0);
  \node[naranjaUNI,anchor=south west,font=\footnotesize,align=left] at (7.7,0.1)
       {centro de gravedad:\\ $4\!\cdot\!\tfrac12+2\!\cdot\!(-1)=0$};

  % desdoblamiento
  \draw[<->,thick] (7.5,{-1.0*\sc})--(7.5,{0.5*\sc});
  \node[anchor=west,font=\footnotesize,align=left] at (7.7,{-0.75*\sc})
       {$\Delta E_{\rm fina}=\tfrac32\,\xi\hbar^2$};
\end{tikzpicture}
\caption{Estructura fina del nivel $2p$ ($\ell=1$, $s=\tfrac12$) por
$\hat H_{SO}=\xi\,\hat{\vec L}\cdot\hat{\vec S}$ con $\xi>0$. La perturbación
reparte el nivel en $^2P_{3/2}$ (subida) y $^2P_{1/2}$ (bajada) de modo que
$\sum_j(2j+1)\Delta E_j=0$: la traza de $\hat H_{SO}$ en el subespacio de
dimensión 6 es nula y la energía media no cambia.}
\label{fig:estructura-fina}
\end{figure}


\subsubsection[Ejercicio 5: Estados de Bell y traza parcial]{Ejercicio Propuesto 5: Producto Tensorial, Estados de Bell y Traza Parcial}
\begin{ejerciciobox}[Ejercicio Propuesto 5: Entrelazamiento Cuántico y Traza Parcial]
    \textbf{Nivel: Avanzado}\\
    Sean $\mathcal{H}_A$ y $\mathcal{H}_B$ dos espacios de Hilbert bidimensionales (qubits) con bases ortonormales $\{\ket{0}, \ket{1}\}$.
    \begin{enumerate}
        \item Demuestre formalmente por contradicción que el \textbf{estado de Bell}:
        \begin{equation}
            \ket{\Phi^+} \equiv \frac{1}{\sqrt{2}}\left( \ket{0}_A \otimes \ket{0}_B + \ket{1}_A \otimes \ket{1}_B \right)
        \end{equation}
        no puede factorizarse como un producto tensorial simple $\ket{\psi_A} \otimes \ket{\psi_B}$, demostrando que es un estado cuántico genuinamente entrelazado.
        \item Construya el operador de densidad global del estado puro $\hat{\rho} \equiv \ket{\Phi^+}\bra{\Phi^+}$ en forma matricial de $4\times 4$ en la base ordenada $\{\ket{00}, \ket{01}, \ket{10}, \ket{11}\}$.
        \item Calcule la matriz de densidad reducida del subsistema $A$ mediante la \textbf{traza parcial} sobre el subsistema $B$:
        \begin{equation}
            \hat{\rho}_A \equiv \mathrm{Tr}_B(\hat{\rho}) = \sum_{k=0}^1 (\hat{I}_A \otimes \bra{k}_B)\hat{\rho}(\hat{I}_A \otimes \ket{k}_B).
        \end{equation}
        Verifique que $\hat{\rho}_A = \frac{1}{2}\hat{I}_2$ y concluya por qué el subsistema individual $A$ se encuentra en un \textbf{estado maximalmente mezclado} (entropía cuántica máxima de von Neumann), a pesar de que el estado global del par es estrictamente puro.
    \end{enumerate}
\end{ejerciciobox}


%% ------------------------------------------------------------
%%  BIBLIOGRAFÍA
%% ------------------------------------------------------------
\clearpage
\section{Bibliografía y Referencias}
\printbibliography[heading=none]

%% ------------------------------------------------------------
%%  ANEXOS
%% ------------------------------------------------------------
\clearpage
%% ============================================================
%%  SECCIÓN 7 — Anexos y Demostraciones Auxiliares
%% ============================================================

\begin{appendices}

\section{Demostración Matricial de la Identidad de Levi-Civita}
\label{anexo:levi_civita}

En esta sección demostraremos de forma elegante la identidad tensorial fundamental:
\begin{equation}
    \varepsilon_{ijk}\varepsilon_{lmk} = \delta_{il}\delta_{jm} - \delta_{im}\delta_{jl},
\end{equation}
sin recurrir a verificar manualmente los $3^4 = 81$ casos combinatorios \cite{sakurai}.

Consideremos una base ortonormal estándar de $\mathbb{R}^3$, $\{\hat{e}_1, \hat{e}_2, \hat{e}_3\}$, donde $\hat{e}_a \cdot \hat{e}_b = \delta_{ab}$. El pseudotensor de Levi-Civita puede definirse geométricamente como el determinante formado por tres vectores base:
\begin{equation}
    \varepsilon_{ijk} = \det \begin{pmatrix} \hat{e}_i & \hat{e}_j & \hat{e}_k \end{pmatrix}.
\end{equation}
Consideremos ahora el producto de dos símbolos de Levi-Civita con índices totalmente libres $(i, j, k)$ y $(l, m, n)$. Como el determinante del producto de dos matrices es el producto de sus determinantes ($\det(A B^T) = \det(A)\det(B)$):
\begin{equation}
    \varepsilon_{ijk}\varepsilon_{lmn} = \det \left[ \begin{pmatrix} \hat{e}_i^T \\ \hat{e}_j^T \\ \hat{e}_k^T \end{pmatrix} \begin{pmatrix} \hat{e}_l & \hat{e}_m & \hat{e}_n \end{pmatrix} \right] = \det \begin{pmatrix}
        \hat{e}_i \cdot \hat{e}_l & \hat{e}_i \cdot \hat{e}_m & \hat{e}_i \cdot \hat{e}_n \\
        \hat{e}_j \cdot \hat{e}_l & \hat{e}_j \cdot \hat{e}_m & \hat{e}_j \cdot \hat{e}_n \\
        \hat{e}_k \cdot \hat{e}_l & \hat{e}_k \cdot \hat{e}_m & \hat{e}_k \cdot \hat{e}_n
    \end{pmatrix}.
\end{equation}
Dado que $\hat{e}_a \cdot \hat{e}_b = \delta_{ab}$, esta matriz de $3 \times 3$ contiene exclusivamente deltas de Kronecker:
\begin{equation}
    \varepsilon_{ijk}\varepsilon_{lmn} = \begin{vmatrix}
        \delta_{il} & \delta_{im} & \delta_{in} \\
        \delta_{jl} & \delta_{jm} & \delta_{jn} \\
        \delta_{kl} & \delta_{km} & \delta_{kn}
    \end{vmatrix}.
    \label{eq:levi_civita_det3x3}
\end{equation}
Para obtener la contracción deseada, igualamos los índices $n = k$ y sumamos sobre $k \in \{1, 2, 3\}$:
\begin{equation}
    \sum_{k=1}^3 \varepsilon_{ijk}\varepsilon_{lmk} = \sum_{k=1}^3 \begin{vmatrix}
        \delta_{il} & \delta_{im} & \delta_{ik} \\
        \delta_{jl} & \delta_{jm} & \delta_{jk} \\
        \delta_{kl} & \delta_{km} & \delta_{kk}
    \end{vmatrix}.
\end{equation}
Desarrollando el determinante por la tercera columna mediante la regla de Laplace:
\begin{equation}
    \sum_{k=1}^3 \varepsilon_{ijk}\varepsilon_{lmk} = \sum_{k=1}^3 \left[ \delta_{ik} (\delta_{jl}\delta_{km} - \delta_{jm}\delta_{kl}) - \delta_{jk} (\delta_{il}\delta_{km} - \delta_{im}\delta_{kl}) + \delta_{kk} (\delta_{il}\delta_{jm} - \delta_{im}\delta_{jl}) \right].
\end{equation}
Evaluemos la suma sobre $k$ término a término aplicando las propiedades del delta de Kronecker:
\begin{enumerate}
    \item Para el primer término:
    \begin{equation}
        \sum_{k=1}^3 \delta_{ik} (\delta_{jl}\delta_{km} - \delta_{jm}\delta_{kl}) = \delta_{jl}\delta_{im} - \delta_{jm}\delta_{il}.
    \end{equation}
    \item Para el segundo término:
    \begin{equation}
        -\sum_{k=1}^3 \delta_{jk} (\delta_{il}\delta_{km} - \delta_{im}\delta_{kl}) = -(\delta_{il}\delta_{jm} - \delta_{im}\delta_{jl}) = -\delta_{il}\delta_{jm} + \delta_{im}\delta_{jl}.
    \end{equation}
    \item Para el tercer término, notemos que $\sum_{k=1}^3 \delta_{kk} = \delta_{11} + \delta_{22} + \delta_{33} = 3$:
    \begin{equation}
        \sum_{k=1}^3 \delta_{kk} (\delta_{il}\delta_{jm} - \delta_{im}\delta_{jl}) = 3 (\delta_{il}\delta_{jm} - \delta_{im}\delta_{jl}).
    \end{equation}
\end{enumerate}
Sumando algebraicamente las tres contribuciones:
\begin{align}
    \sum_{k=1}^3 \varepsilon_{ijk}\varepsilon_{lmk} &= (\delta_{jl}\delta_{im} - \delta_{jm}\delta_{il}) - (\delta_{il}\delta_{jm} - \delta_{im}\delta_{jl}) + 3(\delta_{il}\delta_{jm} - \delta_{im}\delta_{jl}) \nonumber\\
    &= -(\delta_{il}\delta_{jm} - \delta_{im}\delta_{jl}) - (\delta_{il}\delta_{jm} - \delta_{im}\delta_{jl}) + 3(\delta_{il}\delta_{jm} - \delta_{im}\delta_{jl}) \nonumber\\
    &= (-1 - 1 + 3)(\delta_{il}\delta_{jm} - \delta_{im}\delta_{jl}) \nonumber\\
    &= \boxed{\delta_{il}\delta_{jm} - \delta_{im}\delta_{jl}}.
\end{align}
La identidad queda formalmente probada de manera exacta.

\section{Formulario Maestro de Relaciones de Conmutación y Fórmulas Clave}
\label{anexo:formulario}

Para una rápida consulta previa a los exámenes y sesiones de resolución de problemas, la Tabla~\ref{tab:formulario_maestro} sintetiza las fórmulas fundamentales deducidas en este documento.

\begin{table}[H]
    \centering
    \caption{Relaciones operacionales y algebraicas fundamentales de la mecánica cuántica.}
    \label{tab:formulario_maestro}
    \vspace{0.2cm}
    \resizebox{\linewidth}{!}{%
    \begin{tabular}{@{}llp{8.5cm}@{}}
        \toprule
        \textbf{Categoría} & \textbf{Expresión Matemática} & \textbf{Significado Físico / Aplicación} \\
        \midrule
        \textbf{Postulado de Medida} & $\Delta\hat{A} \equiv \hat{A} - \ev{\hat{A}}\hat{I}$ & Operador desviación (fluctuación pura alrededor de la media). \\
        \addlinespace[2pt]
        & $(\Delta A)^2 = \ev{\hat{A}^2} - \ev{\hat{A}}^2$ & Varianza cuántica intrínseca del estado $\ket{\psi}$. \\
        \midrule
        \textbf{Incertidumbre} & $(\Delta A)^2 (\Delta B)^2 \ge \frac{1}{4}|\ev{\{\Delta\hat{A},\Delta\hat{B}\}}|^2 + \frac{1}{4}|\ev{[\hat{A},\hat{B}]}|^2$ & Relación generalizada de Robertson--Schrödinger (incluye covarianza). \\
        \addlinespace[2pt]
        & $\Delta A \, \Delta B \ge \frac{1}{2}|\ev{[\hat{A},\hat{B}]}|$ & Relación de Heisenberg (al despreciar covarianza cuántica). \\
        \midrule
        \textbf{Álgebra de Lie} & $\eu^{\hat{A}}\hat{B}\eu^{-\hat{A}} = \sum_{n=0}^\infty \frac{1}{n!}\mathrm{ad}_{\hat{A}}^n(\hat{B})$ & Lema de Hadamard (transformaciones unitarias de semejanza). \\
        \addlinespace[2pt]
        & $\eu^{\hat{A}}\eu^{\hat{B}} = \eu^{\hat{A}+\hat{B}}\eu^{\frac{1}{2}[\hat{A},\hat{B}]}$ & Fórmula BCH (válida si $[[\hat{A},\hat{B}],\hat{A}]=[[\hat{A},\hat{B}],\hat{B}]=0$). \\
        \midrule
        \textbf{Momento Angular} & $[\hat{x}_i, \hat{p}_j] = \iu\hbar\delta_{ij}\hat{I}$ & Relaciones canónicas cartesianas de Heisenberg. \\
        \addlinespace[2pt]
        & $\hat{L}_k = \varepsilon_{kij}\hat{x}_i\hat{p}_j$ & Definición tensorial del momento angular orbital. \\
        \addlinespace[2pt]
        & $[\hat{L}_i, \hat{L}_j] = \iu\hbar\,\varepsilon_{ijk}\hat{L}_k$ & Álgebra de Lie del grupo de rotaciones $\mathfrak{so}(3)$. \\
        \addlinespace[2pt]
        & $\hat{\vec{L}} \times \hat{\vec{L}} = \iu\hbar\,\hat{\vec{L}} \neq \vec{0}$ & Producto vectorial cuántico de autoacoplamiento. \\
        \addlinespace[2pt]
        & $[\hat{L}^2, \hat{L}_i] = 0, \quad \forall\,i$ & Invariancia del operador Casimir (magnitud total al cuadrado). \\
        \addlinespace[2pt]
        & $\hat{L}_\pm = \hat{L}_x \pm \iu\hat{L}_y$ & Operadores escalera de subida (+) y bajada (-). \\
        \addlinespace[2pt]
        & $[\hat{L}_z, \hat{L}_\pm] = \pm\hbar\hat{L}_\pm$ & Conmutadores fundamentales de escalera. \\
        \addlinespace[2pt]
        & $\hat{L}_\pm\ket{\ell, m} = \hbar\sqrt{\ell(\ell+1)-m(m\pm 1)}\ket{\ell, m\pm 1}$ & Acción espectral normalizada en la base $\{\ket{\ell, m}\}$. \\
        \midrule
        \textbf{Potencial Central} & $\nabla^2 = \frac{1}{r^2}\partial_r(r^2\partial_r) - \frac{\hat{L}^2}{\hbar^2 r^2}$ & Laplaciano esférico con operador Casimir angular. \\
        \addlinespace[2pt]
        & $\hat{H} = -\frac{\hbar^2}{2\mu}\frac{1}{r^2}\partial_r(r^2\partial_r) + \frac{\hat{L}^2}{2\mu r^2} + V(r)$ & Hamiltoniano central con barrera centrífuga cuántica. \\
        \midrule
        \textbf{Producto Tensorial} & $(\hat{A}_1 \otimes \hat{B}_1)(\hat{A}_2 \otimes \hat{B}_2) = (\hat{A}_1\hat{A}_2)\otimes(\hat{B}_1\hat{B}_2)$ & Composición de operadores en espacios producto. \\
        \addlinespace[2pt]
        & $\hat{\vec{J}} = \hat{\vec{J}}_1 + \hat{\vec{J}}_2, \quad [\hat{J}_i, \hat{J}_j] = \iu\hbar\,\varepsilon_{ijk}\hat{J}_k$ & Momento angular total acoplado. \\
        \addlinespace[2pt]
        & $|j_1 - j_2| \le j \le j_1 + j_2$ & Regla triangular de adición de momentos angulares. \\
        \midrule
        \textbf{CCOC} & $\{\hat{H}, \hat{L}^2, \hat{L}_z\}$ (orbital; con espín, $\{\hat{H}, \hat{L}^2, \hat{L}_z, \hat{S}^2, \hat{S}_z\}$) & Base ortonormal común $\ket{n, \ell, m_\ell}$ (orbital) o $\ket{n, \ell, m_\ell, s, m_s}$ en átomos hidrogenoides ($\hat{S}^2$ es constante trivial para $s=1/2$). \\
        \bottomrule
    \end{tabular}%
    }
\end{table}

\section{Respuestas y Guía Metodológica de los Ejercicios Propuestos}
\label{anexo:solucionario}

A continuación se presentan las resoluciones analíticas condensadas de los ejercicios propuestos de la Sección~\ref{sec:ejercicios} para corroborar los resultados obtenidos en el autoestudio.

\subsection[Solución 1: Conmutadores de exponenciales]{Solución al Ejercicio 1: Conmutadores de Exponenciales de Operadores}
\begin{enumerate}
    \item Para $n = 1$, $[\hat{A}, \hat{B}] = 1\hat{A}^0[\hat{A}, \hat{B}]$, lo cual es trivialmente cierto. Suponiendo válido para $n = k$, evaluamos para $k+1$:
    \begin{align}
        [\hat{A}^{k+1}, \hat{B}] &= [\hat{A}\hat{A}^k, \hat{B}] = \hat{A}[\hat{A}^k, \hat{B}] + [\hat{A}, \hat{B}]\hat{A}^k \nonumber\\
        &= \hat{A}(k\hat{A}^{k-1}[\hat{A}, \hat{B}]) + [\hat{A}, \hat{B}]\hat{A}^k = k\hat{A}^k[\hat{A}, \hat{B}] + \hat{A}^k[\hat{A}, \hat{B}] \nonumber\\
        &= (k+1)\hat{A}^k[\hat{A}, \hat{B}],
    \end{align}
    donde se utilizó que $[\hat{A}, \hat{B}]$ conmuta con $\hat{A}$. Por inducción, queda demostrado para todo $n \in \mathbb{N}$.
    \item Expandiendo la exponencial:
    \begin{align}
        [\eu^{\hat{A}}, \hat{B}] &= \sum_{n=0}^\infty \frac{1}{n!}[\hat{A}^n, \hat{B}] = \sum_{n=1}^\infty \frac{1}{n!} n\hat{A}^{n-1}[\hat{A}, \hat{B}] \nonumber\\
        &= [\hat{A}, \hat{B}]\sum_{n=1}^\infty \frac{\hat{A}^{n-1}}{(n-1)!} = [\hat{A}, \hat{B}]\sum_{m=0}^\infty \frac{\hat{A}^m}{m!} = [\hat{A}, \hat{B}]\eu^{\hat{A}}.
    \end{align}
\end{enumerate}

\subsection[Solución 2: Generador de dilataciones $\hat{D}$]{Solución al Ejercicio 2: Hermiticidad e Invariancia del Generador de Dilataciones $\hat{D}$}
\begin{enumerate}
    \item \textit{Hermiticidad y Conmutador con $\hat{L}_i$:} Tomando el adjunto hermítico de $\hat{D} = \frac{1}{2}(\hat{\vec{r}}\cdot\hat{\vec{p}} + \hat{\vec{p}}\cdot\hat{\vec{r}})$:
    \begin{equation}
        \hat{D}\adj = \frac{1}{2}\left( (\hat{\vec{r}}\cdot\hat{\vec{p}})\adj + (\hat{\vec{p}}\cdot\hat{\vec{r}})\adj \right) = \frac{1}{2}\left( \hat{\vec{p}}\cdot\hat{\vec{r}} + \hat{\vec{r}}\cdot\hat{\vec{p}} \right) = \hat{D}.
    \end{equation}
    Además, usando $[\hat{x}_k, \hat{p}_k] = \iu\hbar\hat{I}$ (sumando sobre $k=1,2,3$ da $3\iu\hbar\hat{I}$), se tiene $\hat{\vec{p}}\cdot\hat{\vec{r}} = \hat{\vec{r}}\cdot\hat{\vec{p}} - 3\iu\hbar\hat{I}$, lo que reescribe $\hat{D} = \hat{\vec{r}}\cdot\hat{\vec{p}} - \frac{3}{2}\iu\hbar\hat{I}$.\\
    Dado que el múltiplo escalar conmuta con todo operador, el conmutador con $\hat{L}_i$ se reduce a:
    \begin{equation}
        [\hat{L}_i, \hat{D}] = [\hat{L}_i, \hat{\vec{r}}\cdot\hat{\vec{p}}] = [\hat{L}_i, \hat{x}_k\hat{p}_k] = [\hat{L}_i, \hat{x}_k]\hat{p}_k + \hat{x}_k[\hat{L}_i, \hat{p}_k].
    \end{equation}
    Sustituyendo los resultados de operador vectorial $[\hat{L}_i, \hat{x}_k] = \iu\hbar\varepsilon_{ikm}\hat{x}_m$ y $[\hat{L}_i, \hat{p}_k] = \iu\hbar\varepsilon_{ikm}\hat{p}_m$:
    \begin{align}
        [\hat{L}_i, \hat{D}] &= (\iu\hbar\,\varepsilon_{ikm}\hat{x}_m)\hat{p}_k + \hat{x}_k(\iu\hbar\,\varepsilon_{ikm}\hat{p}_m) \nonumber\\
        &= \iu\hbar\,\varepsilon_{ikm}\left( \hat{x}_m\hat{p}_k + \hat{x}_k\hat{p}_m \right).
    \end{align}
    Como $\hat{x}_m$ y $\hat{p}_k$ conmutan cuando $m \neq k$, la cantidad entre paréntesis es estrictamente \textbf{simétrica} ante el intercambio de $m$ y $k$, mientras que $\varepsilon_{ikm}$ es estrictamente \textbf{antisimétrico}. Por el teorema de contracción simétrico-antisimétrico demostrado en la Sección~\ref{sec:momento_angular}, la contracción se anula idénticamente:
    \begin{equation}
        [\hat{L}_i, \hat{D}] = 0, \quad \forall\,i \in \{1, 2, 3\}.
    \end{equation}
    \item Como $\hat{L}^2 = \hat{L}_i\hat{L}_i$, aplicando Leibniz:
    \begin{equation}
        [\hat{L}^2, \hat{D}] = [\hat{L}_i\hat{L}_i, \hat{D}] = \hat{L}_i[\hat{L}_i, \hat{D}] + [\hat{L}_i, \hat{D}]\hat{L}_i = \hat{L}_i(0) + (0)\hat{L}_i = 0.
    \end{equation}
    Físicamente, $\hat{D}$ es el generador hermítico de dilataciones espaciales (transformaciones de escala $\vec{r} \to \eu^\lambda\vec{r}$); como las dilataciones son isotrópicas (escalares puros que no alteran los ángulos ni la dirección espacial), conmutan con todas las rotaciones del grupo $\mathrm{SO}(3)$.
\end{enumerate}

\subsection[Solución 3: Paquetes con chirp]{Solución al Ejercicio 3: Paquetes con Chirp y Saturación de Incertidumbre}
Para el estado $\psi(x) = C \exp\left(-\frac{x^2}{2\sigma^2} + \iu\frac{\alpha}{2}x^2\right)$:
\begin{enumerate}
    \item La densidad de probabilidad de posición es $|\psi(x)|^2 = |C|^2 \exp(-x^2/\sigma^2)$, idéntica al paquete sin fase. Por ende, $\ev{\hat{x}} = 0$ y $\ev{\hat{x}^2} = \sigma^2/2$, lo que arroja $\Delta x = \sigma/\sqrt{2}$.
    \item Al calcular $\hat{p}_x\psi(x) = -\iu\hbar\frac{\dif\psi}{\dif x} = -\iu\hbar\left(-\frac{x}{\sigma^2} + \iu\alpha x\right)\psi(x) = \left(\hbar\alpha x + \iu\frac{\hbar}{\sigma^2}x\right)\psi(x)$.\\
    Calculando la norma de este estado:
    \begin{equation}
        \ev{\hat{p}_x^2} = \left(\hbar^2\alpha^2 + \frac{\hbar^2}{\sigma^4}\right)\ev{\hat{x}^2} = \left(\hbar^2\alpha^2 + \frac{\hbar^2}{\sigma^4}\right)\frac{\sigma^2}{2} = \frac{\hbar^2}{2\sigma^2}(1 + \alpha^2\sigma^4).
    \end{equation}
    Por ende, $\Delta p_x = \frac{\hbar}{\sqrt{2}\sigma}\sqrt{1 + \alpha^2\sigma^4}$.
    \item La covarianza cuántica resulta ser:
    \begin{equation}
        \mathrm{Cov}(\hat{x}, \hat{p}_x) = \frac{1}{2}\ev{\hat{x}\hat{p}_x + \hat{p}_x\hat{x}} = \hbar\alpha\ev{\hat{x}^2} = \frac{\hbar\alpha\sigma^2}{2}.
    \end{equation}
    Evaluemos el miembro derecho de Robertson--Schrödinger:
    \begin{equation}
        \frac{1}{4}|\ev{\{\hat{x}, \hat{p}_x\}}|^2 + \frac{1}{4}|\ev{[\hat{x}, \hat{p}_x]}|^2 = \left(\frac{\hbar\alpha\sigma^2}{2}\right)^2 + \left(\frac{\hbar}{2}\right)^2 = \frac{\hbar^2}{4}(1 + \alpha^2\sigma^4).
    \end{equation}
    Evaluemos el producto de varianzas:
    \begin{equation}
        (\Delta x)^2 (\Delta p_x)^2 = \left(\frac{\sigma^2}{2}\right)\left(\frac{\hbar^2}{2\sigma^2}(1 + \alpha^2\sigma^4)\right) = \frac{\hbar^2}{4}(1 + \alpha^2\sigma^4).
    \end{equation}
    ¡Ambos lados son idénticos! El estado satura con igualdad matemática exacta la relación de Robertson--Schrödinger, demostrando que cuando existe covarianza no nula ($\alpha \neq 0$), la cota real de dispersión es estrictamente mayor que la cota básica de Heisenberg ($\hbar/2$).
\end{enumerate}

\subsection[Solución 4: Acoplamiento $LS$ y centro de gravedad]{Solución al Ejercicio 4: Acoplamiento Espín-Órbita y Regla del Centro de Gravedad}
\begin{enumerate}
    \item Expandiendo el cuadrado del operador momento angular total $\hat{\vec{J}} = \hat{\vec{L}} + \hat{\vec{S}}$:
    \begin{equation}
        \hat{J}^2 = (\hat{\vec{L}} + \hat{\vec{S}})\cdot(\hat{\vec{L}} + \hat{\vec{S}}) = \hat{L}^2 + \hat{S}^2 + \hat{\vec{L}}\cdot\hat{\vec{S}} + \hat{\vec{S}}\cdot\hat{\vec{L}}.
    \end{equation}
    Como los operadores espaciales conmutan con los operadores de espín ($[\hat{L}_i, \hat{S}_j] = 0$), el producto escalar conmuta: $\hat{\vec{S}}\cdot\hat{\vec{L}} = \hat{\vec{L}}\cdot\hat{\vec{S}}$. Por ende:
    \begin{equation}
        \hat{J}^2 = \hat{L}^2 + \hat{S}^2 + 2\hat{\vec{L}}\cdot\hat{\vec{S}} \implies \hat{\vec{L}}\cdot\hat{\vec{S}} = \frac{1}{2}\left( \hat{J}^2 - \hat{L}^2 - \hat{S}^2 \right).
    \end{equation}
    \item Para $\ell = 1$ y $s = 1/2$, la regla triangular establece $|1 - 1/2| \le j \le 1 + 1/2$, de donde:
    \begin{equation}
        j \in \{1/2, \, 3/2\}.
    \end{equation}
    Las etiquetas espectroscópicas $^{2s+1}L_j$ son:
    \begin{itemize}
        \item Para $j = 3/2$: $^{2}P_{3/2}$, con degeneración $g_{3/2} = 2(3/2) + 1 = 4$.
        \item Para $j = 1/2$: $^{2}P_{1/2}$, con degeneración $g_{1/2} = 2(1/2) + 1 = 2$.
    \end{itemize}
    Nótese que la dimensión total se conserva: $4 + 2 = 6 = (2\cdot 1 + 1)(2\cdot 1/2 + 1)$.
    \item Evaluando el valor propio de $\hat{H}_{SO} = \frac{\xi}{2}(\hat{J}^2 - \hat{L}^2 - \hat{S}^2)$ (donde $\xi \equiv \ev{\xi(r)}_{2p}$ es el valor esperado radial de la función de acoplamiento espín-órbita de la Sección~\ref{subsec:acoplamiento_espin_orbita}) sobre el estado acoplado $\ket{j, M}$ con $\ell=1$ y $s=1/2$:
    \begin{align}
        \Delta E(j) &= \frac{\xi\hbar^2}{2}\left[ j(j+1) - \ell(\ell+1) - s(s+1) \right] \nonumber\\
        &= \frac{\xi\hbar^2}{2}\left[ j(j+1) - 2 - \frac{3}{4} \right] = \frac{\xi\hbar^2}{2}\left[ j(j+1) - \frac{11}{4} \right].
    \end{align}
    \begin{itemize}
        \item Para $j = 3/2$: $j(j+1) = \frac{3}{2}\cdot\frac{5}{2} = \frac{15}{4}$, luego:
        \begin{equation}
            \Delta E(^2P_{3/2}) = \frac{\xi\hbar^2}{2}\left( \frac{15}{4} - \frac{11}{4} \right) = \frac{\xi\hbar^2}{2}(1) = +\frac{1}{2}\xi\hbar^2.
        \end{equation}
        \item Para $j = 1/2$: $j(j+1) = \frac{1}{2}\cdot\frac{3}{2} = \frac{3}{4}$, luego:
        \begin{equation}
            \Delta E(^2P_{1/2}) = \frac{\xi\hbar^2}{2}\left( \frac{3}{4} - \frac{11}{4} \right) = \frac{\xi\hbar^2}{2}(-2) = -\xi\hbar^2.
        \end{equation}
    \end{itemize}
    El desdoblamiento energético total resulta ser:
    \begin{equation}
        \Delta E_{\text{fina}} = \Delta E(^2P_{3/2}) - \Delta E(^2P_{1/2}) = \frac{1}{2}\xi\hbar^2 - (-\xi\hbar^2) = \frac{3}{2}\xi\hbar^2.
    \end{equation}
    \item Evaluando la suma ponderada por las degeneraciones:
    \begin{equation}
        \sum_j (2j+1)\Delta E(j) = 4\left(+\frac{1}{2}\xi\hbar^2\right) + 2(-\xi\hbar^2) = +2\xi\hbar^2 - 2\xi\hbar^2 = 0.
    \end{equation}
    Físicamente, la perturbación interna de espín-órbita redistribuye los niveles de energía pero no altera la energía media (la traza del operador en el subespacio de dimensión 6 es idénticamente nula).

\end{enumerate}

\subsection[Solución 5: Estados de Bell y traza parcial]{Solución al Ejercicio 5: Producto Tensorial, Estados de Bell y Traza Parcial}
\begin{enumerate}
    \item \textit{Demostración de entrelazamiento por contradicción:}
    Supongamos que $\ket{\Phi^+}$ es separable, es decir, existen $\ket{\psi_A} = a_0\ket{0} + a_1\ket{1}$ y $\ket{\psi_B} = b_0\ket{0} + b_1\ket{1}$ tales que:
    \begin{equation}
        \ket{\psi_A} \otimes \ket{\psi_B} = a_0 b_0\ket{00} + a_0 b_1\ket{01} + a_1 b_0\ket{10} + a_1 b_1\ket{11} = \frac{1}{\sqrt{2}}\ket{00} + \frac{1}{\sqrt{2}}\ket{11}.
    \end{equation}
    Igualando coeficientes base a base:
    \begin{equation}
        a_0 b_0 = \frac{1}{\sqrt{2}}, \qquad a_0 b_1 = 0, \qquad a_1 b_0 = 0, \qquad a_1 b_1 = \frac{1}{\sqrt{2}}.
    \end{equation}
    De $a_0 b_0 = 1/\sqrt{2} \neq 0$, se deduce que $a_0 \neq 0$ y $b_0 \neq 0$. De $a_0 b_1 = 0$ con $a_0 \neq 0$, se tiene forzosamente $b_1 = 0$. Pero si $b_1 = 0$, entonces $a_1 b_1 = 0$, lo cual contradice directamente que $a_1 b_1 = 1/\sqrt{2} \neq 0$. La contradicción demuestra que $\ket{\Phi^+}$ no puede escribirse como un producto tensorial simple; por ende, es un \textbf{estado entrelazado}.
    \item \textit{Matriz de densidad global de $4\times 4$:}
    En la base ordenada $\{\ket{00}, \ket{01}, \ket{10}, \ket{11}\}$, el vector de estado es $\ket{\Phi^+} = \frac{1}{\sqrt{2}}\begin{pmatrix} 1 & 0 & 0 & 1 \end{pmatrix}^T$. Por consiguiente:
    \begin{equation}
        \hat{\rho} = \ket{\Phi^+}\bra{\Phi^+} = \frac{1}{2}\begin{pmatrix} 1 \\ 0 \\ 0 \\ 1 \end{pmatrix}\begin{pmatrix} 1 & 0 & 0 & 1 \end{pmatrix} = \frac{1}{2}\begin{pmatrix}
            1 & 0 & 0 & 1 \\
            0 & 0 & 0 & 0 \\
            0 & 0 & 0 & 0 \\
            1 & 0 & 0 & 1
        \end{pmatrix}.
    \end{equation}
    \item \textit{Traza parcial y matriz densidad reducida:}
    Evaluamos la traza parcial sobre el subsistema $B$:
    \begin{align}
        \hat{\rho}_A &= \sum_{k=0}^1 \bra{k}_B \hat{\rho} \ket{k}_B = \bra{0}_B \hat{\rho} \ket{0}_B + \bra{1}_B \hat{\rho} \ket{1}_B \nonumber\\
        &= \frac{1}{2}\ket{0}_A\bra{0}_A + \frac{1}{2}\ket{1}_A\bra{1}_A = \frac{1}{2}\begin{pmatrix} 1 & 0 \\ 0 & 1 \end{pmatrix} = \frac{1}{2}\hat{I}_2.
    \end{align}
    Como $\hat{\rho}_A = \frac{1}{2}\hat{I}_2$, se tiene $\mathrm{Tr}(\hat{\rho}_A^2) = \mathrm{Tr}\left(\frac{1}{4}\hat{I}_2\right) = \frac{2}{4} = \frac{1}{2} < 1$. El subsistema local $A$ se encuentra en un \textbf{estado maximalmente mezclado}: la probabilidad de medir $\ket{0}$ o $\ket{1}$ es exactamente $1/2$ (entropía de von Neumann máxima $S = -\mathrm{Tr}(\rho_A \log_2 \rho_A) = 1$ bit). Toda la certeza del estado global está contenida en el entrelazamiento cuántico no local.
\end{enumerate}

\end{appendices}


\end{document}
